\documentclass[12pt]{article}
\usepackage{lineno}
%\usepackage{ametsoc}
\usepackage{natbib}
\usepackage{a4wide}
\usepackage{longtable}
\usepackage{cancel}
\usepackage{docmute}   % \include of the stand-alone documents below skips their preambles
\usepackage{amsfonts}
\usepackage{amssymb}

\usepackage{color}
\usepackage[pdftex]{graphicx}
\usepackage[pdftex,colorlinks, citecolor=black,linkcolor=black,urlcolor=black]{hyperref}
\definecolor{darkgreen}{rgb}{0.1,0.4,0.1}
\definecolor{darkblue}{rgb}{0.1,0.1,0.3}



\newif\ifdetail
\global\detailfalse
%\global\detailtrue

\def\bem#1{\colorbox{green}{\em #1}}

\newcommand{\beq}  { \begin{eqnarray} }
\newcommand{\eeq}  { \end{eqnarray}}
\newcommand{\beeq}  { \begin{eqnarray*} }
\newcommand{\eeeq}  { \end{eqnarray*}}
\newcommand{\eq }  [1]{Eq.~(\ref{#1})}  % Eq. (#1)
\newcommand{\eqs }  [1]{Eqs.~(\ref{#1})}
\newcommand{\sect}  [1]{Sec.~(\ref{#1})}  % Eq. (#1)
\newcommand{\fig}  [1]{Fig.~\ref{#1}}   % Fig. #1
\renewcommand{\v}[1]{{\mbox{\boldmath$ #1 $}}}  % Vektor durch Fettdruck
\newcommand{\m}    [1]{{\bf  #1 }}              % Matrix in sans serif
\newcommand{\gz}[1] {{{d #1} \over {dz}}}
\newcommand{\pt}[1]{{\partial_t #1}}
\newcommand{\ptt}[1]{{\partial_{tt} #1}}
\newcommand{\pttt}[1]{{\partial_{ttt} #1}}
\newcommand{\pttx}[1]{{{\partial^3 #1} \over {\partial t^2 \partial x}}}
\newcommand{\ptty}[1]{{{\partial^3 #1} \over {\partial t^2 \partial y}}}
\newcommand{\ptx}[1]{{{\partial^2 #1} \over {\partial t \partial x}}}
\newcommand{\pty}[1]{{{\partial^2 #1} \over {\partial t \partial y}}}
\newcommand{\pz}[1]{{\partial_z #1}}
\newcommand{\pzz}[1]{{\partial_{zz} #1}}
\newcommand{\px}[1]{{\partial_x #1}}
\newcommand{\py}[1]{{\partial_y #1}}
\newcommand{\pxi}[1]{{\partial_i #1}}
\newcommand{\pxj}[1]{{\partial_j #1}}
\newcommand{\pxk}[1]{{\partial_k #1}}
\newcommand{\pb}[1]{{{\partial #1} \over {\partial b}}}
%\newcommand{\pxy}[1]{{{\partial^2 #1} \over {\partial x \partial y}}}
\newcommand{\pxy}[1]{\partial_{xy} #1}
\newcommand{\pyy}[1]{\partial_{yy} #1}
\newcommand{\pxx}[1]{\partial_{xx} #1}
\newcommand{\kx}{{\bf k} \times}
\newcommand{\dt}[1]{{{d #1} \over {d t}}}
\newcommand{\Dt}[1]{{{D #1} \over {D t}}}
\newcommand{\Div}[1]{ \nabla \cdot #1 }
\renewcommand{\div}[1]{ \nabla_{yz} \cdot #1 }
\newcommand{\Bar}[1]{ \overline{ #1} }
\newcommand{\rvec} [1]{
    \raisebox{-1.5ex}{$\stackrel{\textstyle #1}{\neg}$} }% rotierter Vektor
\newcommand{\nablq}   {\rvec\nabla}                    % rotated nabla operator
\newcommand{\Vec}[1]{\left(\begin{array}{c} #1 \end{array}\right) }
%\renewcommand{\mho}{\mathcal{W}}
\newcommand{\pn}[1]{{{\partial #1} \over {\partial n}}}
\newcommand{\pnn}[1]{{{\partial #1} \over {\partial m}}}
\newcommand{\ps}[1]{{{\partial #1} \over {\partial s}}}
\newcommand{\p}{{\partial}}
\newcommand{\vn}{{\v \nabla}}
\newcommand{\lk}{\left(}
\newcommand{\rk}{\right)}

\newcommand{\tl}{\tilde}

\newcommand{\I}{\mathcal{I}}
\newcommand{\E}{\mathcal{E}}
\newcommand{\A}{\mathcal{A}}

\newcommand{\ep}{\epsilon}
\newcommand{\erfc}{\mathrm{erfc}}

\begin{document}

\begin{center}
{\LARGE Documentation of pyOM3}
\end{center}

\tableofcontents
\clearpage

\phantomsection
\addcontentsline{toc}{section}{Overview}
\section*{Overview}

pyOM3 is a set of ocean and wave models in Python.
Each model is a package in the directory \verb|pyOM3|.
This document collects the documentation of all packages.
Each documentation is a section of its own and can also be compiled alone.

\begin{longtable}{@{}lp{8.6cm}r@{}}
\hline
Package & Model & Page \\
\hline
\endhead
\verb|pe_simple| & primitive equations on a uniform Cartesian grid & \pageref{doc:pe_simple} \\
\verb|pe_realistic| & pe\_simple with a variable grid, the sphere and isoneutral mixing & \pageref{doc:pe_realistic} \\
\verb|nh_simple| & non-hydrostatic Boussinesq equations on a uniform grid, periodic in $x$ and $y$, with LES closures & \pageref{doc:nh_simple} \\
\verb|nh_realistic| & pe\_realistic with a prognostic $w$ and a non-hydrostatic pressure, Cartesian grid & \pageref{doc:nh_realistic} \\
\verb|wim_simple| & spectral wave model on a Cartesian grid & \pageref{doc:wim_simple} \\
\verb|wim_realistic| & wim\_simple on the sphere & \pageref{doc:wim_realistic} \\
--- & the exact scattering integral $S_{nl}$ of wim\_simple & \pageref{doc:wim_snl_exact} \\
\hline
\end{longtable}

A realistic package is documented by its differences to the simple package.
The documentation of pe\_realistic thus needs the one of pe\_simple.
The same holds for nh\_realistic, which needs pe\_realistic and pe\_simple, and for wim\_realistic.

The routine tests and their results are described in \verb|pyOM3/test/README_test_suite.md|.
The coupled configurations are described in \verb|setups/README.md|.

\phantomsection
\addcontentsline{toc}{section}{Practical tips for programming}
\section*{Practical tips for programming}

This section lists the few things that are needed to write a configuration.
Everything else is in the documentation of the packages.

\subsection*{The model object}

A model is a Python class.
A configuration derives a class from the class \verb|model| of a package.
It overloads the methods it needs, for instance \verb|set_coriolis|, \verb|set_topography|, \verb|set_initial_conditions| and \verb|set_forcing|.
Additional parts such as the TKE closure of pe\_realistic or the LES closures of nh\_simple are further base classes in front, for instance \verb|class model(tke, pe_realistic.model.model)|.

The object is created with the dictionary of parameters and the number of processes in $x$ and $y$,
\begin{verbatim}
   m = model(parameter, n_pes_x = 2, n_pes_y = 2)
\end{verbatim}
Every entry of the dictionary becomes an attribute of the object.
All fields, grid arrays and parameters of the model are attributes of this one object.
In \verb|setups/coupled_common/run.py| the object is called \verb|M|, in the tests \verb|m|.

The time stepping is started with \verb|m.loop(n)| for $n$ time steps.
Before that, the pe family needs \verb|m.prepare_for_run()| after the topography is set.
nh\_simple and the wave models need \verb|m.setup()| after the object is created.
The snapshot classes take the place of \verb|model| when output is wanted.

\subsection*{Arrays and their dimensions}

All fields carry halo points on every side of every axis.
The size of an axis is therefore $N = n + 2 h$ with $n$ cells and $h$ = \verb|halo| halo points.
The interior of a three-dimensional field \verb|a| is \verb|a[h:-h, h:-h, h:-h]|.

The order of the axes differs between the packages.
\begin{longtable}{@{}lll@{}}
\hline
Packages & Order & Shape \\
\hline
\endhead
pe\_simple, pe\_realistic, nh\_realistic & $z, y, x$ & \verb|(Nz, Ny, Nx)| \\
nh\_simple & $x, y, z$ & \verb|(Nx, Ny, Nz)| \\
wim\_simple, wim\_realistic & direction, frequency, $y$, $x$ & \verb|(Nd, Nf, Ny, Nx)| \\
\hline
\end{longtable}
\noindent
Two-dimensional fields of the pe and wim families have the shape \verb|(Ny, Nx)|.
nh\_simple stores its arrays in Fortran order, since its pressure solver needs that.
All other packages use the order of numpy.

In $z$ the index \verb|h| is the lowest cell and the index \verb|-h-1| is the surface cell, in all packages with $z$.
The tendencies have a first axis of length 3 for the time levels of the Adams--Bashforth step.
The tendency of $u$ at the current level is \verb|du[tau]|.

The velocities sit on the faces of the cells, the tracers in the centres.
$u$ with index $i$ sits at the eastern face of cell $i$.
$v$ with index $j$ sits at the northern face of cell $j$.
$w$ with index $k$ sits at the top face of cell $k$.
The shift \verb|rollx(a,-1)| gives the value at $i+1$, and \verb|rollx(a,1)| the value at $i-1$.
The same holds for \verb|rolly| and \verb|rollz|.

\subsection*{The main variables}

\begin{longtable}{@{}lp{10.6cm}@{}}
\hline
Package & Variables \\
\hline
\endhead
pe\_simple & \verb|u|, \verb|v|, \verb|w|, \verb|temp|, \verb|salt|, \verb|rho|, the hydrostatic pressure \verb|p_h|, the surface pressure \verb|p_s|.
  The masks \verb|maskT|, \verb|maskU|, \verb|maskV|, \verb|maskW|, the bottom index \verb|kbot|, the Coriolis parameter \verb|coriolis_t|.
  The grid spacings are the numbers \verb|dx|, \verb|dy|, \verb|dz| \\
pe\_realistic & as pe\_simple.
  The grid spacings are one-dimensional arrays \verb|dxt|, \verb|dxu|, \verb|dyt|, \verb|dyu|, \verb|dzt|, \verb|dzw| and the coordinates \verb|xt|, \verb|xu|, \verb|yt|, \verb|yu|, \verb|zt|, \verb|zw|, all with halo.
  \verb|dxt_x| has the shape \verb|(1,1,Nx)|, \verb|dyt_y| the shape \verb|(1,Ny,1)| and \verb|dzt_z| the shape \verb|(Nz,1,1)|, for use with three-dimensional fields.
  \verb|dx|, \verb|dy| and \verb|dz| are None \\
nh\_realistic & as pe\_realistic, and the non-hydrostatic pressure \verb|p_nh|, the tendency \verb|dw| and the mask \verb|maskWc| \\
nh\_simple & \verb|u|, \verb|v|, \verb|w|, the buoyancy \verb|b|, the pressure \verb|p|.
  The grid spacings are the numbers \verb|dx|, \verb|dy|, \verb|dz| \\
wim\_simple, wim\_realistic & the wave action density \verb|A| \\
\hline
\end{longtable}

\subsection*{Writing to arrays}

\verb|OM.np| is numpy or jax.numpy, after \verb|import pyOM3 as OM|.
The choice is made by the environment variable \verb|PYOM3_BACKEND|.
Arrays of JAX cannot be changed in place.
\verb|a[...] = b| therefore fails with JAX.
\verb|a = OM.modify_array(a, index, b)| works with both and replaces \verb|a[index]| by \verb|b|.
\verb|OM.add_to_array| adds instead.
The result must always be assigned.

After a field is set, \verb|m.apply_bc(a)| fills its halo points.
It exchanges them with the neighbouring processes or wraps them in a periodic direction.
Without it the halo points keep wrong values.

\subsection*{Initial conditions with domain decomposition}

With MPI each process holds the block of its cells and the halo around it.
With JAX one array holds the whole domain, and every shard carries its own halo points inside that array.
A field that is set by hand must follow this layout.

The simplest way in the pe family is \verb|m.set_initial_conditions(u, v, temp, salt, p_s)|.
It takes the fields of the whole domain without halo points, of the shape \verb|(nz, ny, nx)|.
It cuts out the block of the process with MPI and inserts the halo points with JAX.
\verb|p_s| is None for the rigid lid.
Each process then holds the whole field once, which may need too much memory for a large model.

In pe\_realistic and nh\_realistic the coordinates \verb|xt|, \verb|yt| and \verb|zt| already have the layout of the process.
A field built from them is therefore correct with MPI and with JAX,
\begin{verbatim}
   T = 10 + 1e-3*m.zt[:,None,None] + 1e-6*m.xt[None,None,:] \
          + 1e-6*m.yt[None,:,None]
   m.temp = m.apply_bc(OM.np.asarray(T))*m.maskT
\end{verbatim}
The three terms together give the shape \verb|(Nz, Ny, Nx)|.
This is the recommended way.

In pe\_simple and wim\_simple there are no coordinate arrays.
With MPI the local index $i$ has the global index $i - h + \verb|xs_pe| - 1$ of the interior cells, counted from 0.
The same holds in $y$ with \verb|ys_pe|.
The number of interior cells of the process is $\verb|Nx| - 2h$.
In the pe family \verb|xe_pe| should not be used, since it is \verb|Nx| on a single process.
With JAX the global index of a point depends on its shard.
The function \verb|global_index| of \verb|pyOM3/test/pe/config.py| gives it for both cases.

nh\_simple divides the domain in $y$ and $z$ and never in $x$.
The local index $i$ has the global index $i - h + \verb|xs_pe| - 1$, and the same holds with \verb|ys_pe| and \verb|zs_pe|.
\verb|pyOM3/test/nh/config.py| builds its initial fields in this way.

After the fields are set, \verb|apply_bc| fills the halo points.
In the pe family the density follows from \verb|get_rho| of \verb|pyOM3/density.py|.
\verb|set_initial_conditions| does both.
nh\_simple makes the initial velocity free of divergence with \verb|m.pressure_solver()|.
nh\_realistic does this with \verb|m.nh_projection()|.

\clearpage\phantomsection\label{doc:pe_simple}\addcontentsline{toc}{section}{pe\_simple}
\documentclass[12pt]{article}
\usepackage{lineno}
%\usepackage{ametsoc}
\usepackage{natbib}
\usepackage{a4wide}
\usepackage{longtable}
\usepackage{amsfonts}
\usepackage{amssymb}

\usepackage{color}
\usepackage[pdftex]{graphicx}
\usepackage[pdftex,colorlinks, citecolor=black,linkcolor=black,urlcolor=black]{hyperref}
\definecolor{darkgreen}{rgb}{0.1,0.4,0.1}
\definecolor{darkblue}{rgb}{0.1,0.1,0.3}



\newif\ifdetail
\global\detailfalse
%\global\detailtrue

\def\bem#1{\colorbox{green}{\em #1}}

\newcommand{\beq}  { \begin{eqnarray} }
\newcommand{\eeq}  { \end{eqnarray}}
\newcommand{\beeq}  { \begin{eqnarray*} }
\newcommand{\eeeq}  { \end{eqnarray*}}
\newcommand{\eq }  [1]{Eq.~(\ref{#1})}  % Eq. (#1)
\newcommand{\eqs }  [1]{Eqs.~(\ref{#1})}
\newcommand{\sect}  [1]{Sec.~(\ref{#1})}  % Eq. (#1)
\newcommand{\fig}  [1]{Fig.~\ref{#1}}   % Fig. #1
\renewcommand{\v}[1]{{\mbox{\boldmath$ #1 $}}}  % Vektor durch Fettdruck
\newcommand{\m}    [1]{{\bf  #1 }}              % Matrix in sans serif
\newcommand{\gz}[1] {{{d #1} \over {dz}}}
\newcommand{\pt}[1]{{\partial_t #1}}
\newcommand{\ptt}[1]{{\partial_{tt} #1}}
\newcommand{\pttt}[1]{{\partial_{ttt} #1}}
\newcommand{\pttx}[1]{{{\partial^3 #1} \over {\partial t^2 \partial x}}}
\newcommand{\ptty}[1]{{{\partial^3 #1} \over {\partial t^2 \partial y}}}
\newcommand{\ptx}[1]{{{\partial^2 #1} \over {\partial t \partial x}}}
\newcommand{\pty}[1]{{{\partial^2 #1} \over {\partial t \partial y}}}
\newcommand{\pz}[1]{{\partial_z #1}}
\newcommand{\pzz}[1]{{\partial_{zz} #1}}
\newcommand{\px}[1]{{\partial_x #1}}
\newcommand{\py}[1]{{\partial_y #1}}
\newcommand{\pxi}[1]{{\partial_i #1}}
\newcommand{\pxj}[1]{{\partial_j #1}}
\newcommand{\pxk}[1]{{\partial_k #1}}
\newcommand{\pb}[1]{{{\partial #1} \over {\partial b}}}
%\newcommand{\pxy}[1]{{{\partial^2 #1} \over {\partial x \partial y}}}
\newcommand{\pxy}[1]{\partial_{xy} #1}
\newcommand{\pyy}[1]{\partial_{yy} #1}
\newcommand{\pxx}[1]{\partial_{xx} #1}
\newcommand{\kx}{{\bf k} \times}
\newcommand{\dt}[1]{{{d #1} \over {d t}}}
\newcommand{\Dt}[1]{{{D #1} \over {D t}}}
\newcommand{\Div}[1]{ \nabla \cdot #1 }
\renewcommand{\div}[1]{ \nabla_{yz} \cdot #1 }
\newcommand{\Bar}[1]{ \overline{ #1} }
\newcommand{\rvec} [1]{
    \raisebox{-1.5ex}{$\stackrel{\textstyle #1}{\neg}$} }% rotierter Vektor
\newcommand{\nablq}   {\rvec\nabla}                    % rotated nabla operator
\newcommand{\Vec}[1]{\left(\begin{array}{c} #1 \end{array}\right) }
%\renewcommand{\mho}{\mathcal{W}}
\newcommand{\pn}[1]{{{\partial #1} \over {\partial n}}}
\newcommand{\pnn}[1]{{{\partial #1} \over {\partial m}}}
\newcommand{\ps}[1]{{{\partial #1} \over {\partial s}}}
\newcommand{\p}{{\partial}}
\newcommand{\vn}{{\v \nabla}}
\newcommand{\lk}{\left(}
\newcommand{\rk}{\right)}

\newcommand{\tl}{\tilde}

\newcommand{\I}{\mathcal{I}}
\newcommand{\E}{\mathcal{E}}
\newcommand{\A}{\mathcal{A}}

\newcommand{\ep}{\epsilon}

\begin{document}

 \section*{The simple Primitive equation model}

Consider the Boussinesq equations in hydrostatic approximation
\beq \label{eq1}
  \pt{\v u } +   f \rvec{\v u}  +   \vn p =   - \v u \cdot \vn \v u - w \pz{} \v u +  \v F( \v u)~~,~~
    \pt{b}   =  - \v u \cdot \vn b  - w \pz{} b +  M b
     \eeq
  with the diagnostic relations     $\pz p =b$ and $\vn \cdot \v u + \pz w =0$.
  Vectors are two-dimensional and $\rvec{\v u} $ denotes  anticlockwise
 rotation by $90^o$. 
  Pressure is scaled with background density $\rho_0$ and thermodynamics have been
  simplified to a buoyancy equation.
  $\v F$  and $M$ are the friction and mixing operators.
   At the free surface $z=\eta$, we also have the
      condition 
      \beq \pt p_s  = -   g \vn \cdot  \int_{-H}^\eta \v u dz 
       \eeq
  with $p_s = p|_{z=\eta} =  g \eta$.
  Alternatively, we use a rigid lid with $w=0$ at $z=0$
  and $p_s = p|_{z=0}$.
We split full pressure $p$ into hydrostatic part and surface pressure.
\beq
p = p_h(z) + p_s ~,~
p_h =  \int_{-h}^z b  dz'
\eeq
where  $p_s$ is to be determined by the external mode solvers.



\subsection*{Simple implicit external solver for rigid lid}

Velocity and buoyancy and pressure are all co-located in time in this scheme.
The buoyancy equation is integrated as
\beq
  b^{n+1}- b^n = \Delta t \lk  \Bar{\delta b}^{n+1/2} + M (b^{n})    \rk 
\eeq
with the advective buoyancy tendency
\beq
\delta b^n =  - \vn \cdot \v u^{n} b^n - \p_z w^{n} b^n ~~,~~
\Bar{ \delta b}^{n+1/2} = A \delta b^n + B \delta b^{n-1} + C \delta b^{n-2}
 \eeq
 interpolated with the Adam-Bashforth (AB3) scheme  to time level $n+1/2$.
 The biharmonic horizontal mixing (attribute \verb|Khbi|) is treated with an Euler forward step.
 Vertical mixing follows the options described in the section on vertical mixing and friction below.
 
An intermediate velocity $\v u^*$ is calculated from
\beq
\v u^* -\v  u^n = \Delta t \lk    \Bar{\delta \v u}^{n+1/2}   + \v F (\v u^{n}) \rk 
\eeq  
with the momentum tendency $\delta \v u$ excluding the surface pressure gradient and friction
\beq
 \delta \v u^n =  - f \rvec{\v u}^n - \vn  p_h^n  - \v u^n \cdot \vn \v u^n - w^n \pz{} \v u^n 
\eeq
These tendencies are extrapolated with AB3  to time level $n+1/2$ 
\beq
\Bar{ \delta \v u}^{n+1/2} = A \delta \v u^n + B \delta \v u^{n-1} + C \delta \v u^{n-2}
\eeq
Friction is treated with an Euler forward step.
From the vertically integrated continuity equation follows the condition
\beq
\vn \cdot  \int_{-h}^0  \v u^*  =  \Delta t  \vn \cdot h \vn p_s^{n+1}
\eeq
which is solved with a conjugate-gradient iteration for $p_s^{n+1}$
(attributes \verb|epsilon| for the convergence threshold and
\verb|max_iterations|; a BiCGStab variant is implemented as well and is
selected by editing the solver call in \verb|pressure.py|).
The iteration count and the estimated error are available per step as
\verb|congr_itts| and \verb|estimated_error|. The intermediate velocity is  then
corrected as
\beq
 \v u^{n+1} = \v u^* - \Delta t \vn p_s^{n+1}
\eeq
It is possible to include an implicit free surface in the time stepping.
Vertical velocity is calculated from continuity
\beq
 w^{n+1} = -  \int_{-h}^z \vn \cdot \v u^{n+1}
\eeq
This scheme is also used as default in the MITgcm.

\subsection*{Simple split-explicit external mode solver}


   
Velocity is known in this scheme at $t=n+1/2$, while $b$ and $p$ at level $n$.
The momentum equation is discretised as
\beq
\v u^* -\v  u^{n-1/2} = \Delta t \lk - f  \Bar{\rvec{\v u}}^n+   \Bar{\delta \v u}^{n}  - \vn  p_h^n - \vn p_s^n+ \v F (\v u^{n-1/2}) \rk 
\eeq  
with the intermediate velocity $\v u^*$ and 
AB3 extrapolation of Coriolis force and advection terms tendencies
\beq
 \Bar{\rvec{\v u}}^n =  A \rvec{\v u}^{n-1/2} + B \rvec{\v u}^{n-3/2} + C \rvec{\v u}^{n-5/2}
~~,~~ \Bar{ \delta \v u}^{n} = A \delta \v u^{n-1/2} + B \delta \v u^{n-3/2} + C \delta \v u^{n-5/2}
\eeq
and
\beq
 \delta \v u^{n-1/2} =     - \v u^{n-1/2} \cdot \vn \v u^{n-1/2} - w^{n-1/2} \pz{} \v u^{n-1/2}
\eeq
For the hydrostatic and surface pressure $p_h+p_s$ no extrapolation is made, since they are at level $n$ anyways, 
and friction is treated with Euler forward.

Now the momentum tendencies are vertically integrated but without
surface pressure and Coriolis contributions.
\beq
 \v R^n = \int_{-h}^0  \lk  \Bar{\delta \v u}^{n}  - \vn  p_h^n+ \v F(\v u^{n-1/2}) \rk dz 
\eeq
The external mode is integrated with sub-cycling over $m=1, ...M$
($M$ set by the attribute \verb|sub_cycling|) from $n$ to $n+1$ with
a simple, dissipative  scheme proposed by Demange et al (2019) using the forcing $\v R^n$
\beq
\v U^{n+(m+1)/M} - \v U^{n+m/M} &=& \Delta t/M \lk - f  \rvec{\Bar{\v U}}^{n+(m+1/2)/M} -  h \vn p_s^{n+m/M} + 
 \v R^n  \rk
 \\
 p_s^{n+(m+1)/M} - p_s^{n+m/M} &=& - g \Delta t/M \lk  (1+\theta)  \vn \cdot \v U^{n+(m+1)/M} 
 - \theta   \vn \cdot \v U^{n+m/M}
 \rk
\eeq
with the transport $\v U = \int_{-h}^0 \v u\, dz$, and $\Bar{\v U}^{n+(m+1/2)/M}$ interpolated
with AB2.   Demange et al (2019)  found  $\theta=0.14$ to be the best choice in terms of phase error and damping.
In the model $\theta$ is the attribute \verb|fb_theta| with default $0.14$.
Larger values damp grid-scale barotropic waves more strongly.

The sub cycle starts from the transport $\v U^n$ and the surface pressure $p_s^n$ which the
sub cycle of the previous time step ended with.
This transport is model state (\verb|U_bt|, \verb|V_bt|) and part of the restart.
It is computed from the vertical integral of $\v u^n$ only at a cold start.
Two transports result from the sub cycle.
The transport at the end of the sub cycle, $\v U^{n+1} = \v U^{n+M/M}$, starts the next sub cycle.
The average over the sub cycle
\beq
\tl{\v U}^{n+1}=\frac{1}{M} \sum_{m=1}^{M} \v U^{n+m/M}  + \frac{\theta}{M} \lk  \v U^{n+1} - \v U^n  \rk
\eeq
is used to correct the barotropic part of the intermediate velocity
\beq
  \v u^{n+1/2} = \v u^* -1/h \lk  \int_{-h}^0 \v u^* dz - \tl{\v U}^{n+1} \rk
\eeq
The next sub cycle must not start from the average $\tl{\v U}^{n+1}$.
That would pair the averaged transport with the surface pressure at the end of the sub cycle.
This mismatch grows from step to step when the Courant number of the external
gravity wave is large (Banerjee et al.\ 2024, the split-explicit solver of FESOM2).
The buoyancy equation is integrated using $\v u^{n+1/2}$ and $w^{n+1/2} = - \int_{-h}^z \vn \cdot \v u^{n+1/2}$
as
\beq
  b^{n+1}- b^n = \Delta t \lk  \delta b^{n+1/2} + M (b^{n})    \rk 
\eeq
with the buoyancy tendency
\beq
\delta b^n =  - \vn \cdot \v u^{n+1/2} b^n - \p_z w^{n+1/2} b^n
 \eeq
 interpolated with AB3 to time level $n+1/2$.
 This scheme is one option in FESOM2.

In both schemes the advective tracer flux through the surface face is zero by default, as in pyOM2.
Under the linear free surface the vertical velocity $w$ at the surface face is not zero.
A tracer that is uniform in space then changes in the surface cell at the rate $ b\, w|_{z=0} / \Delta z = - b\, \vn \cdot \int_{-h}^0 \v u\, dz / \Delta z$.
With the attribute \verb|free_surface_tracer_flux| the flux through the surface face is
$w\, b$ instead, with $w$ at the surface face and $b$ of the surface cell.
A uniform tracer then stays uniform.
This option is off by default.
The external-mode treatment is selected by the attributes
\verb|implicit_free_surface| and \verb|explicit_free_surface|: both false
gives the rigid lid, the first the implicit free surface (both solved with
the iterative solver above), the second the split-explicit scheme.
  

\section*{Vertical mixing and friction options}

Vertical mixing of buoyancy and vertical friction are organised in a
hierarchy of options similar to pyOM2, selected by the attribute
\verb|vmix_scheme| $\in$ \{\verb|'explicit'|, \verb|'implicit'|\} and two
switches (module \verb|implicit_vmix.py|; the previous flag
\verb|enable_implicit_vmix| maps to the implicit scheme with convective
adjustment, so existing configurations are unchanged).

\emph{Explicit scheme.} Fixed coefficients $\kappa$ (buoyancy) and $A_v$
(momentum) only, integrated forward in time,
\beq
 b^{n+1} = b^n + \Delta t\, \pz{} \lk \kappa\, \pz b^n \rk, \qquad
 \v u^{n+1} = \v u^{n+1} + \Delta t\, \pz{} \lk A_v\, \pz{\v u} \rk
\eeq
subject to the diffusive stability limit, which is checked at start-up
(\verb|check_vmix_stability|, called from \verb|prepare_for_run|):
$ 4 \max(\kappa, A_v)\, \Delta t / \Delta z^2 \le 1/2$; the model refuses to
run otherwise, and also refuses convective adjustment or the TKE closure
under the explicit scheme.

\emph{Implicit scheme.} Both operators are integrated with backward Euler
(unconditionally stable, so arbitrarily large coefficients are allowed),
as a tridiagonal solve per column with the coefficient located on the
$W$ points. Three sub-options for the coefficients:
(i) constant $\kappa$ and $A_v$;
(ii) convective adjustment (\verb|enable_conv_adjust|): background
$\kappa$, but $\kappa = \kappa_{conv}$ (large) where $N^2 < 0$;
(iii) the TKE closure (\verb|enable_tke|), which overrides all constant
coefficients for both buoyancy and momentum as described below.
Vertical friction is applied after the surface pressure correction: since
the vertical operator with zero-flux boundary conditions leaves the
column-integrated momentum unchanged, the barotropic mode is not
perturbed, and $w$ is re-diagnosed afterwards.

\emph{Bottom friction.} Two options act on the velocity $\v u_b$ in the lowest wet cell
of each $U$ and $V$ column. Both are off by default.
Linear friction (\verb|enable_bottom_friction|) has the rate $r = r_{bot}$ in 1/s
(attribute \verb|r_bot|).
Quadratic friction has the rate
$r = C_d |\v u_b| / \Delta z$ with the drag coefficient $C_d$ (attribute \verb|r_quad_bot|).
It is switched on by
\verb|enable_quadratic_bottom_friction|.
Here $|\v u_b|$ is the speed in the bottom cell. The other velocity component is averaged
to the $U$ or $V$ point from the four surrounding points.
If both options are on, the two rates add.
The friction is an implicit step on the bottom cell,
\beeq
 \v u_b^{new} = \v u_b / \lk 1 + \Delta t\, r \rk
\eeeq
In the rigid-lid and implicit free-surface scheme it is applied after the vertical friction.
In the split-explicit scheme it is applied before the external mode is integrated.
There the same bottom stress also acts on the transport inside the sub cycle.
Every sub step divides $\v U$ by $1 + (\Delta t/M)\, r\, \Delta z / h$.

\subsection*{TKE closure}

The closure of Gaspar et al.\ (1990) is ported from pyOM2
(\verb|tke.py|, mixin class). All closure constants are configuration
attributes. They are set before \verb|init_tke|.
The defaults are the pyOM2 \verb|tke_module| defaults $c_k = 0.1$, $c_\epsilon = 0.7$, $\alpha = 1$,
$\ell_{min} = 10^{-12}\,$m and $\kappa_M^{max} = 100\,$m$^2$/s.
The floors $\kappa_M^{min}$ and $\kappa_H^{min}$ of viscosity and diffusivity are zero by default, as in pyOM2.
With the TKE closure on, the constant coefficients $A_v$ and $\kappa$ of the hierarchy above are not used.
A background viscosity or diffusivity is therefore a choice of the configuration.
It is made by setting the floors. Note that the
pyOM2 configurations of the acc family deviate from the module defaults and
set $\alpha = 30$ (a proxy for vertical TKE transport by pressure work,
which spreads the surface-injected TKE substantially deeper) and
$\ell_{min} = 10^{-8}\,$m; for consistency with those experiments both
should be set accordingly (\verb|alpha_tke|, \verb|mxl_min|). The turbulent
kinetic energy $e$ is a prognostic variable on the $W$ points,
\beq
 \pt e = \pz{} \lk \alpha\, \kappa_M \pz e \rk
 + K_{diss} - P_{diss} - c_\epsilon\, e^{3/2} / \ell
\eeq
with the mixing length $\ell = \sqrt{2 e}/N$, limited such that it cannot
grow faster than the distance to the boundaries (two sweeps, pyOM2's
\verb|tke_mxl_choice=2|), and the diffusivities
\beq
 \kappa_M = \min( \kappa_M^{max},\, c_k\, \ell \sqrt{e} ), \qquad
 \kappa_H = \kappa_M / \max(1, \min(10,\, 6.6\, Ri)), \qquad
 Ri = N^2 \kappa_M / K_{diss}
\eeq
There is no separate convective branch.
Where $N^2 < 0$ the Richardson number is negative.
The Prandtl number is then one and $\kappa_H = \kappa_M$, exactly as in pyOM2.
This replaces the convective-adjustment option.
The ratio $K_{diss}/\kappa_M$ in $Ri$ is the squared vertical shear of the friction operator.
With \verb|mix_eulerian| (see below) it is $\pz{\bar{\v u}} \cdot \pz{} \lk \bar{\v u} + \v u_S \rk$.
Without it, it is the squared shear of the model velocity.
$N^2$ lives on the $W$ points, i.e.\ on the top face of each cell.
At the surface face there is no cell above, so $N^2$ is copied from the face below, as in pyOM2.
For \verb|eq_of_state| $=5$ the density of the upper cell is evaluated at the pressure of the
lower cell and compared with the density of the lower cell (adiabatic displacement, as in pyOM2).

Two sets of limits for the mixing length are available. They are selected by \verb|mxl_choice|.
The default \verb|mxl_choice| $=1$ is the pyOM2 form named above.
Its sweeps run in the order of \verb|tke.f90|.
First a downward sweep from the surface limits $\ell_k \le \ell_{k+1} + \Delta z$.
It starts from the unlimited value at the surface face.
Then the value at the surface face is limited to $\ell \le \ell_{min} + z_0 + \Delta z$,
with the roughness length $z_0$ given below.
Finally an upward sweep from the bottom limits $\ell_k \le \ell_{k-1} + \Delta z$.
The same length enters the diffusivity and the dissipation.
With \verb|mxl_choice| $=2$ the limits of NEMO 4.2 are used (\verb|nn_mxl| $=3$ with \verb|ln_mxl0|).
All lengths live on the $W$ points, i.e.\ at the top face of each tracer cell.
The index $n$ counts these points upward. Here $n=1$ is the lowest point of the array
and $n_s$ the surface face.
The starting point is the buoyancy length
\beeq
 \ell_n = \max \lk \sqrt{2 e} / \sqrt{\max(N^2, 10^{-12}\,\mbox{s}^{-2})},\ \ell_{min0} \rk
\eeeq
with $\ell_{min0} = 0.04\,$m (attribute \verb|mxl_min0|).
This floor is applied at all $W$ points, also at dry points where $e=0$.
The surface value is then reduced to the surface length, $\ell_{n_s} = \min(\ell_{n_s}, \ell_0)$, with $\ell_0$ given below.
Since $\ell_0$ is set before the downward sweep, it bounds the whole near-surface profile and not only the surface point.
Two profiles are formed from $\ell_n$. Going downward from the surface the length is limited by the distance to the surface,
\beeq
 \ell^{up}_n = \min \lk \ell_n,\ \ell^{up}_{n+1} + \Delta z \rk, \quad n = n_s-1, \ldots, 1, \qquad \ell^{up}_{n_s} = \ell_{n_s}
\eeeq
Going upward from the bottom of the array the length is limited by the distance to the bottom,
\beeq
 \ell^{dwn}_n = \min \lk \ell_n,\ \ell^{dwn}_{n-1} + \Delta z \rk, \quad n = 2, \ldots, n_s, \qquad \ell^{dwn}_1 = \ell_1
\eeeq
The two sweeps run over the whole array, dry points included.
In a column whose lowest wet cell lies above the bottom of the array, the $W$ point below the lowest wet cell is dry.
There $\ell_n = \ell_{min0}$, so the lowest wet point gets $\ell^{dwn} \le \ell_{min0} + \Delta z$.
In a column that is wet down to the bottom of the array the lowest point gets $\ell^{dwn}_1 = \ell_1$ without this limit.
The diffusivity uses the geometric mean of the two profiles and the dissipation uses their minimum,
\beeq
 \ell_k = \sqrt{\ell^{up}\, \ell^{dwn}} , \qquad \ell_\epsilon = \min \lk \ell^{up},\ \ell^{dwn} \rk
\eeeq
$\ell_k$ replaces $\ell$ in the diffusivity $\kappa_M$, and $\ell_\epsilon$ replaces $\ell$ in the dissipation $c_\epsilon e^{3/2}/\ell$.
The constant $c_\epsilon$ multiplies the dissipation here.
Writing the dissipation as $e^{3/2}/\ell_\epsilon$ with $\ell_\epsilon = \min(\ell^{up}, \ell^{dwn})/c_\epsilon$ is the same.
The surface length $\ell_0$ has two forms.
Without a roughness length it follows Charnock's relation
\beeq
 \ell_0 = \kappa_{vk}\, \beta\, |\v \tau| / g
\eeeq
with the kinematic surface stress $\v \tau$, i.e.\ the stress divided by $\rho_0$,
the von K\'arm\'an constant $\kappa_{vk} = 0.4$ (attribute \verb|mxl_kappa|)
and the Charnock constant $\beta$ (attribute \verb|mxl_beta|).
The default $\beta = 2\cdot 10^5$ is the value of NEMO. It gives $\ell_0 = 0.8\,$m at a stress of $0.1\,$Pa.
The value $\beta = 2\cdot 10^4$ gives $\ell_0 = 0.08\,$m at the same stress.
The modulus $|\v \tau|$ is formed pointwise from \verb|taux| and \verb|tauy|,
or taken from \verb|tau_modulus| if the configuration provides it.
With a roughness length $z_0 > 0$ (attribute \verb|z0_surf|) the surface length is $\ell_0 = \kappa_{vk}\, z_0$ instead.
The coupled configurations set $z_0 = 1.6\, H_s$ with the significant wave height $H_s$ from the wave model.
In both forms $\ell_0 \ge \ell_{min0}$.
One halo exchange of $\ell_0$ is made per time step.
A configuration whose stress carries valid halo values can avoid it by setting \verb|tau_halo_ok|.

The discretisation follows pyOM2: backward Euler with the dissipation
linearised in the diagonal ($c_\epsilon \sqrt{e^n}/\ell$ acting on
$e^{n+1}$), a surface $W$ cell of half thickness $\Delta z/2$, and the
surface flux boundary condition added to that cell,
\beq
 F_e(0) = \left\{
 \begin{array}{ll}
   \Phi_{oc} & \mbox{wave dissipation flux from the wave model} \\
   (|\v \tau|)^{3/2} & \mbox{classical parameterisation}
 \end{array} \right.
\eeq
selected by \verb|tke_surf_from_waves| (default on, falling back to the
classical form if no wave flux is present; $\v \tau$ kinematic; the first
form closes the energy cycle between the wave and ocean components, the
second is the pyOM2 default and serves as the control experiment).
The classical form carries the factor one of pyOM2.
Craig and Banner (1994) give $\alpha_{cb} |\v \tau|^{3/2}$ with $\alpha_{cb} \approx 100$.
pe\_simple has no attribute for this factor.
A configuration that wants it overloads \verb|_surface_forcing|,
as \verb|setups/coupled_common/exp_ocean.py| does with \verb|EXP_ALPHA_CB|. No
diffusive TKE, tracer or momentum flux passes through the surface face
itself: the surface forcing enters only as the explicit flux above, which
also removes a spurious shear production against the zero-velocity halo
above the surface (a phantom no-slip lid).
For wave-driven injection a surface roughness length $z_0$ can be set
(attribute \verb|z0_surf|, default $0$, e.g.\
$1.6\,H_s$ from the wave model in the sense of Mellor \& Blumberg 2004).
With \verb|mxl_choice| $=1$ it only raises the limit of the value at the surface face.
The downward sweep runs before that limit is applied, so the mixing length below the surface face does not change.
With \verb|mxl_choice| $=2$ it sets the surface length $\ell_0$ described above.
If the
surface cell is drained negative by buoyancy fluxes, the missing energy is
added back and booked in \verb|tke_surf_corr| as in pyOM2.
There is no other floor on $e$, as in pyOM2.
A floor $e \ge e_{min}$ is applied only if \verb|tke_min| $> 0$ is set.
The initial value of $e$ in all wet cells is \verb|tke_init| $= 10^{-12}\,$m$^2$/s$^2$, as in the pyOM2 configurations.

Vertical friction is
computed implicitly with $\kappa_M$ averaged to the $U$/$V$ points as in
pyOM2 ($\overline{\kappa_M}^i$).
A momentum flux exists only across a face between two wet $U$ (or $V$) cells, as in pyOM2.
TKE is likewise exchanged only between two wet $W$ levels.
The shear production is diagnosed
from the same discrete operator with the updated velocities,
\beq
 K_{diss} = \Bar{\, \overline{\kappa_M}^i \lk \pz{u^{n+1}} \rk^2 \,}^i
          + \Bar{\, \overline{\kappa_M}^j \lk \pz{v^{n+1}} \rk^2 \,}^j
\eeq
The backward-Euler friction step is $\v u^{n+1} - \v u^n = \Delta t\, \pz{} \lk \overline{\kappa_M}\, \pz{\v u^{n+1}} \rk$.
Multiplying it by $\v u^{n+1}$ and integrating over the column gives the kinetic energy it removes,
\beeq
 \int_{-h}^0 \frac{|\v u^n|^2 - |\v u^{n+1}|^2}{2 \Delta t}\, dz
 = \int_{-h}^0 \overline{\kappa_M} \lk \pz{\v u^{n+1}} \rk^2 dz
 + \int_{-h}^0 \frac{|\v u^{n+1} - \v u^n|^2}{2 \Delta t}\, dz
\eeeq
The first term on the right is $K_{diss}$ and is the forcing of the TKE equation.
The second term is the numerical dissipation of the backward-Euler step.
It is of order $\Delta t$ and is not put into the TKE equation.
The kinetic energy removed from the mean flow therefore reappears in the TKE forcing
up to this term. $P_{diss} = \kappa_H N^2$ is the
exact transfer of the implicit buoyancy mixing, evaluated with the
\emph{updated} buoyancy (the tracer solve precedes the TKE step; using the
pre-mixing $N^2$ over-drains the TKE whenever $\kappa_H \Delta t/\Delta z^2$
is not small).
At the surface face $P_{diss}$ is the surface buoyancy flux instead, as in pyOM2,
\beeq
 P_{diss}(0) = - \frac{g}{\rho_0} \lk \frac{\p \rho}{\p T} F_T + \frac{\p \rho}{\p S} F_S \rk
\eeeq
Here $F_T$ and $F_S$ are the surface fluxes of temperature and salt into the ocean
(K\,m/s and psu\,m/s, positive into the ocean).
The configuration provides them as the attributes \verb|forc_temp_surface| and \verb|forc_salt_surface|.
Surface cooling is thus a source of TKE.
Without these attributes the term is zero.
For \verb|eq_of_state| $=0$ the temperature is the buoyancy $b$ and $\p \rho / \p b = -\rho_0/g$ is used.
pyOM2 uses zero there and has no surface buoyancy source of TKE for this equation of state.
The surface $W$ cell has the thickness $\Delta z/2$.
The TKE that enters it per unit area and time is therefore
$F_e(0) + (\Delta z/2)\, J_s$ plus the correction \verb|tke_surf_corr| described above.
Here $J_s = - P_{diss}(0)$ is the surface buoyancy loss.
In stand-alone column
tests the discrete TKE budget (surface flux $+$ shear production $=$
dissipation $+$ buoyancy destruction $+$ storage) closes to machine
precision.

When the model velocity is the Lagrangian mean $\v u_L$ of the coupled
wave--ocean formulation, the attribute \verb|mix_eulerian| (default off,
requires the Stokes velocity to be present) lets the friction operator act
on the Eulerian mean velocity $\bar{\v u} = \v u_L - \v u_S$: the
down-gradient closure parameterises Eulerian fluctuations, so mixing
$\v u_L$ would spuriously diffuse the Stokes shear. The Stokes velocity is
subtracted before the implicit solve and added back afterwards, and the
shear-production booking becomes flux times \emph{Lagrangian} shear, split
into the non-negative Eulerian part $K_{diss}^E = \kappa_M (\pz{\bar{\v u}})^2$
and the sign-indefinite Stokes (Langmuir) part
$K_{diss}^L = \kappa_M \pz{\bar{\v u}} \cdot \pz{\v u_S}$, kept separately
in the diagnostics with $K_{diss} = K_{diss}^E + K_{diss}^L$ feeding the TKE
equation. Without \verb|mix_eulerian| this reduces bit-identically to the
plain booking above.

The calling sequence per time step is: diffusivities from the current
$e$ and $N^2$ at the start of the step (so buoyancy mixing and friction
use the same coefficients the transfers are paired with), vertical
friction after the pressure correction, and the TKE step at the end of
the time step.
The density is evaluated once per time step, as in pyOM2.
This happens after the vertical mixing of the tracers and before the TKE step.
$N^2$ and $P_{diss}$ of the TKE step use this updated density.
The hydrostatic pressure of the next time step is computed from the same density.


\section*{Model parameters and options}

All parameters and options are attributes of the model instance, set either
through the parameter dictionary passed at construction or by assigning the
attribute before \verb|prepare_for_run|; options with a default may be left
unset. The tables list every attribute the model reads.

\subsection*{Required parameters (no default)}

\begin{longtable}{@{}llp{8.4cm}@{}}
\hline
Attribute & Unit & Meaning \\
\hline
\endhead
\verb|nx|, \verb|ny|, \verb|nz| & --- & global grid size (cells, without halo
  points) \\
\verb|dx|, \verb|dy|, \verb|dz| & m & grid spacing; $\Delta z$ is a uniform
  scalar \\
\verb|dt| & s & time step \\
\verb|eq_of_state| & --- & equation of state: $0$ (temperature carries the
  buoyancy, no salinity), $1$ (linear $\rho(S,T)$, salinity prognostic),
  $5$ (GSW/TEOS-10, salinity prognostic, the pressure is the depth of the
  cell centre below the surface in dbar), $100$ (ideal-gas atmosphere) \\
\verb|implicit_free_surface| & --- & implicit free surface if true \\
\verb|explicit_free_surface| & --- & split-explicit free surface if true;
  both false: rigid lid \\
\verb|periodic_in_x|, \verb|periodic_in_y| & --- & lateral boundaries
  periodic (true) or closed (false) \\
\verb|epsilon| & --- & convergence threshold of the iterative surface-pressure
  solver (rigid lid / implicit free surface) \\
\verb|max_iterations| & --- & iteration cap of that solver \\
\verb|sub_cycling| & --- & number of barotropic sub-cycles $M$ (split-explicit
  scheme only) \\
\hline
\end{longtable}

\noindent
The constructor additionally takes \verb|n_pes_x|, \verb|n_pes_y| (processor
grid, default $1 \times 1$), \verb|halo_points| (default 1) and
\verb|comm_in| (default: MPI world communicator). The constants
\verb|grav| $=9.81$ and \verb|rho_0| $=1024$ are set internally (note: they
overwrite any value passed in the parameter dictionary). Coriolis parameter,
topography (\verb|kbot|) and initial conditions are set through the template
methods \verb|set_coriolis|, \verb|set_topography| and
\verb|set_initial_conditions|, overloaded in the configuration.

\subsection*{Forcing hooks and coupling fields}

\begin{longtable}{@{}lp{10.6cm}@{}}
\hline
Attribute / method & Meaning \\
\hline
\endhead
\verb|set_forcing(tau)| & optional user method; adds forcing to the momentum
  tendencies \verb|du[tau]|, \verb|dv[tau]| (enters the AB interpolation).
  Surface fluxes of heat and salt are added to \verb|dtemp[tau]| and
  \verb|dsalt[tau]| in the same way. Both schemes keep these tracer
  increments \\
\verb|set_forcing_direct()| & optional user method; direct momentum increment
  bypassing the AB weights, applied before the pressure solve. Used for the
  Stokes tendency $\Delta \v u_S$ of the Lagrangian-mean ($\v u_L$)
  formulation, which must be a plain forward increment.
  It is called only in the rigid-lid and implicit free-surface scheme.
  The split-explicit scheme does not call it \\
\verb|forc_temp_surface| & surface flux of temperature into the ocean
  (K\,m/s). It gives the surface buoyancy source $P_{diss}(0)$ of the TKE
  equation. It does not force the tracer. That is done in \verb|set_forcing| \\
\verb|forc_salt_surface| & the same for salt (psu\,m/s) \\
\verb|tau_modulus| & optional modulus of the kinematic surface stress at the
  $T$ points, used by \verb|mxl_choice| $=2$ instead of the pointwise modulus
  of \verb|taux|, \verb|tauy| \\
\verb|taux|, \verb|tauy| & kinematic wind stress; also the fallback surface
  TKE forcing $|\v \tau|^{3/2}$ \\
\verb|phi_oc| & wave-breaking energy flux from the wave model; surface TKE
  flux when \verb|tke_surf_from_waves| is on \\
\verb|u_S|, \verb|v_S| & Stokes velocity from the wave model; required by
  \verb|mix_eulerian| \\
\hline
\end{longtable}

\subsection*{Options with defaults}

\begin{longtable}{@{}llp{5.9cm}@{}}
\hline
Attribute & Default & Meaning \\
\hline
\endhead
\verb|superbee_advection| & unset & if set true, tracer advection uses the
  Superbee flux limiter instead of the second-order fluxes (momentum
  advection is always second order). Note: setting it to false raises an
  error; leave it unset for the second-order scheme \\
\verb|Khbi| & unset & biharmonic horizontal tracer diffusion coefficient
  (m$^4$/s); active if present, forward Euler \\
\verb|Ahbi| & unset & biharmonic horizontal friction coefficient (m$^4$/s);
  active if present, forward Euler. In the rigid-lid and implicit
  free-surface scheme it is applied to $\v u$ before the AB update. In the
  split-explicit scheme it is a tendency in the forcing $\v R^n$ of the
  external mode \\
\verb|vmix_scheme| & \verb|'explicit'| & vertical mixing/friction hierarchy:
  \verb|'explicit'| (forward Euler, coefficient checked against the
  diffusive stability limit) or \verb|'implicit'| (backward Euler,
  unconditionally stable). If unset, the legacy flag
  \verb|enable_implicit_vmix| decides \\
\verb|kappa_T| & $0$ & background vertical diffusivity (m$^2$/s) \\
\verb|A_v| & $0$ & background vertical viscosity (m$^2$/s) \\
\verb|enable_conv_adjust| & on iff implicit & convective adjustment:
  $\kappa = \kappa_{conv}$ where $N^2 < 0$ (implicit scheme only) \\
\verb|kappa_conv| & $1$ & convective-adjustment diffusivity (m$^2$/s) \\
\verb|enable_tke| & off & TKE closure; overrides the constant coefficients
  for both buoyancy and momentum (implicit scheme only) \\
\verb|enable_bottom_friction| & off & linear bottom friction on the lowest
  wet cell, rate \verb|r_bot| \\
\verb|r_bot| & $0$ & rate of the linear bottom friction (1/s) \\
\verb|enable_quadratic_bottom_friction| & off & quadratic bottom friction on
  the lowest wet cell, rate $C_d |\v u_b|/\Delta z$ \\
\verb|r_quad_bot| & $0$ & drag coefficient $C_d$ of the quadratic bottom
  friction \\
\verb|fb_theta| & $0.14$ & dissipation parameter $\theta$ of the
  split-explicit sub cycle \\
\verb|free_surface_tracer_flux| & off & tracer flux through the surface face.
  Off means zero flux, as in pyOM2. On means $w\, b$ with the surface-cell
  value, so a uniform tracer stays uniform \\
\hline
\end{longtable}

\subsection*{TKE closure constants (read at \texttt{init\_tke})}

\begin{longtable}{@{}llp{8.2cm}@{}}
\hline
Attribute & Default & Meaning \\
\hline
\endhead
\verb|c_k| & $0.1$ & diffusivity constant, $\kappa_M = c_k \ell \sqrt{e}$ \\
\verb|c_eps| & $0.7$ & dissipation constant, $c_\epsilon e^{3/2}/\ell$ \\
\verb|alpha_tke| & $1$ & vertical TKE diffusion $\alpha \kappa_M$; the pyOM2
  acc configurations use $30$ \\
\verb|mxl_min| & $10^{-12}$ & mixing-length floor (m); acc configurations use
  $10^{-8}$ \\
\verb|kappaM_min| & $0$ & viscosity floor (m$^2$/s), as in pyOM2. A background
  viscosity is set here \\
\verb|kappaM_max| & $100$ & viscosity cap (m$^2$/s) \\
\verb|kappaH_min| & $0$ & diffusivity floor (m$^2$/s). A background
  diffusivity is set here \\
\verb|tke_min| & $0$ & TKE floor (m$^2$/s$^2$), applied only if positive \\
\verb|tke_init| & $10^{-12}$ & initial TKE in wet cells (m$^2$/s$^2$) \\
\verb|mxl_choice| & $1$ & limits of the mixing length. $1$ is pyOM2,
  $2$ is NEMO 4.2 with a separate length for the dissipation \\
\verb|mxl_kappa| & $0.4$ & von Karman constant of the surface length
  (\verb|mxl_choice| $=2$) \\
\verb|mxl_beta| & $2\cdot 10^5$ & Charnock constant $\beta$ of the surface
  length (\verb|mxl_choice| $=2$) \\
\verb|mxl_min0| & $0.04$ & minimum surface and bottom length (m)
  (\verb|mxl_choice| $=2$) \\
\verb|tau_halo_ok| & off & set when \verb|taux|, \verb|tauy| carry valid halo
  values. Saves the halo exchange of the surface length
  (\verb|mxl_choice| $=2$) \\
\verb|z0_surf| & $0$ & surface roughness length (m), e.g.\ $1.6\,H_s$.
  Limits the value at the surface face (\verb|mxl_choice| $=1$) or sets the
  surface length (\verb|mxl_choice| $=2$) \\
\verb|tke_surf_from_waves| & on & surface TKE flux from the wave dissipation
  $\Phi_{oc}$; falls back to $|\v \tau|^{3/2}$ if no wave flux is present \\
\verb|mix_eulerian| & off & vertical friction acts on the Eulerian mean
  velocity $\v u_L - \v u_S$ ($\v u_L$ formulation), with the
  $K_{diss}^E$/$K_{diss}^L$ booking split \\
\hline
\end{longtable}

\subsection*{Restart state}

The model has no restart routine of its own. A configuration that continues a
run must save the following fields and set them again before the next step.
A restart is then bit-identical to the uninterrupted run.
This is tested by \verb|pyOM3/test/pe/restart.py|, which holds the reference list.
\begin{itemize}
\item Always the fields \verb|u|, \verb|v|, \verb|w|, \verb|temp|, \verb|salt|, \verb|rho|, \verb|p_h| and \verb|p_s|.
\item Always the three time levels of the tendencies \verb|du|, \verb|dv|, \verb|dtemp| and \verb|dsalt|.
\item The time-level pointers \verb|tau|, \verb|taum1|, \verb|taum2| and the AB coefficients of the driving loop.
The method \verb|loop| starts every call with a forward Euler step and the pointers $2,1,0$.
A restart must therefore drive the time steps itself.
\item With the TKE closure \verb|tke| and \verb|K_diss_v|.
The diffusivities of a step are built from $K_{diss}$ of the previous step through the Richardson number.
Without \verb|K_diss_v| the first step after the restart differs.
\item With the split-explicit scheme \verb|du_cor|, \verb|dv_cor|, \verb|U_bt| and \verb|V_bt|.
\item With \verb|mxl_choice| $=2$ \verb|taux| and \verb|tauy| (or \verb|tau_modulus|).
The closure uses them at the start of the step, before \verb|set_forcing| sets the new stress.
\end{itemize}

\subsection*{Snapshot diagnostics}

Two snapshot classes are available; either replaces the model class as the
base of the configuration and overrides \verb|prepare_for_run| and
\verb|diagnose|. Both read the attribute \verb|snapint| (snapshot interval in
seconds, required) and take the \verb|prepare_for_run| arguments
\verb|snap_file_name| and \verb|show_each_time_step| (default on: print every
time step; off: print at output steps only). When the iterative
surface-pressure solver is active, both print its per-step iteration count
and accumulate the total count between outputs.

\emph{snap\_nc4} writes a single netCDF file (default name
\verb|'snap.cdf'|), created with the T- and U-grid coordinate axes and the
variables $u$, $v$, $w$, temp, $p_h$, $\rho$ and $p_s$; salinity is included
when the equation of state carries it (\verb|eq_of_state| $>0$ and
$\ne 100$). At every output step each PE in turn writes its halo-free
sub-block into the global arrays of the shared file (serialised with
barriers). For the rigid-lid/implicit path the current time level of the
surface pressure is written, for the split-explicit path the plain 2D field.
This class works with the MPI domain decomposition and with
single-process JAX, but not in a multihost JAX environment, since all
processes must be able to open the same file.

\emph{snap\_raw} writes one \verb|.npy| file per variable, per PE and per
output step into a directory tree
\verb|<name>/pe=<rank>/n=<step>/| (default name \verb|'raw_snap'|; an
existing directory of that name is removed at start-up), with the
coordinate axes and the processor-grid layout
(\verb|decomposition.npy|) stored once at the top level. Each process only
touches its own subdirectory, and JAX sharded arrays are written through
their local shard, so this class also works in a multihost JAX
environment. The variable set is $u$, $v$, $w$, temp, salt, $p_h$ and
$p_s$ (salt always, $\rho$ not included). The fragmented output is meant
to be reassembled afterwards using the stored decomposition, e.g.\ into
zarr.


\end{document}



\clearpage\phantomsection\label{doc:pe_realistic}\addcontentsline{toc}{section}{pe\_realistic}
\documentclass[12pt]{article}
\usepackage{lineno}
%\usepackage{ametsoc}
\usepackage{natbib}
\usepackage{a4wide}
\usepackage{longtable}
\usepackage{amsfonts}
\usepackage{amssymb}

\usepackage{color}
\usepackage[pdftex]{graphicx}
\usepackage[pdftex,colorlinks, citecolor=black,linkcolor=black,urlcolor=black]{hyperref}
\definecolor{darkgreen}{rgb}{0.1,0.4,0.1}
\definecolor{darkblue}{rgb}{0.1,0.1,0.3}



\newif\ifdetail
\global\detailfalse
%\global\detailtrue

\def\bem#1{\colorbox{green}{\em #1}}

\newcommand{\beq}  { \begin{eqnarray} }
\newcommand{\eeq}  { \end{eqnarray}}
\newcommand{\beeq}  { \begin{eqnarray*} }
\newcommand{\eeeq}  { \end{eqnarray*}}
\newcommand{\eq }  [1]{Eq.~(\ref{#1})}  % Eq. (#1)
\newcommand{\eqs }  [1]{Eqs.~(\ref{#1})}
\newcommand{\sect}  [1]{Sec.~(\ref{#1})}  % Eq. (#1)
\newcommand{\fig}  [1]{Fig.~\ref{#1}}   % Fig. #1
\renewcommand{\v}[1]{{\mbox{\boldmath$ #1 $}}}  % Vektor durch Fettdruck
\newcommand{\m}    [1]{{\bf  #1 }}              % Matrix in sans serif
\newcommand{\gz}[1] {{{d #1} \over {dz}}}
\newcommand{\pt}[1]{{\partial_t #1}}
\newcommand{\ptt}[1]{{\partial_{tt} #1}}
\newcommand{\pttt}[1]{{\partial_{ttt} #1}}
\newcommand{\pttx}[1]{{{\partial^3 #1} \over {\partial t^2 \partial x}}}
\newcommand{\ptty}[1]{{{\partial^3 #1} \over {\partial t^2 \partial y}}}
\newcommand{\ptx}[1]{{{\partial^2 #1} \over {\partial t \partial x}}}
\newcommand{\pty}[1]{{{\partial^2 #1} \over {\partial t \partial y}}}
\newcommand{\pz}[1]{{\partial_z #1}}
\newcommand{\pzz}[1]{{\partial_{zz} #1}}
\newcommand{\px}[1]{{\partial_x #1}}
\newcommand{\py}[1]{{\partial_y #1}}
\newcommand{\pxi}[1]{{\partial_i #1}}
\newcommand{\pxj}[1]{{\partial_j #1}}
\newcommand{\pxk}[1]{{\partial_k #1}}
\newcommand{\pb}[1]{{{\partial #1} \over {\partial b}}}
%\newcommand{\pxy}[1]{{{\partial^2 #1} \over {\partial x \partial y}}}
\newcommand{\pxy}[1]{\partial_{xy} #1}
\newcommand{\pyy}[1]{\partial_{yy} #1}
\newcommand{\pxx}[1]{\partial_{xx} #1}
\newcommand{\kx}{{\bf k} \times}
\newcommand{\dt}[1]{{{d #1} \over {d t}}}
\newcommand{\Dt}[1]{{{D #1} \over {D t}}}
\newcommand{\Div}[1]{ \nabla \cdot #1 }
\renewcommand{\div}[1]{ \nabla_{yz} \cdot #1 }
\newcommand{\Bar}[1]{ \overline{ #1} }
\newcommand{\rvec} [1]{
    \raisebox{-1.5ex}{$\stackrel{\textstyle #1}{\neg}$} }% rotierter Vektor
\newcommand{\nablq}   {\rvec\nabla}                    % rotated nabla operator
\newcommand{\Vec}[1]{\left(\begin{array}{c} #1 \end{array}\right) }
%\renewcommand{\mho}{\mathcal{W}}
\newcommand{\pn}[1]{{{\partial #1} \over {\partial n}}}
\newcommand{\pnn}[1]{{{\partial #1} \over {\partial m}}}
\newcommand{\ps}[1]{{{\partial #1} \over {\partial s}}}
\newcommand{\p}{{\partial}}
\newcommand{\vn}{{\v \nabla}}
\newcommand{\lk}{\left(}
\newcommand{\rk}{\right)}

\newcommand{\tl}{\tilde}

\newcommand{\I}{\mathcal{I}}
\newcommand{\E}{\mathcal{E}}
\newcommand{\A}{\mathcal{A}}

\newcommand{\ep}{\epsilon}

\begin{document}

 \section*{The realistic primitive equation model pe\_realistic}

pe\_realistic solves the same equations as pe\_simple with the same time stepping.
This text describes only the differences to pe\_simple.
Everything that is not mentioned here is as described in the documentation of pe\_simple.

The differences are the following.
The grid spacings may vary in all three directions.
The horizontal coordinates may be longitude and latitude on the sphere, as in pyOM2.
The equation of state may be the one of TEOS-10.
Isoneutral mixing is available.
Horizontal momentum friction uses the discrete operators of pyOM2.
Tracer advection can be taken out of the Adams--Bashforth step.

pyOM2 is the reference for all discrete operators that are new here.
They were ported from pyOM2 and compared with pyOM2 in twin experiments.
The tests and their results are listed in \verb|pyOM3/pe_realistic/README_test_pe_realistic.md|.

The domain decomposition, the halo exchange and the iterative solver for the surface pressure are not part of the package.
They are imported from pe\_simple.
The grid is built by the class \verb|grid| in \verb|grid.py|.

\subsection*{Grid}

The model domain consists of $n_x \times n_y \times n_z$ cells.
The horizontal coordinates $x$ and $y$ are either Cartesian coordinates in metres or longitude and latitude in degrees.
The attribute \verb|coord_degree| selects longitude and latitude.
The vertical coordinate $z$ is in metres.

The widths of the cells are given by the attributes \verb|dxt|, \verb|dyt| and \verb|dzt|.
Each of them is either a number or an array of length $n_x$, $n_y$ or $n_z$ without halo points.
The numbers \verb|dx|, \verb|dy| and \verb|dz| of pe\_simple are accepted instead.
Without either of them the model stops with an error.

A cell has its centre at a $T$ point and its faces at $U$ points.
In $x$ the cell $i$ has the centre $x^t_i$, the width $\Delta x^t_i$ and the eastern face $x^u_i$.
The distance between the centres $x^t_i$ and $x^t_{i+1}$ is $\Delta x^u_i$.
The same holds in $y$ with $y^t_j$, $\Delta y^t_j$, $y^u_j$, $\Delta y^u_j$.
In $z$ the face on top of cell $k$ is the $W$ point $z^w_k$, and the distance between the centres $z^t_k$ and $z^t_{k+1}$ is $\Delta z^w_k$.

The grid is built from the widths as in pyOM2 (\verb|u_centered_grid| of \verb|numerics.f90|).
The faces are summed from the widths.
The centres follow from the faces by
\beq \label{eq:ucentred}
 x^u_i = x^u_{i-1} + \Delta x^t_i ~,~~
 x^t_{i+1} = 2 x^u_i - x^t_i ~,~~
 \Delta x^u_i = x^t_{i+1} - x^t_i
\eeq
with the start $x^t_0 = x^u_0 - \Delta x^t_0/2$ and the end $\Delta x^u_n = 2 \Delta x^t_n - \Delta x^u_{n-1}$ at the last cell $n$.
This recursion keeps the discrete gradient and divergence adjoint on a stretched grid.

On a stretched grid $x^t_i$ is not the midpoint between the faces $x^u_{i-1}$ and $x^u_i$.
This follows from \eq{eq:ucentred}.
It matters wherever a depth is taken from $z^t$, for instance in the equation of state below.

In $x$ and $y$ the widths are first extended by the halo points.
In a periodic direction the halo widths are copies of their periodic images.
In a closed direction they repeat the width of the edge cell.
The recursion \eq{eq:ucentred} then runs over the extended array.
The first interior $U$ point is placed at \verb|x_origin| and \verb|y_origin| (default 0).
In a periodic direction the halo coordinates are copies of the coordinates of their periodic images.
They are used only for the metric factors below.

In $z$ the recursion \eq{eq:ucentred} runs over the $n_z$ interior cells only, as in pyOM2.
The coordinates are then extended linearly into the halo, and the halo widths repeat the edge widths.
The array \verb|dzt| is given from the bottom to the surface.
The cell $k$ = \verb|halo| is the lowest cell and the cell $k$ = \verb|Nz-halo-1| is the surface cell, as in pe\_simple.

The end of the recursion, $\Delta z^w_n = 2 \Delta z^t_n - \Delta z^w_{n-1}$, can be zero or negative.
This happens for instance when $\Delta z^t$ grows by a constant ratio from cell to cell.
The model then stops with an error.
$\Delta z^t$ must therefore vary smoothly.

The origin of $z$ is the bottom face of the lowest cell by default.
$z$ is then the height above the bottom, as in pe\_simple.
With the attribute \verb|set_zero_at_surface| the origin is the surface face instead, as in pyOM2.

With \verb|coord_degree| the coordinates $x$ and $y$ are longitude $\lambda$ and latitude $\phi$ in degrees.
The widths are given in degrees as well.
The coordinates are built from the widths in degrees.
After that the widths are converted to metres with the factor $a \pi / 180$.
The radius of the earth is $a = 6370\,$km (attribute \verb|radius|).
$\pi$ is the double precision value.
pyOM2 uses a single precision value, which gives relative differences of $3 \cdot 10^{-8}$ in all metric quantities.

The metric factors are $\cos \phi$ at the $T$ rows (\verb|cost|), $\cos \phi$ at the $U$ rows (\verb|cosu|) and $\tan \phi/a$ at the $T$ rows (\verb|tantr|).
In Cartesian coordinates they are 1, 1 and 0.
The horizontal areas of the cells are
\beq
 A^t = \cos \phi^t\, \Delta x^t \Delta y^t ~,~~
 A^u = \cos \phi^t\, \Delta x^u \Delta y^t ~,~~
 A^v = \cos \phi^u\, \Delta x^t \Delta y^u
\eeq
for the $T$, $U$ and $V$ cells, with $\Delta x$ and $\Delta y$ in metres.

The attributes \verb|dx|, \verb|dy| and \verb|dz| are set to None after the grid is built.
Code from pe\_simple that still uses them therefore stops with an error.

The grid arrays are one-dimensional.
With MPI every process holds its part including the halo points.
With JAX every shard holds its own halo points, as all fields do in pe\_simple.
With JAX $n_x$ and $n_y$ must be divisible by the number of shards in $x$ and $y$.

\subsection*{Discrete operators}

All operators of pe\_simple are written in flux form.
In pe\_realistic every flux is multiplied by the width of the face it crosses.
Every flux divergence is divided by the volume of the cell.
On a uniform Cartesian grid every operator reduces to the one of pe\_simple.
This is tested to round-off.

The vertical velocity at the top face of cell $k$ follows from continuity,
\beq
 w_k = - \sum_{k' \le k} \Delta z^t_{k'} \lk
 \frac{u_i - u_{i-1}}{\cos \phi^t_j\, \Delta x^t_i}
 + \frac{\cos \phi^u_j\, v_j - \cos \phi^u_{j-1}\, v_{j-1}}{\cos \phi^t_j\, \Delta y^t_j} \rk
\eeq
with $u$ and $v$ of cell $k'$ and the sum over the wet cells of the column.

The tracer fluxes are those of pe\_simple.
The meridional flux is multiplied by $\cos \phi^u$.
The flux divergence is divided by $\cos \phi^t \Delta x^t$, $\cos \phi^t \Delta y^t$ and $\Delta z^t$.
The Superbee scheme uses the Courant numbers $|u| \Delta t /(\cos \phi^t \Delta x^t)$, $|\cos \phi^u v| \Delta t / (\cos \phi^t \Delta y^t)$ and $|w| \Delta t / \Delta z^t$.
The WENO scheme uses the coefficients of pe\_simple, those of a uniform grid, as Oceananigans does.
A uniform tracer then stays uniform and the content is conserved.
Where the widths of the cells change, the face value is of lower than fifth order.

The momentum advection is that of pe\_simple with volume transports.
The transports through the faces are $u \Delta z^t \Delta y^t$, $v \Delta z^t \cos \phi^u \Delta x^t$ and $w A^t$.
The divergence is divided by the volume $A^u \Delta z^t$ of the $U$ cell or $A^v \Delta z^t$ of the $V$ cell.

The Coriolis force averages $f v$ to the $U$ point and $f u$ to the $V$ point with the widths of the cells as weights, as in pyOM2,
\beq
 (\pt u)_{cor} &=& \frac{1}{4 \Delta x^u_i} \lk
   f_i\, \Delta x^t_i \lk v_{i,j} + v_{i,j-1} \rk
 + f_{i+1}\, \Delta x^t_{i+1} \lk v_{i+1,j} + v_{i+1,j-1} \rk \rk
 \\
 (\pt v)_{cor} &=& - \frac{1}{4 \cos \phi^u_j\, \Delta y^u_j} \lk
   f_j \cos \phi^t_j \Delta y^t_j \lk u_{i,j} + u_{i-1,j} \rk \right. \nonumber \\
 && \left. + f_{j+1} \cos \phi^t_{j+1} \Delta y^t_{j+1} \lk u_{i,j+1} + u_{i-1,j+1} \rk \rk
\eeq
with the Coriolis parameter $f$ at the $T$ points.
On a uniform grid this is the average of pe\_simple.

On the sphere the momentum equations contain the metric terms $+u v \tan \phi /a$ and $-u^2 \tan \phi /a$.
They are discretised with the same weights as the Coriolis force,
\beq
 (\pt u)_{met} &=& \frac{1}{8 \Delta x^u_i} \lk
   \lk \frac{\tan \phi}{a}\, \Delta x^t (u_i + u_{i-1})(v_j + v_{j-1}) \rk_i
 + \lk \cdots \rk_{i+1} \rk
 \\
 (\pt v)_{met} &=& - \frac{1}{8 \cos \phi^u_j\, \Delta y^u_j} \lk
   \lk \frac{\tan \phi}{a} \cos \phi^t \Delta y^t (u_i + u_{i-1})^2 \rk_j
 + \lk \cdots \rk_{j+1} \rk
\eeq
where $(\cdots)$ repeats the bracket before it at the next $T$ point.
In Cartesian coordinates both terms vanish.

The pressure gradients are $-(p_{i+1}-p_i)/(\cos \phi^t \Delta x^u_i)$ and $-(p_{j+1}-p_j)/\Delta y^u_j$.
This holds for the hydrostatic pressure $p_h$ and for the surface pressure $p_s$.

The hydrostatic pressure is integrated downward from the surface cell with the distances $\Delta z^w$,
\beq
 p_h(n_z) = \frac{g}{2 \rho_0} \rho(n_z)\, \Delta z^w(n_z) ~,~~
 p_h(k) = p_h(k+1) + \frac{g}{2 \rho_0} \lk \rho(k+1) + \rho(k) \rk \Delta z^w(k)
\eeq
where $n_z$ denotes the surface cell here.

The surface pressure of the rigid lid and of the implicit free surface follows from the vertically integrated continuity equation, as in pe\_simple.
The depth of the water column at the $U$ and $V$ points is $h^u = \sum \Delta z^t$ and $h^v = \sum \Delta z^t$ over the wet $U$ and $V$ cells.
The discrete equation for $p_s$ is
\beeq
 && \frac{1}{\cos \phi^t \Delta x^t_i} \lk
   \frac{h^u_i (p_{i+1} - p_i)}{\cos \phi^t \Delta x^u_i}
 - \frac{h^u_{i-1} (p_i - p_{i-1})}{\cos \phi^t \Delta x^u_{i-1}} \rk \\
 &+& \frac{1}{\cos \phi^t_j \Delta y^t_j} \lk
   \frac{\cos \phi^u_j h^v_j (p_{j+1} - p_j)}{\Delta y^u_j}
 - \frac{\cos \phi^u_{j-1} h^v_{j-1} (p_j - p_{j-1})}{\Delta y^u_{j-1}} \rk
 = \frac{1}{\Delta t}\, \vn \cdot \int \v u^* dz
\eeeq
with the divergence on the right discretised as for $w$ above.
The implicit free surface adds $-p_s/(g \Delta t^2)$ on the left as in pe\_simple.
The equation is solved with the conjugate-gradient solver of pe\_simple.

The rigid lid does not converge on a domain that is periodic in $x$ and contains land.
pyOM2 uses a streamfunction with island integrals there, which pe\_realistic does not have.
Such domains need the implicit free surface.

The vertical mixing of tracers and momentum uses the widths $\Delta z^t$ and the distances $\Delta z^w$,
\beeq
 (\pt T)_{mix} = \frac{1}{\Delta z^t_k} \lk
   \kappa_k \frac{T_{k+1} - T_k}{\Delta z^w_k}
 - \kappa_{k-1} \frac{T_k - T_{k-1}}{\Delta z^w_{k-1}} \rk
\eeeq
with the coefficient $\kappa_k$ at the top face of cell $k$.
The stability limit of the explicit scheme becomes $4 \max(\kappa, A_v) \Delta t / \min(\Delta z^t \Delta z^w) \le 1/2$.

\subsection*{Horizontal friction}

Horizontal momentum friction uses the operators of pyOM2 (\verb|harmonic_friction.f90| and \verb|biharmonic_friction.f90|).
They place the fluxes and the metric factors on the rows of the $U$ and $V$ points.
pe\_simple uses an operator at the $T$ points for $u$ and $v$ instead.
On a grid with varying $\Delta y$ or on the sphere that operator is not exact for $u$ and $v$.

For $u$ the harmonic operator with coefficient $A$ is
\beeq
 F_u &=& \frac{f^x_i - f^x_{i-1}}{\cos \phi^t_j\, \Delta x^u_i}
       + \frac{f^y_j - f^y_{j-1}}{\cos \phi^t_j\, \Delta y^t_j} \\
 f^x_i &=& A \, (\cos \phi^t_j)^p \, \frac{u_{i+1} - u_i}{\cos \phi^t_j\, \Delta x^t_{i+1}} ~,~~
 f^y_j = A \, (\cos \phi^u_j)^p \cos \phi^u_j \, \frac{u_{j+1} - u_j}{\Delta y^u_j}
\eeeq
A flux exists only between two wet $U$ cells.

For $v$ the harmonic operator is
\beeq
 F_v &=& \frac{f^x_i - f^x_{i-1}}{\cos \phi^u_j\, \Delta x^t_i}
       + \frac{f^y_j - f^y_{j-1}}{\cos \phi^u_j\, \Delta y^u_j} \\
 f^x_i &=& A \, (\cos \phi^u_j)^p \, \frac{v_{i+1} - v_i}{\cos \phi^u_j\, \Delta x^u_i} ~,~~
 f^y_j = A \, (\cos \phi^t_{j+1})^p \cos \phi^t_{j+1} \, \frac{v_{j+1} - v_j}{\Delta y^t_{j+1}}
\eeeq
A flux exists only between two wet $V$ cells.

The harmonic friction is $\v u \leftarrow \v u + \Delta t\, (F_u, F_v)$ with $A$ = \verb|Ah|.
The exponent $p$ is \verb|hor_friction_cosPower| (default 0).
It is applied only in the rigid-lid and implicit free-surface scheme.

The biharmonic friction applies the harmonic operator twice, $\v u \leftarrow \v u - \Delta t\, F(F(\v u))$.
Each pass uses $A = \sqrt{A_{hbi}}$ with $A_{hbi}$ = \verb|Ahbi|.
The exponent $p$ is \verb|biha_friction_cosPower| (default 0).
Since $(\cos \phi)^p$ enters each pass, the viscosity scales with $(\cos \phi)^{2p}$.
In the split-explicit scheme the biharmonic friction enters the forcing of the external mode as the tendency $-F(F(\v u))$ of the same operator.

The biharmonic mixing of tracers (\verb|Khbi|) uses the operator at the $T$ points with the metric factors of the tracer fluxes above.

The bottom friction is that of pe\_simple.
The thickness of the bottom cell is $\Delta z^t$ of that cell instead of $\Delta z$.

\subsection*{Split-explicit external mode}

The split-explicit scheme is that of pe\_simple.
The Coriolis force on the transports uses the weights of the Coriolis force above.
The pressure gradient and the divergence of the transports carry the metric factors above.

On the sphere the external mode also contains the metric terms.
They are the vertical integrals of $+u v \tan \phi/a$ and $-u^2 \tan \phi /a$ with $u = U/h^u$ and $v = V/h^v$.
They are discretised as the metric terms above and multiplied by $h^u$ and $h^v$.
The metric terms of the baroclinic part are already part of the Coriolis tendency of the three-dimensional velocity.

\subsection*{Equation of state}

The attribute \verb|eq_of_state| selects the equation of state.
The value 0 means that temperature is buoyancy.
The value 1 is the linear equation of state.
The value 5 is the 48-term equation of state of TEOS-10 (GSW), as in pyOM2.
The value 100 is the ideal gas of the atmosphere.
The density is the anomaly to $\rho_0 = 1024\,$kg/m$^3$, as in pyOM2.

The pressure in the equation of state of TEOS-10 is the depth of the cell below the surface face in dbar, as in pyOM2.
It is $z^w$ of the surface face minus $z^t$ of the cell.
It is not $z^t + \Delta z^t/2$, since $z^t$ is not the midpoint of the cell on a stretched grid.
The value 100 needs the default origin of $z$ at the bottom.

$N^2$ for the value 5 uses the adiabatic displacement as in pe\_simple.
The density is evaluated once per time step as in pe\_simple.

\subsection*{TKE closure}

The TKE closure is that of pe\_simple with the following changes.
The mixing length sweeps use the widths $\Delta z^t$ instead of $\Delta z$.
The surface value of \verb|mxl_choice| $=1$ is limited by $\ell_{min} + z_0 + \Delta z^t$ of the surface cell.
The vertical shear and $N^2$ use the distances $\Delta z^w$.
The surface $W$ cell has the thickness $\Delta z^w/2$ of the surface face.

The shear production $K_{diss}$ is computed at the $U$ and $V$ points and averaged to the $T$ points.
The average is weighted with the areas, as in pyOM2,
\beq
 K^t_i = \frac{A^u_i K^u_i + A^u_{i-1} K^u_{i-1}}{2 A^t_i}
\eeq
and correspondingly in $y$ with $A^v$.
An average without the areas would change $K_{diss}$ by 6\% at $60^\circ$ latitude.

\subsection*{Isoneutral mixing}

Isoneutral mixing is switched on by \verb|enable_neutral_diffusion|.
It is a port of \verb|isoneutral.f90| and \verb|isoneutral_diffusion.f90| of pyOM2.
It uses the triads of Griffies et al., as pyOM2 does.

The isoneutral slopes follow from the gradients of temperature and salinity and the derivatives of density.
For \verb|eq_of_state| $=0$ the derivative with respect to temperature is $-\rho_0/g$.
pyOM2 uses zero there.
For the value 1 the derivatives are the constants of the linear equation of state.
For the value 5 they are the derivatives of the TEOS-10 equation of state.
The value 100 is not possible.

Steep slopes are tapered with
\beq
 \tau(s) = \frac{1}{2} \lk 1 + \tanh \frac{s_c - |s|}{\Delta s} \rk
\eeq
with $s_c$ = \verb|iso_slopec| (default $10^{-3}$) and $\Delta s$ = \verb|iso_dslope| (default $8 \cdot 10^{-4}$).

The diffusivity is the three-dimensional field \verb|K_iso|.
It is set to the constant \verb|K_iso_0| at the start unless the configuration has set the field.
The horizontal diffusivity in the tensor is $\max(K_{steep}, K_{iso}\, \tau)$ with $K_{steep}$ = \verb|K_iso_steep|.
The field can be changed during the run, for instance by another parameterisation.

With \verb|enable_skew_diffusion| the skew flux of Gent and McWilliams is added, as in pyOM2.
Its diffusivity is the three-dimensional field \verb|K_gm|.
It is set to the constant \verb|K_gm_0| at the start unless the configuration has set the field.
On the lateral faces the triads then use $K_{iso} - K_{gm}$.
On the top faces they use $K_{iso} + K_{gm}$.

The horizontal and mixed parts of the tensor are applied with a forward Euler step.
The vertical part $K_{33}$ is applied with the implicit vertical diffusion of the model.
Isoneutral mixing acts on temperature and, for an equation of state with salinity, on salinity.

At the start the model checks the stability of the explicit part, as pyOM2 does.
The quantity
\beeq
 \delta = \min \lk \frac{\cos \phi^t \Delta x^t \Delta z^t}{4 K_{iso,0} \Delta t},\ \frac{\Delta y^t \Delta z^t}{4 K_{iso,0} \Delta t} \rk
\eeeq
is the steepest slope that is stable.
The model stops with an error if $\delta$ is smaller than $s_c$.

The following parts of pyOM2 are not ported.
These are the dissipation diagnostics of the isoneutral and skew fluxes, the treatment in the mixed layer and the friction of the TEM formulation.

\subsection*{Tracer advection outside the Adams--Bashforth step}

With thin cells on a fine horizontal grid the vertical Courant number $|w| \Delta t/\Delta z^t$ can limit the time step.
The Superbee scheme is stable up to a Courant number of 1 in a forward step.
Inside the Adams--Bashforth step with the weights $1.6$ and $-0.6$ it becomes unstable above about $0.45$.
The attribute \verb|vadv_sub_cycling| $= M$ therefore takes the tracer advection out of the Adams--Bashforth step.
The default $M = 0$ keeps the advection in the Adams--Bashforth step.

With $M > 0$ the advection uses the velocities $u$, $v$ and $w$ at the end of the time step.
It consists of a step in $x$, a step in $y$ and $M$ steps in $z$ with the time step $\Delta t/M$.
Each of them is a forward step of the Superbee scheme, as in MITgcm.
The vertical Courant number may then be as large as $M$.
The Adams--Bashforth step adds only the forcing of the tracers.
The option needs \verb|superbee_advection|.
It cannot be combined with \verb|weno_advection|, and the model stops with an error then.
The steps outside the Adams--Bashforth step are forward steps of the Superbee scheme.
The WENO scheme is written for the Adams--Bashforth step.

All of the advection has to leave the Adams--Bashforth step, not the vertical part alone.
Horizontal and vertical advection cancel in a flow along isopycnals.
A step in which one of them is forward and the other is part of the Adams--Bashforth step is unstable in such a flow.
For the same reason the vertical advection of momentum stays in the Adams--Bashforth step.

The velocities at the end of the time step are needed.
With the velocities at the start of the time step the forward tracer step and the Adams--Bashforth step of momentum amplify internal waves.

A step in one direction changes a uniform tracer, since the flow in that direction alone is not free of divergence.
Each step therefore carries the volume $\varrho$ that the cell would have after the steps done so far (Easter 1993, Lin and Rood 1996).
It starts with $\varrho = 1$ and changes in each step in $x$ as
\beeq
 \varrho \leftarrow \varrho - \frac{\Delta t}{\cos \phi^t \Delta x^t} \lk u_i - u_{i-1} \rk
\eeeq
and correspondingly in $y$ and $z$.
The fluxes are formed with the tracer value $T = q/\varrho$ and change the tracer content $q$.
The content is conserved to round-off, and a uniform tracer stays uniform.

The flux through the surface face is handled in the steps in $z$.
It is zero, or $w\, T$ of the surface cell with \verb|free_surface_tracer_flux|.
Its mean over the $M$ steps is kept in \verb|vadv_fs_flux| for a configuration that books the tracer budget.

The option adds no new state to the restart.
The tendencies \verb|dtemp| and \verb|dsalt| then contain the forcing only.
The first step after switching the option on or off should therefore not use the old tendencies.

\subsection*{Time step}

The order of the calls in a time step of the split-explicit scheme is
\begin{enumerate}
\item diffusivities of the TKE closure (\verb|set_tke_diffusivities|)
\item momentum tendencies, Coriolis force and hydrostatic pressure (\verb|momentum_tendency_expl|)
\item forcing of the configuration (\verb|set_forcing|)
\item bottom friction
\item external mode with sub cycling (\verb|solve_expl_surface|)
\item vertical friction
\item vertical velocity
\item tracer advection and biharmonic tracer mixing, or the advection outside the Adams--Bashforth step
\item Adams--Bashforth update of the tracers
\item isoneutral mixing
\item vertical mixing of the tracers
\item density
\item TKE step (\verb|integrate_tke|)
\end{enumerate}

The order of the calls in a time step of the rigid lid and of the implicit free surface is
\begin{enumerate}
\item diffusivities of the TKE closure
\item tracer advection and biharmonic tracer mixing
\item momentum tendencies and hydrostatic pressure (\verb|momentum_tendency_impl|)
\item forcing of the configuration (\verb|set_forcing|)
\item biharmonic and harmonic momentum friction
\item Adams--Bashforth update of the velocities
\item direct forcing of the configuration (\verb|set_forcing_direct|)
\item surface pressure and its gradient (\verb|solve_for_pressure|)
\item vertical friction
\item bottom friction
\item vertical velocity
\item Adams--Bashforth update of the tracers, or the advection outside the Adams--Bashforth step
\item isoneutral mixing
\item vertical mixing of the tracers
\item density
\item TKE step
\end{enumerate}
In both schemes the steps are those of pe\_simple, with the isoneutral mixing, the harmonic friction and the advection outside the Adams--Bashforth step added.

\section*{Model parameters and options}

The parameters and options of pe\_simple are also those of pe\_realistic.
The tables list only the differences.

\subsection*{Required parameters}

\begin{longtable}{@{}llp{8.4cm}@{}}
\hline
Attribute & Unit & Meaning \\
\hline
\endhead
\verb|dxt|, \verb|dyt| & m or $^\circ$ & widths of the cells in $x$ and $y$, a number or an array without halo.
  Degrees with \verb|coord_degree|. Replace \verb|dx| and \verb|dy| of pe\_simple, which are still accepted \\
\verb|dzt| & m & widths of the cells in $z$, a number or an array without halo, from the bottom to the surface.
  Replaces \verb|dz| of pe\_simple, which is still accepted \\
\verb|eq_of_state| & --- & equation of state, $0$, $1$, $5$ (TEOS-10) or $100$ \\
\hline
\end{longtable}

\noindent
The configuration must set the Coriolis parameter with \verb|set_coriolis|.
The field \verb|coriolis_t| is zero otherwise.

\subsection*{Options with defaults}

\begin{longtable}{@{}llp{7.6cm}@{}}
\hline
Attribute & Default & Meaning \\
\hline
\endhead
\verb|coord_degree| & off & longitude and latitude in degrees instead of Cartesian coordinates \\
\verb|x_origin|, \verb|y_origin| & $0$ & position of the first interior $U$ point \\
\verb|radius| & $6370\,$km & radius of the earth (m) \\
\verb|set_zero_at_surface| & off & origin of $z$ at the surface face instead of the bottom face \\
\verb|Ah| & unset & harmonic momentum friction (m$^2$/s), active if present.
  Rigid lid and implicit free surface only \\
\verb|hor_friction_cosPower| & $0$ & exponent $p$ of $(\cos \phi)^p$ in the harmonic friction \\
\verb|biha_friction_cosPower| & $0$ & exponent $p$ of $(\cos \phi)^p$ in each pass of the biharmonic friction \\
\verb|enable_neutral_diffusion| & off & isoneutral mixing \\
\verb|enable_skew_diffusion| & off & skew flux of Gent and McWilliams \\
\verb|K_iso_0| & $0$ & isoneutral diffusivity (m$^2$/s), initial value of \verb|K_iso| \\
\verb|K_gm_0| & $0$ & skew diffusivity (m$^2$/s), initial value of \verb|K_gm| \\
\verb|K_iso_steep| & $0$ & lower limit of the tapered horizontal diffusivity (m$^2$/s) \\
\verb|iso_slopec| & $10^{-3}$ & slope at which the taper is one half \\
\verb|iso_dslope| & $8 \cdot 10^{-4}$ & width of the taper \\
\verb|K_iso|, \verb|K_gm| & unset & three-dimensional diffusivities, set from \verb|K_iso_0| and \verb|K_gm_0| if not given \\
\verb|vadv_sub_cycling| & $0$ & number $M$ of vertical steps of the tracer advection outside the Adams--Bashforth step.
  $0$ keeps the advection in the Adams--Bashforth step \\
\verb|weno_advection| & off & as in pe\_simple, with the coefficients of a uniform grid.
  Not together with \verb|vadv_sub_cycling| \\
\verb|weno_ref| & $1$ & as in pe\_simple \\
\hline
\end{longtable}

\subsection*{Snapshot diagnostics}

The snapshot classes \verb|snap_nc4| and \verb|snap_raw| are those of pe\_simple.
They write the coordinates of the grid above.
These are $x^t$, $x^u$, $y^t$, $y^u$, $z^t$ and $z^w$, in degrees with \verb|coord_degree|.
$z^w$ is written under the name \verb|zu|.

\subsection*{Limits}

\begin{itemize}
\item The grid must not include the poles. The zonal width $\cos \phi\, \Delta x$ goes to zero there, and $\tan \phi$ grows without bound.
\item The rigid lid does not converge on a domain that is periodic in $x$ and contains land.
\item TEM friction, the EKE model and IDEMIX of pyOM2 are not part of pe\_realistic.
\item With varying $\Delta x^t$ in a periodic direction pe\_realistic and pyOM2 build different $x^t$ and $\Delta x^u$ in the halo.
\item $\Delta z^t$ must vary smoothly, see the grid above.
\item \verb|weno_advection| and \verb|vadv_sub_cycling| exclude each other.
\end{itemize}

\end{document}

\clearpage\phantomsection\label{doc:nh_simple}\addcontentsline{toc}{section}{nh\_simple}
\documentclass[12pt]{article} 
\usepackage{natbib}
\usepackage{a4wide}
\usepackage{lineno}
\usepackage{cancel}
\usepackage{longtable}

\newif\ifbuch
\global\buchtrue
\global\buchfalse

\newif\ifdetail
\global\detailfalse
%\global\detailtrue

\usepackage{color}
\usepackage[pdftex]{graphicx}
\usepackage[pdftex,colorlinks,
                citecolor=darkgreen,linkcolor=darkblue,urlcolor=darkblue]{hyperref}
\definecolor{darkgreen}{rgb}{0.1,0.4,0.1}
\definecolor{darkblue}{rgb}{0.1,0.1,0.3}

\usepackage{amssymb}

\newcommand{\beq}  { \begin{eqnarray*} }
\newcommand{\eeq}  { \end{eqnarray*}}
\newcommand{\beeq}  { \begin{eqnarray*} }
\newcommand{\eeeq}  { \end{eqnarray*}}
\newcommand{\eq }  [1]{Eq.~(\ref{#1})}  % Eq. (#1)
\newcommand{\sect}  [1]{Sec.~(\ref{#1})}  % Eq. (#1)
\newcommand{\fig}  [1]{Fig.~\ref{#1}}   % Fig. #1
\renewcommand{\v}[1]{{\mbox{\boldmath$ #1 $}}}  % Vektor durch Fettdruck
\newcommand{\m}    [1]{{\bf  #1 }}              % Matrix in sans serif
\newcommand{\gz}[1] {{{d #1} \over {dz}}}
\newcommand{\pt}[1]{{\partial_t #1}}
\newcommand{\ptt}[1]{{{\partial^2 #1} \over {\partial t^2}}}
\newcommand{\pz}[1]{{\partial_z #1}}
\newcommand{\pzz}[1]{{\partial_{zz} #1}}
\newcommand{\px}[1]{{\partial_x #1}}
\newcommand{\py}[1]{{\partial_y #1}}
\newcommand{\p}{{\partial}}
\newcommand{\vn}{{\v \nabla}}
\newcommand{\tl}{\tilde}
\newcommand{\lk}{\left(}
\newcommand{\rk}{\right)}
\newcommand{\rvec} [1]{
    \raisebox{-1.5ex}{$\stackrel{\textstyle #1}{\neg}$} }% rotierter Vektor
\newcommand{\Vec}[1]{\left(\begin{array}{c} #1 \end{array}\right) }
\newcommand{\Bar}[1]{ \overline{ #1} }

\begin{document} 

\section*{Documentation for the simple non-hydrostatic model nh\_simple}

nh\_simple is a finite-difference sibling of the pseudo-spectral model ps3D:
it integrates the same scaled non-hydrostatic Boussinesq equations on the
same kind of domain (horizontally always periodic, vertically periodic or
bounded by rigid lids), but discretises them on an Arakawa C grid in
physical space and uses the Fourier transform only for the pressure
inversion. The LES closures are ported from ps3D. Within pyOM3, nh\_simple
is one of the ocean components that can be coupled to the spectral surface
wave model wim\_simple; the hooks of the Lagrangian-mean formulation of
that coupling are noted where they appear.

\subsection*{Continuous equations}

The scaled non-hydrostatic Boussinesq equations (see the scaling section
of the ps3D documentation for the derivation) are
\beq
 \pt{\v u } + f \rvec{\v u}  +   \vn p =   
  - Ro \lk \v u \cdot \vn \v u +  w \pz{} \v u \rk
 ~~,~~
    \pt{b}  + N^2  w    =  -Ro \lk  \v u \cdot \vn b  +  w \pz{} b \rk
 \\
   \delta^2 \pt w +   \pz p -  b =  - Ro \, \delta^2  \lk   \v u \cdot \vn w +  w \pz{} w \rk
\eeq
with the diagnostic relation $\vn \cdot \v u + \pz w = 0$. Vectors are
horizontal, $\rvec{\v u}$ denotes anticlockwise rotation by $90^o$,
$Ro$ is the Rossby number and $\delta = H/L$ the aspect ratio;
$f$ and the background stability frequency $N$ are constants ($f$-plane,
uniform stratification). The buoyancy $b$ is the deviation from the
background stratification $N^2 z/Ro$; the term $N^2 w$ is the advection of
this background. Kinetic and potential energy are
$\v u^2/2 + \delta^2 w^2/2$ and $b^2/2N^2$. Setting $Ro = 1$ and
$\delta^2 = 1$ recovers the plain dimensional equations, which is the
configuration required by the LES closures below.

\subsection*{Grid and spatial discretisation}

All fields live on an Arakawa C grid with uniform spacings
$\Delta x, \Delta y, \Delta z$: $b$ and $p$ at cell centres, $u,v,w$ on the
respective cell faces. The momentum advection is the second-order
energy-conserving flux form, e.g.\ for $u$
\beq
 \lk \pt u \rk_{adv} = -Ro \lk \delta_x^- \Bar{\Bar u^{x+} \Bar u^{x+}}
 + \delta_y^- \Bar{ \Bar u^{y+} \Bar v^{x+}}
 + \delta_z^- \Bar{ \Bar u^{z+} \Bar w^{x+}} \rk
\eeq
in the finite-difference and averaging operator notation of the ps3D
documentation ($\delta^\pm$ forward/backward differences, $\Bar{\cdot}^\pm$
two-point averages), and analogously for $v$ and $w$ (the $w$ tendency
carries the factor $\delta^2$). The Coriolis force uses the standard
four-point average of the C grid.

For the buoyancy advection three flux-form schemes are available, selected
by attribute: the second-order centred scheme (default, energy conserving),
the second-order scheme with Superbee flux limiter
(\verb|buoyancy_superbee_advection|, from MITgcm; monotone, and its
implicit scalar mixing can replace the explicit SGS buoyancy diffusion of
the LES closures, see below), and the third-order upwind scheme
(\verb|buoyancy_upwind3_advection|, from MITgcm). Momentum advection is
always the energy-conserving second-order form.

\subsection*{Pressure solver}

The velocity is first advanced with all terms except the pressure gradient
to an intermediate value $\v u^*, w^*$; the pressure then enforces the
divergence constraint,
\beq
 \lk \delta_x^- \delta_x^+ + \delta_y^- \delta_y^+
   + \delta_z^- \delta_z^+/\delta^2 \rk p =
 \frac{1}{\Delta t} \lk \delta_x^- u^* + \delta_y^- v^* + \delta_z^- w^* \rk
\eeq
followed by the projection
$\v u = \v u^* - \Delta t\, \vn^+ p$ and
$w = w^* - (\Delta t/\delta^2)\, \delta_z^+ p$.
As shown in the ps3D documentation, the Fourier transform of the discrete
one-dimensional operator $\delta^-\delta^+$ is
$-\hat k^+ \hat k^- = -2\lk 1 - \cos k \Delta \rk/\Delta^2$, which
approaches $-k^2$ for $k \Delta \to 0$; the discrete Poisson equation is
therefore diagonal in spectral space with the eigenvalue
$\hat k^2 = \hat k_x^+ \hat k_x^- + \hat k_y^+ \hat k_y^-
+ \hat k_z^+ \hat k_z^-/\delta^2$,
so the discrete divergence of the projected velocity vanishes to round-off.

Two cases are implemented, selected by \verb|periodic_in_z|. For the
triple-periodic domain a 3D FFT (parallel, through the p3dfft library)
inverts the full eigenvalue; the mean mode $\hat k^2 = 0$ is set to zero.
For rigid lids at top and bottom, a 2D FFT in $x,y$ is combined with a
tridiagonal solve in $z$ (Thomas algorithm) with Neumann boundary
conditions at the lids; the horizontal-mean column
($\hat k_x = \hat k_y = 0$), for which the tridiagonal system is singular,
is integrated directly by the second-order recurrence
$p_{k+1} = 2 p_k - p_{k-1} + \delta^2 \Delta z^2 f_k$ with $p_0 = 0$.

\subsection*{Boundary conditions}

The domain is always periodic in $x$ and $y$ (halo wrap or MPI exchange).
In $z$ it is either periodic (\verb|periodic_in_z| true) or bounded by
rigid free-slip lids: all vertical fluxes --- advective, frictional,
diffusive and the LES fluxes --- are set to zero on the top and bottom
faces, and $w$ vanishes there.

\subsection*{Friction and mixing}

Harmonic friction with constant coefficients acts on the momentum when the
attributes are present: with both \verb|Ah| and \verb|Av| the full operator
$A_h \vn_h^2 + A_v \p_{zz}$, with only one of them the corresponding
horizontal-only or vertical-only part. Harmonic buoyancy mixing requires
both \verb|Kh| and \verb|Kv|. Biharmonic friction and mixing
($-A_{hbi}\vn_h^4$ etc.) are activated by the pairs \verb|Ahbi|,
\verb|Avbi| for momentum and \verb|Khbi|, \verb|Kvbi| for buoyancy. All
these operators are applied with a forward Euler step outside the
Adams--Bashforth interpolation of the adiabatic tendencies, so their
coefficients are subject to the usual explicit diffusive stability limits.

\subsection*{Time stepping}

The adiabatic tendencies (advection, Coriolis, buoyancy force) are
interpolated with the multi-step Adams--Bashforth scheme,
\beq
(c_A, c_B, c_C) = (1,0,0) ~\to~ \lk \frac{3}{2}, -\frac{1}{2}, 0\rk
 ~\to~ \lk \frac{23}{12}, -\frac{16}{12}, \frac{5}{12} \rk
\eeq
i.e.\ a forward Euler start, one AB2 step and third-order Adams--Bashforth
from the third step on, with the tendencies held in a rotating three-level
buffer. Friction, mixing and the LES closures update the fields directly
with forward Euler before the Adams--Bashforth increment; the pressure
projection concludes the step. A user-defined forcing can be added to the
tendencies at the current time level through the optional method
\verb|set_forcing(tau)| (it then enters the Adams--Bashforth
interpolation), or applied as a direct increment through
\verb|set_forcing_direct()| after the Adams--Bashforth update and before
the pressure solve. The latter bypasses the AB weights and is the hook for
the Stokes tendency $\Delta \v u_S$ of the Lagrangian-mean ($\v u_L$)
formulation of the wave coupling, which must be a plain forward increment
and must precede the projection so that it acts on its divergence.

\subsection*{LES closures}

The LES closures of ps3D (\verb|les_closures.f90|) are ported as the mixin
class \verb|les| (module \verb|les.py|): the Smagorinsky closure, the
filtered-structure-function closure of Ducros et al.\ (1996), and a
Deardorff-type prognostic TKE closure, plus an optional Peclet bound on
the horizontal viscosities. The closures are formulated dimensionally and
refuse to run unless $Ro = 1$ and $\delta^2 = 1$. One call of
\verb|les_closure()| per time step (after the harmonic friction, before
the Adams--Bashforth increment, e.g.\ hooked through \verb|set_forcing|)
computes the strain-rate tensor on its natural C-grid positions,
evaluates the enabled closure(s), and applies the anisotropic viscous and
diffusive flux divergences to $u,v,w$ and $b$ with a forward Euler step.

\emph{Smagorinsky} (\verb|enable_smagorinsky|):
$\nu = (c_s \Delta)^2 \sqrt{2 S_{ij} S_{ij}}$ with
$\Delta = (\Delta x \Delta y \Delta z)^{1/3}$ and
$c_s = $ \verb|smag_coeff|; the diffusivity is $\nu/Pr$ with the turbulent
Prandtl number \verb|prandtl|. The viscosity is evaluated separately at
the four staggered positions where its fluxes are needed (cell centre,
horizontal corner, above $u$, above $v$).

\emph{Ducros et al.\ (1996)} (\verb|enable_ducros_96|): the viscosity is
built from the second-order structure function of the high-pass filtered
velocity (two iterated horizontal Laplace filters, three with
\verb|enable_three_level_filter|).

\emph{Prognostic TKE} (\verb|enable_tke_closure|): a Deardorff-type
closure with prognostic subgrid TKE on the $w$ grid, a
stratification-limited length scale
$\ell = \min(\Delta,\, 0.76 \sqrt{e/N^2})$, dissipation
$(0.19 + 0.51 \ell/\Delta)\, e^{3/2}/\ell$, TKE diffusion with
coefficient \verb|c_eke| times the eddy diffusivity, and shear and
buoyancy production diagnosed from the actual discrete mixing fluxes of
the previous step. Viscosity and diffusivity are
$\nu = c_m \ell \sqrt{e}$ and $\kappa = (1 + 2\ell/\Delta)\, \nu$ with
$c_m = $ \verb|c_m|; $N^2$ always includes the background $N_0^2$.

\emph{Stability caps.} All closures cap their viscosities and
diffusivities against the explicit diffusive limit of the forward Euler
update. The caps are direction-aware: each staggered viscosity is capped
with the most restrictive grid spacing among the directions its fluxes
act in (e.g.\ the cell-centre viscosity by $\min(\Delta x, \Delta y,
\Delta z)^2/(4\Delta t)$). This deviates deliberately from the original
ps3D caps, which use geometric means of the spacings and are valid only
for near-isotropic grids: for $\Delta x \gg \Delta z$ they admit
vertically unstable viscosities.

\emph{Boundary treatment.} With rigid lids the vertical shear seen by the
closures is zeroed on the lids (\verb|les_freeslip_strain|, default on),
consistent with the free-slip zero-flux condition of the mixing;
otherwise the zero halo above the lid produces a spurious shear that
inflates the viscosity in the top and bottom cells.

\emph{Scalar mixing.} The explicit SGS buoyancy diffusion can be switched
off (\verb|les_mix_b|) when a flux-limited buoyancy advection scheme
(Superbee) provides the scalar SGS mixing implicitly; under the TKE
closure the resolved buoyancy production is then zero.

\emph{Peclet bound} (\verb|enable_peclet_bound|): where the grid Peclet
number of the resolved flow exceeds \verb|Peclet_limit|, the horizontal
viscosities are raised proportionally. The bound uses the prognostic
(advecting) velocity, i.e.\ $\v u_L$ in the Lagrangian-mean formulation.

\emph{Coupled wave--ocean setups.} When the prognostic velocity is the
Lagrangian mean $\v u_L$, the attribute \verb|les_mix_eulerian| makes the
strains and mixing fluxes act on the Eulerian velocity
$\v u = \v u_L - \v u_S$ (attributes \verb|u_S|, \verb|v_S| and optionally
\verb|w_S| supplied by the coupler), with the tendency added to the
prognostic $\v u_L$ --- the down-gradient closure parameterises Eulerian
fluctuations, so mixing $\v u_L$ would spuriously diffuse the Stokes
shear. This is the LES counterpart of the \verb|mix_eulerian| option of
the TKE closure in pe\_simple.

\subsection*{Domain decomposition}

The parallelisation follows the pencil decomposition of the p3dfft
library used by the pressure solver: the $x$ direction is kept local on
every processor, and the constructor arguments \verb|n_pes_x| and
\verb|n_pes_y| define the processor grid over the $y$ and $z$ directions
(internally the blocks are addressed as $y$/$z$ blocks). Model arrays are
Fortran-ordered $(N_x, N_y, N_z)$ with halo points on all sides; the
pressure solver requires this ordering. Halo exchanges wrap periodically
in $x$ and $y$; in $z$ they wrap only for the triple-periodic case.

\subsection*{Snapshot diagnostics}

Two snapshot classes live in the module \verb|snapshot.py|; either
replaces the model class as the base of the configuration and overrides
\verb|setup| and \verb|diagnose|. Both read the attribute \verb|snapint|
(snapshot interval in seconds, required), take the setup arguments
\verb|snap_file_name| and \verb|show_each_time_step| (default on: print
the model time every step; off: only at output steps), and write at
$n = 0$ and whenever $n \bmod (\mathrm{snapint}/\Delta t) = 0$. Both
record the variables $u,v,w,b,p$ on their staggered coordinate axes
(\verb|xt|/\verb|yt|/\verb|zt| for cell centres and
\verb|xu|/\verb|yu|/\verb|zu| for the faces, shifted by half a spacing);
the vertical axis is negative downwards with the surface at $z = 0$ and
the cell centres at $z = -\Delta z/2, \ldots$. Halo points are stripped
before writing.

\emph{snap\_nc4} writes a single netCDF file (default \verb|'snap.cdf'|).
At start-up it detects whether the netCDF4/HDF5 build supports parallel
I/O (\verb|__has_parallel4_support__| or pnetcdf) and reports the result.
In parallel mode the file is created and opened collectively by all PEs,
and every PE writes its subdomain into the global arrays in a single
collective pass per snapshot. Without parallel support the original
serialised path is used: PE0 creates the file and extends the time
dimension, then each PE in turn opens the file and writes its block,
synchronised with barriers --- correct on any build, but increasingly
slow with many PEs. The internal Fortran ordering is transposed on the
fly, so the file layout (time, $z$, $y$, $x$) is the same in both modes.

\emph{snap\_zarr\_gather} writes a single zarr store (default
\verb|'snap.zarr'|) with field arrays of shape (time, $n_x$, $n_y$,
$n_z$), always $x$-first regardless of the internal ordering, one time
slice per chunk. The design goal is a Lustre-friendly file count: with
the default \verb|z_chunks| $= 1$ each snapshot adds one chunk file per
field --- of the order of six files per snapshot regardless of the number
of PEs --- instead of the one-file-per-PE-per-field fragmentation of a
per-block chunking, which can overwhelm the metadata server of a parallel
file system at large PE counts. Every PE participates in the collective
gather of each field (\verb|gather_array_mpi|), but only PE0 touches the
store, writing a full horizontal slab of $\lceil n_z/\verb|z_chunks|
\rceil$ levels per chunk; increase \verb|z_chunks| only if a single field
slab is too large to hold and write in one piece. Fields are gathered and
written one at a time to bound the PE0 memory to one global field. The
time axis and the field arrays are resized (appended) per snapshot, so
the store remains valid after an interrupted run; an existing store of
the same name is removed at start-up. The on-disk format is pinned to
zarr v2 even when a zarr-python v3 library writes it, so the output stays
readable by v2 environments; the cosmetic consolidated-metadata warning
of v3 is silenced. Compared to the netCDF path this trades the gather
(one global field on PE0) for independence from parallel-I/O library
support.

\section*{Model parameters and options}

All parameters and options are attributes of the model instance, set
either through the parameter dictionary passed at construction or by
assigning the attribute before the time loop; options with a default may
be left unset.

\subsection*{Required parameters (no default)}

\begin{longtable}{@{}llp{8.4cm}@{}}
\hline
Attribute & Unit & Meaning \\
\hline
\endhead
\verb|nx|, \verb|ny|, \verb|nz| & --- & global grid size (cells, without
  halo points) \\
\verb|dx|, \verb|dy|, \verb|dz| & m & grid spacing (uniform) \\
\verb|dt| & s & time step \\
\verb|f0| & 1/s & Coriolis parameter ($f$-plane) \\
\verb|N0| & 1/s & background stability frequency \\
\verb|Ro| & --- & Rossby number of the scaled equations ($1$ for the
  dimensional form; the LES closures require $1$) \\
\verb|dsqr| & --- & squared aspect ratio $\delta^2$ ($1$ for the fully
  non-hydrostatic dimensional form; the LES closures require $1$) \\
\verb|periodic_in_z| & --- & triple-periodic domain if true; rigid
  free-slip lids at top and bottom if false. $x$ and $y$ are always
  periodic \\
\hline
\end{longtable}

\noindent
The constructor additionally takes \verb|n_pes_x|, \verb|n_pes_y|
(processor grid over the $y$ and $z$ pencils, default $1 \times 1$;
their product must equal the communicator size), \verb|halo_points|
(default 1) and \verb|comm_in| (default: MPI world communicator).

\subsection*{Forcing hooks and coupling fields}

\begin{longtable}{@{}lp{10.6cm}@{}}
\hline
Attribute / method & Meaning \\
\hline
\endhead
\verb|set_forcing(tau)| & optional user method; adds forcing to the
  tendencies at time level \verb|tau| (enters the Adams--Bashforth
  interpolation). Also the natural hook to call \verb|les_closure()| \\
\verb|set_forcing_direct()| & optional user method; direct increment after
  the Adams--Bashforth update and before the pressure solve, bypassing the
  AB weights. Used for the Stokes tendency $\Delta \v u_S$ of the
  Lagrangian-mean ($\v u_L$) formulation \\
\verb|u_S|, \verb|v_S|, \verb|w_S| & Stokes velocity from the wave model
  ($w_S$ optional); required by \verb|les_mix_eulerian| \\
\hline
\end{longtable}

\subsection*{Options with defaults}

\begin{longtable}{@{}llp{7.9cm}@{}}
\hline
Attribute & Default & Meaning \\
\hline
\endhead
\verb|buoyancy_superbee_advection| & off & buoyancy advection with the
  Superbee flux limiter instead of the centred second-order scheme \\
\verb|buoyancy_upwind3_advection| & off & third-order upwind buoyancy
  advection \\
\verb|Ah|, \verb|Av| & unset & harmonic viscosities (m$^2$/s); either one
  alone activates the horizontal-only / vertical-only operator \\
\verb|Kh|, \verb|Kv| & unset & harmonic buoyancy diffusivities (m$^2$/s);
  both must be present \\
\verb|Ahbi|, \verb|Avbi| & unset & biharmonic viscosities (m$^4$/s); both
  must be present \\
\verb|Khbi|, \verb|Kvbi| & unset & biharmonic buoyancy diffusivities
  (m$^4$/s); both must be present \\
\hline
\end{longtable}

\subsection*{LES options (read at \texttt{init\_les})}

\begin{longtable}{@{}llp{7.9cm}@{}}
\hline
Attribute & Default & Meaning \\
\hline
\endhead
\verb|enable_smagorinsky| & off & Smagorinsky closure \\
\verb|enable_ducros_96| & off & Ducros et al.\ (1996)
  structure-function closure \\
\verb|enable_tke_closure| & off & Deardorff-type prognostic TKE closure \\
\verb|enable_three_level_filter| & off & three instead of two iterated
  Laplace filters (Ducros closure) \\
\verb|enable_peclet_bound| & off & raise horizontal viscosities where the
  grid Peclet number exceeds \verb|Peclet_limit| \\
\verb|les_mix_eulerian| & off & strains and mixing act on the Eulerian
  velocity $\v u_L - \v u_S$ ($\v u_L$ formulation) \\
\verb|les_freeslip_strain| & on & zero the vertical shear seen by the
  closure at the rigid lids \\
\verb|les_mix_b| & on & explicit SGS buoyancy diffusion; switch off when a
  flux-limited advection scheme provides the scalar SGS mixing \\
\verb|smag_coeff| & $0.148$ & Smagorinsky coefficient \\
\verb|prandtl| & $0.47$ & turbulent Prandtl number \\
\verb|min_smag| & $0$ & viscosity floor (m$^2$/s) \\
\verb|Peclet_limit| & $50$ & grid Peclet number bound \\
\verb|c_m| & $0.1$ & TKE-closure viscosity constant \\
\verb|c_eke| & $2$ & TKE diffusion constant \\
\verb|tke_min| & $10^{-8}$ & TKE floor (m$^2$/s$^2$) \\
\verb|les_epsln| & $10^{-20}$ & numerical floor of the closure \\
\hline
\end{longtable}

\subsection*{Snapshot diagnostics}

Both snapshot classes (module \verb|snapshot.py|) read the attribute
\verb|snapint| (snapshot interval in seconds, required) and take the setup
arguments \verb|snap_file_name| (defaults \verb|'snap.cdf'| /
\verb|'snap.zarr'|) and \verb|show_each_time_step| (default on);
\verb|snap_zarr_gather| additionally takes \verb|z_chunks| (default 1:
one chunk file per field and snapshot; larger values split the vertical
into $\lceil n_z/\verb|z_chunks| \rceil$-level chunks).

\end{document}

\clearpage\phantomsection\label{doc:nh_realistic}\addcontentsline{toc}{section}{nh\_realistic}
\documentclass[12pt]{article}
\usepackage{lineno}
%\usepackage{ametsoc}
\usepackage{natbib}
\usepackage{a4wide}
\usepackage{longtable}
\usepackage{amsfonts}
\usepackage{amssymb}

\usepackage{color}
\usepackage[pdftex]{graphicx}
\usepackage[pdftex,colorlinks, citecolor=black,linkcolor=black,urlcolor=black]{hyperref}
\definecolor{darkgreen}{rgb}{0.1,0.4,0.1}
\definecolor{darkblue}{rgb}{0.1,0.1,0.3}



\newif\ifdetail
\global\detailfalse
%\global\detailtrue

\def\bem#1{\colorbox{green}{\em #1}}

\newcommand{\beq}  { \begin{eqnarray} }
\newcommand{\eeq}  { \end{eqnarray}}
\newcommand{\beeq}  { \begin{eqnarray*} }
\newcommand{\eeeq}  { \end{eqnarray*}}
\newcommand{\eq }  [1]{Eq.~(\ref{#1})}  % Eq. (#1)
\newcommand{\eqs }  [1]{Eqs.~(\ref{#1})}
\newcommand{\sect}  [1]{Sec.~(\ref{#1})}  % Eq. (#1)
\newcommand{\fig}  [1]{Fig.~\ref{#1}}   % Fig. #1
\renewcommand{\v}[1]{{\mbox{\boldmath$ #1 $}}}  % Vektor durch Fettdruck
\newcommand{\m}    [1]{{\bf  #1 }}              % Matrix in sans serif
\newcommand{\gz}[1] {{{d #1} \over {dz}}}
\newcommand{\pt}[1]{{\partial_t #1}}
\newcommand{\ptt}[1]{{\partial_{tt} #1}}
\newcommand{\pttt}[1]{{\partial_{ttt} #1}}
\newcommand{\pttx}[1]{{{\partial^3 #1} \over {\partial t^2 \partial x}}}
\newcommand{\ptty}[1]{{{\partial^3 #1} \over {\partial t^2 \partial y}}}
\newcommand{\ptx}[1]{{{\partial^2 #1} \over {\partial t \partial x}}}
\newcommand{\pty}[1]{{{\partial^2 #1} \over {\partial t \partial y}}}
\newcommand{\pz}[1]{{\partial_z #1}}
\newcommand{\pzz}[1]{{\partial_{zz} #1}}
\newcommand{\px}[1]{{\partial_x #1}}
\newcommand{\py}[1]{{\partial_y #1}}
\newcommand{\pxi}[1]{{\partial_i #1}}
\newcommand{\pxj}[1]{{\partial_j #1}}
\newcommand{\pxk}[1]{{\partial_k #1}}
\newcommand{\pb}[1]{{{\partial #1} \over {\partial b}}}
%\newcommand{\pxy}[1]{{{\partial^2 #1} \over {\partial x \partial y}}}
\newcommand{\pxy}[1]{\partial_{xy} #1}
\newcommand{\pyy}[1]{\partial_{yy} #1}
\newcommand{\pxx}[1]{\partial_{xx} #1}
\newcommand{\kx}{{\bf k} \times}
\newcommand{\dt}[1]{{{d #1} \over {d t}}}
\newcommand{\Dt}[1]{{{D #1} \over {D t}}}
\newcommand{\Div}[1]{ \nabla \cdot #1 }
\renewcommand{\div}[1]{ \nabla_{yz} \cdot #1 }
\newcommand{\Bar}[1]{ \overline{ #1} }
\newcommand{\rvec} [1]{
    \raisebox{-1.5ex}{$\stackrel{\textstyle #1}{\neg}$} }% rotierter Vektor
\newcommand{\nablq}   {\rvec\nabla}                    % rotated nabla operator
\newcommand{\Vec}[1]{\left(\begin{array}{c} #1 \end{array}\right) }
%\renewcommand{\mho}{\mathcal{W}}
\newcommand{\pn}[1]{{{\partial #1} \over {\partial n}}}
\newcommand{\pnn}[1]{{{\partial #1} \over {\partial m}}}
\newcommand{\ps}[1]{{{\partial #1} \over {\partial s}}}
\newcommand{\p}{{\partial}}
\newcommand{\vn}{{\v \nabla}}
\newcommand{\lk}{\left(}
\newcommand{\rk}{\right)}

\newcommand{\tl}{\tilde}

\newcommand{\I}{\mathcal{I}}
\newcommand{\E}{\mathcal{E}}
\newcommand{\A}{\mathcal{A}}

\newcommand{\ep}{\epsilon}
\begin{document}

 \section*{The realistic non-hydrostatic model nh\_realistic}

nh\_realistic is pe\_realistic with a prognostic vertical velocity and a non-hydrostatic pressure.
This text describes only the differences to pe\_realistic.
Everything that is not mentioned here is as described in the documentation of pe\_realistic and pe\_simple.

The differences are the following.
The vertical velocity $w$ has its own momentum equation.
A three-dimensional pressure $p_{nh}$ makes the velocity free of divergence in every cell.
The horizontal coordinates must be Cartesian.

The package contains only what is new.
The grid, the masks, the tracers, the equation of state, the vertical mixing, the TKE closure, the isoneutral mixing, the surface pressure and the split-explicit external mode are imported from pe\_realistic.
The file \verb|momentum.py| holds the equation for $w$.
The file \verb|vmix.py| holds the vertical friction of $w$.
The files \verb|pressure.py| and \verb|solver.py| hold the equation for $p_{nh}$ and its solver.
The file \verb|model.py| holds the time step.

The formulation follows pyOM2 with \verb|enable_hydrostatic| off and the non-hydrostatic model of Oceananigans.
The solver for $p_{nh}$ follows Oceananigans.
The tests and their results are listed in the section on the nhr family of \verb|pyOM3/test/README_test_suite.md|.

\subsection*{Grid}

The grid is that of pe\_realistic with Cartesian coordinates.
The widths \verb|dxt|, \verb|dyt| and \verb|dzt| are in metres and may vary in all three directions.
The model stops with an error if \verb|coord_degree| is set.
On the sphere the equation for $w$ would need the metric term $-(u^2+v^2)/a$ and the Coriolis terms with $\cos \phi$.
nh\_realistic has neither of them.

The metric factors \verb|cost| and \verb|cosu| of pe\_realistic are therefore 1, and \verb|tantr| is 0.
The kernels imported from pe\_realistic still contain them.
The new kernels of nh\_realistic do not.
The horizontal area of a $T$ cell is $A^t = \Delta x^t \Delta y^t$, and its volume is $V^t = A^t \Delta z^t$.

The vertical velocity $w_k$ sits at the top face of the $T$ cell $k$.
This face is the centre of the $W$ cell $k$.
The $W$ cell $k$ reaches from the centre $z^t_k$ to the centre $z^t_{k+1}$.
Its height is $\Delta z^w_k$, and its volume is $A^t \Delta z^w_k$.

The mask \verb|maskW| of pe\_realistic is open at the surface face, since the free surface needs a vertical velocity there.
nh\_realistic adds the mask \verb|maskWc|, which is \verb|maskW| with the surface face closed.
The equation for $w$, its friction and the equation for $p_{nh}$ use \verb|maskWc|.
The surface face thus belongs to the free surface and not to the equation for $w$.

\subsection*{Pressure}

All pressures are divided by $\rho_0$, as in pe\_realistic.
The pressure is the sum
\beq
 p = p_h + p_s + p_{nh}
\eeq
of the hydrostatic pressure $p_h$, the surface pressure $p_s$ and the non-hydrostatic pressure $p_{nh}$.
$p_h$ and $p_s$ are computed as in pe\_realistic.

The momentum equations are
\beq \label{eq:mom}
 \pt{\v u} + \dots = - \vn_h \lk p_h + p_s + p_{nh} \rk ~,~~
 \pt w + \v u \cdot \vn w = - \pz{p_{nh}} + F_w
\eeq
where $\v u = (u,v)$ is the horizontal velocity, the dots stand for all other terms of pe\_realistic, and $F_w$ is the friction of $w$.

The equation for $w$ carries neither the buoyancy nor $p_h$ nor $p_s$.
The full equation would be $\pt w + \v u \cdot \vn w = - \pz p - g \rho/\rho_0 + F_w$.
$p_s$ does not depend on $z$.
$p_h$ is integrated downward from the surface with
\beq \label{eq:ph}
 p_h(k) = p_h(k+1) + \frac{g}{2 \rho_0} \lk \rho(k+1) + \rho(k) \rk \Delta z^w_k
\eeq
which is the integration of pe\_realistic.
\eq{eq:ph} gives at the $W$ point $k$ the vertical gradient
\beq
 \frac{p_h(k+1) - p_h(k)}{\Delta z^w_k} = - \frac{g}{2 \rho_0} \lk \rho(k+1) + \rho(k) \rk
\eeq
The right hand side is the buoyancy force averaged to the $W$ point.
The vertical gradient of $p_h$ and the buoyancy force therefore cancel exactly, also on a stretched grid.
Only $p_{nh}$ is left in the equation for $w$.

With $p_{nh} = 0$ and $w$ taken from continuity the model is pe\_realistic.
The attribute \verb|enable_nonhydrostatic| (default on) selects the non-hydrostatic model.
With it off the time step is that of pe\_realistic to the last bit.
This is tested.

\subsection*{Advection and friction of $w$}

The advection of $w$ is the energy-conserving second-order form of the momentum advection of pe\_realistic.
The transports through the faces of a $T$ cell are
\beq
 U_i = \Delta y^t \Delta z^t u_i ~,~~
 V_j = \Delta x^t \Delta z^t v_j ~,~~
 W_k = A^t w_k
\eeq
with the masks of the faces.
The transport through a face of the $W$ cell $k$ is the mean of the transports through the same face of the $T$ cells $k$ and $k+1$.
The fluxes of $w$ are
\beq
 f^x_i = - \frac{1}{4} \lk U_{i,k} + U_{i,k+1} \rk \lk w_i + w_{i+1} \rk ~,~~
 f^z_k = - \frac{1}{4} \lk W_k + W_{k+1} \rk \lk w_k + w_{k+1} \rk
\eeq
and $f^y$ correspondingly with $V$.
$f^z_k$ sits at the $T$ point $k+1$.
The advective tendency of $w_k$ is the divergence of these fluxes divided by the volume $A^t \Delta z^w_k$ of the $W$ cell.

The divergence of the transports through the $W$ cell $k$ is the mean of the divergences of the $T$ cells $k$ and $k+1$.
It vanishes when the velocity is free of divergence.
The advection then conserves $w^2$.
On a uniform grid it is the advection of $w$ of nh\_simple.

The equation for $w$ has no Coriolis term, as in pyOM2.

The horizontal friction of $w$ uses the coefficients \verb|Ahbi| and \verb|Ah| of the horizontal friction of $u$ and $v$.
It uses the operator at the $T$ points of pe\_realistic, since $w$ sits at the $T$ points in the horizontal.
For $w$ this operator is exact.
A flux exists only between two wet $W$ cells.

The vertical friction of $w$ has the flux at the $T$ point $k$ between the $W$ cells $k-1$ and $k$,
\beq
 F_k = A_k \frac{w_k - w_{k-1}}{\Delta z^t_k}
\eeq
The tendency of $w_k$ is $(F_{k+1} - F_k)/\Delta z^w_k$.
The roles of $\Delta z^t$ and $\Delta z^w$ are thus swapped with respect to the vertical friction of $u$ and $v$.
A flux exists only between two wet $W$ cells.

The coefficient $A_k$ follows the vertical friction of $u$ and $v$.
It is the constant \verb|A_v| with \verb|vmix_scheme| explicit or implicit.
With the TKE closure it is the viscosity of the closure averaged from the two $W$ points to the $T$ point.
The step is a forward step for \verb|vmix_scheme| explicit and an implicit step otherwise.
The implicit step is the tridiagonal solve of pe\_realistic on the $W$ grid.
The energy that the friction of $w$ removes is not given to the TKE closure.

\subsection*{Projection}

The velocities are first advanced with all terms except the gradient of $p_{nh}$.
This gives $u^*$, $v^*$ and $w^*$.
The surface pressure $p_s$ is already part of $u^*$ and $v^*$.
The vertical friction is also part of them.

The vertical velocity at the surface face is taken from continuity as in pe\_realistic,
\beq \label{eq:wsurf}
 w_{surf} = - \frac{1}{A^t} \sum_k \lk U_i - U_{i-1} + V_j - V_{j-1} \rk
\eeq
with $u^*$ and $v^*$ in the transports and the sum over the wet cells of the column.
With the free surface $w_{surf}$ is the change of the surface elevation in time.
With the rigid lid it is zero to the tolerance of the surface pressure solver.

The divergence of the $T$ cell is
\beq
 D = \frac{1}{V^t} \lk U_i - U_{i-1} + V_j - V_{j-1} + W_k - W_{k-1} \rk
\eeq
The velocity is corrected by
\beq \label{eq:proj}
 u_i = u^*_i - \Delta t \frac{p_{i+1} - p_i}{\Delta x^u_i} ~,~~
 v_j = v^*_j - \Delta t \frac{p_{j+1} - p_j}{\Delta y^u_j} ~,~~
 w_k = w^*_k - \Delta t \frac{p_{k+1} - p_k}{\Delta z^w_k}
\eeq
with $p = p_{nh}$, on the wet faces and the interior faces in $z$.
The surface face is closed in \eq{eq:proj}.
$w_{surf}$ therefore stays as it is.

The pressure $p_{nh}$ makes $D$ of the corrected velocity vanish in every wet cell.
Inserting \eq{eq:proj} into $D = 0$ gives
\beq \label{eq:pnh}
 a^x_i \lk p_{i+1} - p_i \rk - a^x_{i-1} \lk p_i - p_{i-1} \rk
 + a^y_j \lk p_{j+1} - p_j \rk - a^y_{j-1} \lk p_j - p_{j-1} \rk && \nonumber \\
 + a^z_k \lk p_{k+1} - p_k \rk - a^z_{k-1} \lk p_k - p_{k-1} \rk
 &=& \frac{V^t D^*}{\Delta t}
\eeq
with $D^*$ the divergence of $u^*$, $v^*$, $w^*$ and $w_{surf}$.
The coefficients are
\beq
 a^x_i = \frac{\Delta y^t \Delta z^t}{\Delta x^u_i} ~,~~
 a^y_j = \frac{\Delta x^t \Delta z^t}{\Delta y^u_j} ~,~~
 a^z_k = \frac{A^t}{\Delta z^w_k}
\eeq
on the wet faces and zero on the closed faces.
\eq{eq:pnh} is the volume times the divergence of the gradient of $p_{nh}$.

The coefficient of a face is the same in the equations of the two cells next to it.
The matrix of \eq{eq:pnh} is therefore symmetric, also on a stretched grid with topography.
This is tested.
Without the volume factor the matrix would not be symmetric on a stretched grid.

All faces at the boundary of the wet region are closed in \eq{eq:pnh}.
A solution exists only if the right hand side summed over the wet cells vanishes.
Summed over a column, the transports through the inner faces cancel.
What is left is the sum of the horizontal transports of the column and $A^t w_{surf}$.
With \eq{eq:wsurf} this sum is zero.
The right hand side thus sums to zero in every column, and \eq{eq:pnh} has a solution.

\eq{eq:pnh} fixes $p_{nh}$ only up to a constant.
After every iteration the mean of $p_{nh}$ over the wet cells is set to zero, as in Oceananigans.
The same is done with the residual.

\subsection*{Solver}

\eq{eq:pnh} is solved with the method of conjugate gradients with a preconditioner.
The solver is written for JAX and MPI in the same way as the solver for $p_s$ of pe\_simple.
The first guess is $p_{nh}$ of the previous time step.

The preconditioner solves the vertical part of \eq{eq:pnh} exactly in every column.
It drops the horizontal coefficients $a^x$ and $a^y$.
The vertical part is the tridiagonal system
\beq
 - a^z_{k-1} z_{k-1} + (1 + \epsilon) \lk a^z_k + a^z_{k-1} \rk z_k - a^z_k z_{k+1} = r_k
\eeq
for the preconditioned residual $z$ with the residual $r$.
It is solved with the Thomas algorithm.
The factor $1+\epsilon$ with $\epsilon = 0.01$ makes the system of each column regular, as in Oceananigans.
In a dry cell and in a wet cell without a wet cell above or below, the preconditioner is the identity.
The preconditioner is that of Marshall et al.\ (1997).

The preconditioner suits grids with $\Delta z \ll \Delta x$.
There the vertical coefficients $a^z$ are much larger than the horizontal ones.
In the test configuration they are larger by a factor of $3 \cdot 10^4$ to $2 \cdot 10^7$.
The preconditioner reduces the number of iterations there from 1000 to 3500 to 50 to 150 per time step.
For $\Delta z \approx \Delta x$ it helps less.

The residual $r$ is the right hand side minus the left hand side of \eq{eq:pnh}.
The divergence of the velocity corrected with the current $p_{nh}$ is $\Delta t\, r/V^t$.
The iteration stops when $\Delta t \max |r|/V^t$ is smaller than \verb|epsilon_nh| (default $10^{-12}\,$s$^{-1}$).
\verb|epsilon_nh| is thus the largest divergence that the velocity keeps.
The iteration also stops after \verb|max_iterations_nh| iterations (default \verb|max_iterations|).

The stopping condition is met by the vertical structure of $p_{nh}$.
The horizontal structure of $p_{nh}$ is less well determined when $\Delta z \ll \Delta x$.
It still changes with the round-off of the sums in the solver.
Runs with different domain decompositions therefore differ by about $10^{-8}$ relative with the default \verb|epsilon_nh|.
With \verb|epsilon_nh| $= 10^{-16}$ they differ by about $10^{-10}$.

\subsection*{Tracers}

The tracers are advected with $u$, $v$ and $w$ of the end of the previous time step, as in pe\_realistic.
In pe\_realistic $w$ follows from continuity and the velocity is free of divergence to round-off.
In nh\_realistic the velocity is free of divergence to \verb|epsilon_nh|.
A uniform tracer therefore stays uniform only to \verb|epsilon_nh| $\Delta t$ per time step.
The content of a tracer is still conserved to round-off, since the advection is in flux form.
The WENO scheme of pe\_realistic works in nh\_realistic unchanged.
For a uniform tracer it gives the same tendency as the mean of the two cells, the tracer times the divergence that the solve leaves.

\subsection*{Time step}

The order of the calls in a time step of the split-explicit scheme is
\begin{enumerate}
\item diffusivities of the TKE closure (\verb|set_tke_diffusivities|)
\item momentum tendencies, Coriolis force and hydrostatic pressure (\verb|momentum_tendency_expl|)
\item advective tendency of $w$ (\verb|w_tendency|)
\item forcing of the configuration (\verb|set_forcing|)
\item horizontal friction of $w$ and Adams--Bashforth update of $w$
\item bottom friction
\item external mode with sub cycling (\verb|solve_expl_surface|)
\item vertical friction of $u$, $v$ and of $w$
\item surface vertical velocity and the solve for $p_{nh}$ (\verb|nh_projection|)
\item tracer advection and biharmonic tracer mixing, or the advection outside the Adams--Bashforth step
\item Adams--Bashforth update of the tracers
\item isoneutral mixing
\item vertical mixing of the tracers
\item density
\item TKE step (\verb|integrate_tke|)
\end{enumerate}

The order of the calls in a time step of the rigid lid and of the implicit free surface is
\begin{enumerate}
\item diffusivities of the TKE closure
\item tracer advection and biharmonic tracer mixing
\item momentum tendencies and hydrostatic pressure (\verb|momentum_tendency_impl|)
\item advective tendency of $w$ (\verb|w_tendency|)
\item forcing of the configuration (\verb|set_forcing|)
\item biharmonic and harmonic friction of $u$, $v$ and of $w$
\item Adams--Bashforth update of $u$, $v$ and of $w$
\item direct forcing of the configuration (\verb|set_forcing_direct|)
\item surface pressure and its gradient (\verb|solve_for_pressure|)
\item vertical friction of $u$, $v$ and of $w$
\item bottom friction
\item surface vertical velocity and the solve for $p_{nh}$ (\verb|nh_projection|)
\item Adams--Bashforth update of the tracers, or the advection outside the Adams--Bashforth step
\item isoneutral mixing
\item vertical mixing of the tracers
\item density
\item TKE step
\end{enumerate}
In both schemes the steps are those of pe\_realistic.
The steps for $w$ are added.
The solve for $p_{nh}$ replaces the vertical velocity from continuity.

The forcing of the configuration may add to the tendency \verb|dw| at the level \verb|tau|, as it does for \verb|du| and \verb|dv|.
The tendency then enters the Adams--Bashforth step.

The solve for $p_{nh}$ comes after the vertical friction and the bottom friction.
The velocity at the end of the time step is therefore free of divergence.
The surface pressure is solved before.
It removes the depth-averaged divergence, which is the mode with the slowest convergence in \eq{eq:pnh}.

The restart needs the fields of pe\_realistic and in addition \verb|dw| and \verb|p_nh|.
Without \verb|p_nh| the solver starts from another first guess.
The run is then no longer identical to the last bit.

\section*{Model parameters and options}

The parameters and options of pe\_realistic are also those of nh\_realistic.
The tables list only the differences.

\subsection*{Required parameters}

\begin{longtable}{@{}llp{8.4cm}@{}}
\hline
Attribute & Unit & Meaning \\
\hline
\endhead
\verb|dxt|, \verb|dyt| & m & widths of the cells in $x$ and $y$, as in pe\_realistic but in metres only \\
\hline
\end{longtable}

\noindent
\verb|coord_degree| must be off.

\subsection*{Options with defaults}

\begin{longtable}{@{}llp{7.6cm}@{}}
\hline
Attribute & Default & Meaning \\
\hline
\endhead
\verb|enable_nonhydrostatic| & on & prognostic $w$ and the solve for $p_{nh}$. Off gives pe\_realistic \\
\verb|epsilon_nh| & $10^{-12}$ & largest divergence (s$^{-1}$) that the solve for $p_{nh}$ leaves \\
\verb|max_iterations_nh| & \verb|max_iterations| & largest number of iterations of the solve for $p_{nh}$ \\
\verb|Ah|, \verb|Ahbi| & unset & as in pe\_realistic, also acting on $w$ \\
\verb|A_v| & $0$ & as in pe\_realistic, also acting on $w$ \\
\verb|weno_advection|, \verb|weno_ref| & off, $1$ & WENO tracer advection as in pe\_realistic.
  It needs \verb|halo_points| $=3$ \\
\hline
\end{longtable}

\noindent
The number of iterations of the last solve for $p_{nh}$ is \verb|congr_itts_nh|.
The divergence it leaves is \verb|estimated_error_nh|.

\subsection*{Snapshot diagnostics}

The snapshot classes \verb|snap_nc4| and \verb|snap_raw| are those of pe\_realistic.
\verb|snap_nc4| also writes $p_{nh}$ under the name \verb|p_nh|.
$w$ is written on the $W$ points as in pe\_realistic.
It is now the prognostic $w$.

\subsection*{Limits}

\begin{itemize}
\item The horizontal coordinates must be Cartesian, see the grid above.
\item The equation for $w$ has no Coriolis term.
\item The energy that the friction of $w$ removes is not given to the TKE closure.
\item With the free surface $w_{surf}$ does not enter the advection of $w$, since the surface face is closed in the equation for $w$. The advection of $w$ then does not conserve $w^2$ exactly in the uppermost $W$ cell.
\item The preconditioner suits $\Delta z \ll \Delta x$. For $\Delta z \approx \Delta x$ the number of iterations grows. A preconditioner with the FFT of a uniform grid, as in Oceananigans, is not available.
\item Runs with different domain decompositions are not identical to the last bit, see the solver above.
\item The LES closures of nh\_simple are not part of nh\_realistic.
\end{itemize}

\end{document}

\clearpage\phantomsection\label{doc:wim_simple}\addcontentsline{toc}{section}{wim\_simple}
\documentclass[12pt]{article}
\usepackage{lineno}
%\usepackage{ametsoc}
\usepackage{natbib}
\usepackage{a4wide}
\usepackage{longtable}
\usepackage{amsfonts}
\usepackage{amssymb}

\usepackage{color}
\usepackage[pdftex]{graphicx}
\usepackage[pdftex,colorlinks, citecolor=black,linkcolor=black,urlcolor=black]{hyperref}
\definecolor{darkgreen}{rgb}{0.1,0.4,0.1}
\definecolor{darkblue}{rgb}{0.1,0.1,0.3}



\newif\ifdetail
\global\detailfalse
%\global\detailtrue

\def\bem#1{\colorbox{green}{\em #1}}

\newcommand{\beq}  { \begin{eqnarray} }
\newcommand{\eeq}  { \end{eqnarray}}
\newcommand{\beeq}  { \begin{eqnarray*} }
\newcommand{\eeeq}  { \end{eqnarray*}}
\newcommand{\eq }  [1]{Eq.~(\ref{#1})}  % Eq. (#1)
\newcommand{\eqs }  [1]{Eqs.~(\ref{#1})}
\newcommand{\sect}  [1]{Sec.~(\ref{#1})}  % Eq. (#1)
\newcommand{\fig}  [1]{Fig.~\ref{#1}}   % Fig. #1
\renewcommand{\v}[1]{{\mbox{\boldmath$ #1 $}}}  % Vektor durch Fettdruck
\newcommand{\m}    [1]{{\bf  #1 }}              % Matrix in sans serif
\newcommand{\gz}[1] {{{d #1} \over {dz}}}
\newcommand{\pt}[1]{{\partial_t #1}}
\newcommand{\ptt}[1]{{\partial_{tt} #1}}
\newcommand{\pttt}[1]{{\partial_{ttt} #1}}
\newcommand{\pttx}[1]{{{\partial^3 #1} \over {\partial t^2 \partial x}}}
\newcommand{\ptty}[1]{{{\partial^3 #1} \over {\partial t^2 \partial y}}}
\newcommand{\ptx}[1]{{{\partial^2 #1} \over {\partial t \partial x}}}
\newcommand{\pty}[1]{{{\partial^2 #1} \over {\partial t \partial y}}}
\newcommand{\pz}[1]{{\partial_z #1}}
\newcommand{\pzz}[1]{{\partial_{zz} #1}}
\newcommand{\px}[1]{{\partial_x #1}}
\newcommand{\py}[1]{{\partial_y #1}}
\newcommand{\pxi}[1]{{\partial_i #1}}
\newcommand{\pxj}[1]{{\partial_j #1}}
\newcommand{\pxk}[1]{{\partial_k #1}}
\newcommand{\pb}[1]{{{\partial #1} \over {\partial b}}}
%\newcommand{\pxy}[1]{{{\partial^2 #1} \over {\partial x \partial y}}}
\newcommand{\pxy}[1]{\partial_{xy} #1}
\newcommand{\pyy}[1]{\partial_{yy} #1}
\newcommand{\pxx}[1]{\partial_{xx} #1}
\newcommand{\kx}{{\bf k} \times}
\newcommand{\dt}[1]{{{d #1} \over {d t}}}
\newcommand{\Dt}[1]{{{D #1} \over {D t}}}
\newcommand{\Div}[1]{ \nabla \cdot #1 }
\renewcommand{\div}[1]{ \nabla_{yz} \cdot #1 }
\newcommand{\Bar}[1]{ \overline{ #1} }
\newcommand{\rvec} [1]{
    \raisebox{-1.5ex}{$\stackrel{\textstyle #1}{\neg}$} }% rotierter Vektor
\newcommand{\nablq}   {\rvec\nabla}                    % rotated nabla operator
\newcommand{\Vec}[1]{\left(\begin{array}{c} #1 \end{array}\right) }
%\renewcommand{\mho}{\mathcal{W}}
\newcommand{\pn}[1]{{{\partial #1} \over {\partial n}}}
\newcommand{\pnn}[1]{{{\partial #1} \over {\partial m}}}
\newcommand{\ps}[1]{{{\partial #1} \over {\partial s}}}
\newcommand{\p}{{\partial}}
\newcommand{\vn}{{\v \nabla}}
\newcommand{\lk}{\left(}
\newcommand{\rk}{\right)}

\newcommand{\tl}{\tilde}

\newcommand{\I}{\mathcal{I}}
\newcommand{\E}{\mathcal{E}}
\newcommand{\A}{\mathcal{A}}

\newcommand{\ep}{\epsilon}
\newcommand{\erfc}{\mathrm{erfc}}
 
 
\begin{document}

 \section*{The spectral wave model wim\_simple (Wave Interaction Model, WIM)}
 
wim\_simple is a spectral surface wave model, coupled to an ocean model in an energetically
and dynamically consistent way.
Basis is the wave action spectral density equation for deep water surface gravity waves
\beq \label{eqA}
 \pt{A}  +  \dot{\v x} \cdot  \vn_{\v x}  A   +\dot{\v k} \cdot  \vn_{\v k} A  &=& S = S_{in} + S_{out} + S_{nl}
\eeq
with wave action density $A(\v k,\v x,t)$, the horizontal wavenumber vector $\v k = (k_x,k_y)$,
the dispersion relation $\Omega=\sigma + \v U \cdot \v k$ with intrinsic frequency
$\sigma(\v k) = \sqrt{ g |\v k| }$ and Doppler shift velocity $\v U$,
transport velocities $  \dot{\v x} =   \vn_{\v k} \Omega$ and $\dot{\v k}= - \vn_{\v x} \Omega$,
action input by wave growth $S_{in}$, action dissipation by wave breaking $S_{out}$,
 and non-linear action transfer $S_{nl}$
and operators $\vn_{\v x} = (\p_x,\p_y)$ and $\vn_{\v k} = (\p_{k_x}, \p_{k_y})$.
Only deep water waves are considered, with $\vn_{\v x} \sigma = 0$ and $\p_t \sigma = 0$,
but the concept extends readily to waves at finite depth.
Since the scattering integral $S_{nl}$ conserves energy and momentum of the waves,
it holds that $\int \sigma S_{nl} d \v k = 0$ and $\int \v k S_{nl} d \v k = 0$.


The surface wave model predicts the Stokes velocity $\v u_S$ and the wave stress $\v \tau_o$
to force an ocean model.
Considering the Lagrangian ocean velocity $\v u_L = \v u + \v u_S$ and its kinetic energy,
 the consistent energy cycle becomes 
\beq \label{eq_K}
\p_t \int_{-\infty}^0  |\v u_L|^2/2 dz + \vn_{\v x}  \cdot \int_{-\infty}^0  \v u^L_h  \pi  dz  &=&   C  \\  \label{eq_W}
\p_t  \int E d \v k +  \vn_{\v x}  \cdot  \int \lk  \dot{\v x}   E  + \dot{ \v x}  A (\v U \cdot \v k)  \rk d \v k
   &=& \int (\Omega S_{in} +\v U \cdot \v k \, S_{nl}) d \v k +\int \sigma S_{out}   d \v k   -C  \nonumber
\eeq
with the wave energy spectral density $E = \sigma A$ and the exchange term 
\beq \label{eq_C}
   C&=&  \int_{-\infty}^0  \v u^L_h \cdot   \p_t \v u_S  dz +\int_{-\infty}^0  \v u^L_h  \cdot    \p_z \v \tau  dz 
   = \int  \v U \cdot \v k (\p_t A -S_{out}) d \v k
   \eeq
   and the setting, following Vanneste and Young for the Doppler shift velocity $\v U$,
\beq \label{eq_Q}
  \v U(\v x,\v k,t) = \int_{-\infty}^0 \v u^L_h Q  dz  ~~,~~\v u^S_h(\v x,z,t)  = \int Q A \v k d \v k~~,~~Q(k,z) = 2 k e^{2 kz}    
  \\
\p_z \v \tau(\v x,z,t) = -  \int  Q    \v k S_{out} d \v k~,~ \v \tau_o(\v x,t) =\int_{-\infty}^0 \p_z \v \tau dz =- \int   \v k S_{out} d \v k
\eeq
with $\int_{-\infty}^0 Q \, dz = 1$ and the ocean momentum forcing $ \v \tau_o =-\int   \v k S_{out} d \v k$,
i.e.\ the stress on the ocean is the dissipation of wave momentum.
Note that $\v U$ contains both the Eulerian mean and the Stokes horizontal velocity components.
Waves and ocean exchange energy by the effect of the Doppler shift and the ocean stress, both
given by $C$; other source terms of the wave energy are the wave growth by the wind,
$\int \Omega S_{in} d \v k$, and the dissipation by wave breaking, $\int \sigma S_{out} d \v k$,
which may be directed into a turbulent ocean energy reservoir (see below).
The cycle closes exactly up to the term $\int \v U \cdot \v k \, S_{nl} d \v k$, which represents
the work done on the sheared mean flow by the vertical redistribution of Stokes momentum
through the nonlinear transfer. Only the baroclinic part of $\v U$ contributes to this term,
since $\int \v k S_{nl} d \v k = 0$. The discrete interaction approximation used for $S_{nl}$
evaluates resonant quadruplets based on the intrinsic frequency and thus enforces
$\int \sigma S_{nl} d \v k = 0$, such that with a background current this term appears as a
spurious source; since it is small compared to the other sources, it is ignored.


\subsection*{Relation to WAM}

The concept and discretisation of wim\_simple follows closely the example of WAM
found in WAM (cycle 7) at \url{https://github.com/myWAveModel/WAM},
except for simplified action sources.
Instead of the canonical wavenumber coordinates $\v k$ used above,
frequency $\nu = \sigma/(2 \pi)$ and wavenumber angle $\theta$
with $k_x = k \sin \theta$, $k_y = k \cos \theta$ and $k=|\v k|$ are used as phase dimensions,
covered by an equidistant grid in $\theta$ and a logarithmic grid in $\nu$ as in WAM.
Different to WAM, however, $\nu$ is here the \emph{intrinsic} frequency, while in WAM
the extrinsic frequency is used; the intrinsic frequency is better suited to
represent wave--mean flow interaction, following the example of WAVEWATCH III.
Instead of the spectral energy density $E(\v k) = \sigma A(\v k)$, the model carries
the WAM-like spectral density $E(\theta,\nu)$, related by
$ \int_{-\infty}^0  |\v u_L|^2/2 dz= \int E(\v k) d \v k = g \int E(\theta,\nu) d \theta d\nu  $.
Units of $E(\v k) d \v k$ are $m^3/s^2$,
thus $E(\theta,\nu)  d \theta d\nu$ has units $m^2$ (the surface elevation variance)
and $E(\theta,\nu)$ has units $m^2 s $,
while $E(\v k)$ has units $m^5/s^2$; the factor $g$ shows up accordingly in the
definitions of $\v u_S$ and $\v \tau$ below.
\ifdetail
Coordinate transformation is not necessary but given here for completeness.
 From $(k_x,k_y)$ to ($\theta,k$) 
with $k_x=k \cos \theta$, $k_y=k \sin \theta$
\beq
E(\v k) d \v k  = k E(\v k) d \theta d k
\eeq
since (take absolute value)
\beq
\left |   \begin{array}{cc}
\p k_x /\p \theta & \p k_x/ \p k  \\
\p k_y /\p \theta & \p k_y/ \p k  \\
\end{array} \right|  = 
\left |   \begin{array}{cc}
-k \sin \theta & \cos \theta  \\
k \cos \theta &  \sin \theta\\
\end{array} \right| 
= - k \sin^2 \theta - k \cos^2 \theta = -k
\eeq
 Transformation from  $(\theta,k)$ to  $(\theta,\sigma)$ with $k = \sigma^2/g$ by
\beq
   d \theta d k =  d \theta \frac{\p k}{\p \sigma} d \sigma = 2 (\sigma/g)  d \theta  d \sigma
   ~~,~~
   d \v k =   2 (\sigma^3/g^2)  d \theta  d \sigma
\eeq 
\fi




The model equation formulated for $A(\theta,\nu)$ reads
\beq \label{eq6}
% \pt{ A}    +  \dot{\v x} \cdot  \vn_{\v x}  A  +\dot{\theta} \p_{\theta} A   +\dot{\sigma} \p_{\sigma} A
% &=& S_{in} + S_{out} + S_{nl} 
% \\
  \pt{ A}    +  \vn_{\v x}  \cdot   (\dot{\v x} A ) +\p_{\theta}  (\dot{\theta} A  )
 + \p_{\nu} ( \dot{\nu} A)
 &=&S_{in} + S_{out} + S_{nl} 
\eeq
We need expressions for group velocity $\dot{\v x}$ and the drift velocities $\dot{\theta}$
and $\dot \nu$.

\subsubsection*{Group velocity}

Since $\v U = \v U(k)$, the group velocity is given by
\ifdetail
\beq
 \dot x &=& \p_{k_x} \Omega = \p_{k_x} \sigma  + \p_{k_x}  Uk_x  + \p_{k_x}  V k_y  
 =   \sqrt{g} \p_{k_x}  (k_x^2 + k_y^2)^{1/4}+ U + k_x \p_{k_x}  U + k_y\p_{k_x}  V 
 \\ &=&  (c /2)  k_x/k + U + \v k \cdot \p_{k_x}  \v U 
 \\
 \dot y &=& \p_{k_y} \Omega = (c /2)  k_y/k+ \p_{k_y}  Uk_x  + \p_{k_y}  V k_y  
 =  (c /2)  k_y/k+ V + \v k \cdot \p_{k_y}  \v U   
 \eeq
 \fi
 \beq
 \dot{\v x} =  (c /2)  \v s + \v U + \v k \cdot \vn_{\v k}  \v U 
 \eeq
with the phase velocity  $c = \sigma/k = \sqrt{g} /\sqrt{k} = g /\sigma $
and the wave direction vector $\v s = (\sin \theta,\cos \theta)$.
\ifdetail
Note that
\beq 
\vn_{\v k}  ( \v  k \cdot    \v U)   = \Vec{  \p_{k_x} (U k_x + V k_y) \\ \p_{k_y} (U k_x + V k_y) }
= \v U + \Vec{   k_x  \p_{k_x} U  + k_y   \p_{k_x}  V \\  k_x \p_{k_y} U  +k_y  \p_{k_y} V   }
\\
= 
 \v U + \Vec{   (k_x  \p_{k} U  + k_y    \p_{k}  V) k_x/k  \\  
 (k_x \p_{k} U  +k_y  \p_{k} V) k_y/k    }
\eeq
with $ \p_{k_x} k  = % \p_{k_x} (k_x^2+k_y^2)^{1/2} = 1/2  (k_x^2+k_y^2)^{-1/2}  2 k_x = 
k_x/k$
and  $\v U = \v U(k)$
\fi
Since the wavenumber gradient of the Doppler velocity $\v U$
is purely radial (it points along $\v k$),
\beq \label{eq:kgradkU}
\v k \cdot \vn_{\v k} \v U
= \Vec{ (k_x \p_k U + k_y \p_k V)\, k_x/k \\
        (k_x \p_k U + k_y \p_k V)\, k_y/k }
= \lk \v s \cdot \p_k \v U \rk \, \v k
= k \lk \sin\theta\, \p_k U + \cos\theta\, \p_k V \rk \, \v s ,
\eeq
with  $\p_k \v U = (\p_k U, \p_k V)$. The full group velocity reads
\beq \label{eq:xdot_full}
\dot{\v x} = (c/2)\,\v s + \v U + k \lk \sin\theta\, \p_k U + \cos\theta\, \p_k V \rk \, \v s .
\eeq
With $k = (2\pi \nu)^2/g$ one has
$k\,\p_k U_i = k \p_\nu U_i \p \nu/\p k = \frac{\nu}{2}\,\p_\nu U_i$
%with $2 \pi \nu =  \sqrt{g k} $ and $  (2 \pi) \p \nu/\p k = g/2 /\sqrt{g k}  $ 
%and $k   (2 \pi) \p \nu/\p k = \sqrt{g k} /2 =\pi f $
%
 so the radial correction can be written with
a single scalar factor
\beq \label{eq:xdot_disc}
R = \frac{\nu}{2}\lk \v s \cdot \p_\nu \v U  \rk ,
\qquad
\dot{\v x} %= \Vec{ (c/2)\sin\theta + U + R\,\sin\theta  \\ (c/2)\cos\theta + V + R\,\cos\theta  }
= (c/2) \v s  + \v U + R \v s
\eeq
i.e.\ the same directional projection $R\,\v s$ enters both components. This reduces
to the pure group velocity $\dot{\v x} = (c/2)\v s + \v U$ whenever the current has no
frequency dependence ($\p_\nu \v U = 0$), e.g.\ for a barotropic current.

\subsubsection*{Other drift velocities}

Since $\v U = \v U(k)$, the material derivative of $\v U$ along a ray contains,
besides $\p_t \v U + \dot{\v x} \cdot \vn_{\v x} \v U$,  also the contribution
$\dot{\v k} \cdot \vn_{\v k} \v U$, so  $\dot \sigma$
becomes
with %$k_x = k \sin \theta$ and $k_y = k \cos \theta$ and 
$\dot{\v k} =  - (\vn \v U )\cdot \v k $
\ifdetail
\beq
\dot \sigma  &=& 2 \pi \dot \nu =   \p_t \Omega - \dot{\v U}  \cdot \v k  - \v U \cdot \dot{\v k}
\\
&=&  \p_t \Omega -  \dot U k_x + \dot V k_y  - \v U \cdot \dot{\v k}
\\ &=&  \p_t \Omega -  (\p_t U + \v u \cdot \vn_x U + \dot{\v k} \cdot \vn_{\v k} U) k_x + \dot V k_y  - \v U \cdot \dot{\v k}
\\ &=& (\p_t \v U ) \cdot \v k 
-  (\p_t \v U ) \cdot \v k -  ( \dot{\v x} \cdot \vn_x \v U) \cdot \v k 
-(\dot{\v k} \cdot \vn_{\v k} \v U) \cdot \v k
-  \v U \cdot \dot{\v k} 
\\ &=& -  ( \dot{\v x} \cdot \vn_x \v U) \cdot \v k 
-(\dot{\v k} \cdot \vn_{\v k} \v U) \cdot \v k
+  \v U \cdot  ( (\vn \v U )\cdot \v k )
\\
&=&  -   \dot{\v x} \cdot ((\vn_x \v U)  \cdot \v k) +  \v U \cdot  ( (\vn \v U )\cdot \v k )
-(\dot{\v k} \cdot \vn_{\v k} \v U) \cdot \v k
\\ &=& -( \v U  -   \dot{\v x} ) \cdot \dot \v k-(\dot{\v k} \cdot \vn_{\v k} \v U) \cdot \v k
\eeq
\fi
\beq \label{eq:omdot_Uk}
\dot \sigma &=& \p_t \Omega - \dot{\v U} \cdot \v k - \v U \cdot \dot{\v k}
 =  ( \dot{\v x}   - \v U ) \cdot \dot{\v k}
   - ( \dot{\v k} \cdot \vn_{\v k} \v U ) \cdot \v k
\eeq
\ifdetail
\beq
 ( \dot{\v k} \cdot \vn_{\v k} \v U ) \cdot \v k 
 &=&  (\dot{\v k} \cdot \vn_{\v k} U) k_x +  (\dot{\v k} \cdot \vn_{\v k}) V k_y
 \\ &=&  (\dot{k_x} \p_{k_x} U +(\dot{k_y} \p_{k_y} U ) k_x +  (\dot{\v k} \cdot \vn_{\v k}) V k_y
 \\ &=&  (\dot{k_x} \p_{k} U k_x/k +(\dot{k_y} \p_{k} U k_y/k) k_x +  (\dot{\v k} \cdot \vn_{\v k}) V k_y
 \\ &=& \p_{k} U  (\dot{\v k} \cdot \v k) k_x /k  + \p_{k} V  (\dot{\v k} \cdot \v k )  k_y/k
 =  k (\dot{\v k} \cdot \v s)  ( \p_k \v U \cdot \v s)
\eeq  
with $ \p_{k_x} k  = k_x/k$.
\fi
With $\vn_{\v k} U_i = (\p_k U_i) \, \v k /k$ and $k \, \p_k U_i = (\nu/2) \p_\nu U_i$
one finds
$( \dot{\v k} \cdot \vn_{\v k} \v U ) \cdot \v k
 = ( \dot{\v k} \cdot \v s ) \, k \lk \v s \cdot \p_k \v U \rk
 %=  ( \dot{\v k} \cdot \v s ) \,  \lk \v s \cdot (\nu/2) \p_\nu \v U \rk
 = R \, ( \v s \cdot \dot{\v k} )$
with $R$ from \eq{eq:xdot_disc}. Since also
$\dot{\v x} - \v U = (c/2 + R) \, \v s$ by \eq{eq:xdot_disc}, we obtain
\beq \label{eq:omdot_final}
\dot \sigma = 2 \pi \dot \nu = 
 (  (c/2 + R) \, \v s ) \cdot \dot{\v k}
   -  R \, ( \v s \cdot \dot{\v k} ) 
=
(c/2) \, \v s \cdot \dot{\v k}
\eeq
Finally,
\ifdetail
\beq
\tan \theta &=& k_x/k_y ~~\to ~~ \dot \theta /\cos^2\theta = (1/k_y) \dot k_x - (k_x/k_y^2) \dot k_y 
\eeq
\fi
\beq
\dot \theta = ( \cos\theta   \dot k_x -    \sin \theta \dot k_y ) k^{-1}
\eeq
results
using $\tan \theta = k_x/k_y$ and 
$\dot k_x = -  k_x \p_x U - k_y \p_x V $ and  $\dot k_y = - k_x \p_y  U   -k_y \p_y  V $.

\ifdetail
\subsubsection*{Energy transfer to mean flow}

Multiplying \eq{eq6} by $\sigma$ gives the budget of the intrinsic energy
$E = \sigma A$. Since $\sigma = 2 \pi \nu$ is the spectral coordinate, it
commutes with $\p_t$, $\vn_{\v x}$ and $\p_\theta$; only the frequency
advection generates
$\sigma \, \p_\nu ( \dot \nu A) = \p_\nu ( \dot \nu E) - \dot \sigma A$
so that
\beq \label{eq:ebudget}
\pt{E} + \vn_{\v x} \cdot ( \dot{\v x} E ) + \p_\theta ( \dot \theta E )
 + \p_\nu ( \dot \nu E )
 &=& \sigma \dot A + \dot \sigma A
\eeq
All terms on the left are conservative transport of intrinsic energy; on the
right, $\sigma \dot A$ is the energy of the source terms and $\dot \sigma A$
is the local energy exchange with the mean flow mediated by propagation.
Multiplying \eq{eq6} instead by $\Omega(\v x, \theta, f, t)$
yields 
\ifdetail
\beq 
  \Omega  \pt{ A}    +  \Omega  \vn_{\v x}  \cdot   (\dot{\v x} A ) 
  + \Omega \p_{\theta}  (\dot{\theta} A  )
 + \Omega \p_{\nu} ( \dot{\nu} A)
 &=& \Omega  \dot A 
 \\
  \pt{ \Omega   A}  - A \p_t \Omega 
  +  \vn_{\v x}  \cdot ( \Omega  \dot{\v x} A )  
  + \p_{\theta}  (\Omega  \dot{\theta} A  )
  +\p_{\nu} (  \Omega  \dot{\nu} A) 
  - A ( \dot{\v x}    \cdot \vn_x  \Omega 
   + \dot{\theta} \p_{\theta} \Omega  
  +  \dot{\nu}  \p_\nu  \Omega )
   &=& \Omega  \dot A 
   \eeq
   \beq
  \pt{ \Omega   A} 
  +  \vn_{\v x}  \cdot ( \Omega  \dot{\v x} A )  
  + \p_{\theta}  (\Omega  \dot{\theta} A  )
  +\p_{\nu} (  \Omega  \dot{\nu} A) 
    &=& \Omega  \dot A  + A \p_t \Omega 
\eeq
\fi
\beq \label{eq:absE}
 \p_t   E + \vn_{\v x} \cdot ( \dot{\v x} \Omega A )
 + \p_\theta ( \dot \theta \Omega A ) + \p_\nu ( \dot \nu \Omega A )
 &=& \Omega \dot A  - \v U \cdot \v k \p_t A
\eeq
using $\dot{\v x} \cdot \vn_{\v x} \Omega + \dot \theta \, \p_\theta \Omega
 + \dot \nu \, \p_\nu \Omega = 0$.

Except for flux divergencies the
exchange $C$ thus proceeds through three routes: (i) the refraction transfer
$ \dot \sigma A$ to the ocean,
carried by the frequency advection,
(ii) the source-mediated part
$\v U \cdot \v k \, ( S_{in} + S_{nl} ) $,
handed over as $\p_t \v u_S$-work, and (iii) the dissipation part routed
through the stress $\v \tau$ related to $\v U \cdot \v k \,  S_{out}$.
In a horizontally homogeneous column all propagation terms vanish, $\p_t A = \dot
A$, and $C$ collapses to routes (ii) and (iii): energy flows to the ocean
with $\dot \sigma \equiv 0$, since each action quantum created or destroyed
by the sources carries the mean-flow energy $\v U \cdot \v k$ alongside its
intrinsic energy $\sigma$.


\fi
\subsection*{High frequency cut-off}

Since the growth rates $\gamma_{in}$ and $\gamma_{out}$ increase strongly with frequency,
which would require a very small time step at large $\nu$,
the wave model is only calculated up to  a cut-off frequency $\nu_c$,
and for the higher frequencies $\nu > \nu_c$ 
a Phillips spectrum with a $\nu^{-5}$ tail of the wave
spectrum is  applied, the so-called high-frequency tail. It has the form 
$E_{hi}(\nu)  =   E(\nu_c,\theta) \nu_c^5 \nu^{-5}$. 
The cut-off frequency $\nu_c$ is calculated as in WAMDI (1988) 
\beq
\nu_c = \max( 2.5 \bar \nu, 3 \nu_{pm})
\eeq
with the mean frequency 
$2 \pi \bar \nu =  \int E d \theta d\nu/ \int (E/\sigma) d \theta d\nu$, the ''wind sea'' frequency $2 \pi \nu_{pm} =  g/(28 u_{\star})$,
and the friction velocity $u_\star$.
For $\nu > \nu_c$ it is assumed that \eq{eq6} becomes local and steady, such that the
source terms balance locally in space and time and in $\nu$ and $\theta$, i.e.\ $S |_{\nu>\nu_c}=0$.
The per-column prognostic cut-off and the imposed $\nu^{-5}$ tail are activated by the
attribute \verb|mask_energy| (default off); the resulting prognostic/diagnostic bin mask
is also the basis of the equilibrium tail-flux booking of \verb|tail_flux| $=$ \verb|'input'|
described below, which therefore requires \verb|mask_energy| to be active.


\subsection*{Source terms}

The source terms $S_{in}$ and $S_{out}$ are parameterized by growth rates,
$ S_{in} = \gamma_{in} A$ and $S_{out}=  \gamma_{out} A$, 
for which the simplest forms  following  WAMDI (1988) are used
\beq \label{eq:gammas}
 \gamma_{out} = -  \beta_{o} g^{-4}  \bar E^2 \bar \sigma^7 \, \,     \sigma^2  ~&,&~
 \gamma_{in} =   \sigma\frac{\rho_a}{\rho_o} \times \, 
    0.25 \max\left( 0,
     28 (u^*/c)  \cos(\theta - \theta*) -1  \right) )
 \eeq
with mean energy $\bar E = \int E d \theta d\nu$,  mean frequency
$\bar \sigma = \bar E/ \int (E/\sigma) d \theta d\nu$,
air,water density  $\rho_a,\rho_o$, atmospheric friction velocity magnitude $u^{\star}$,
with direction $ \theta^*$, deep water phase velocity $c = g/  \sigma$,
and the parameter $\beta_{o} = 2.56 $.
Other parameterisations are implemented as options.
The mean frequency is clipped to the resolved spectral range,
\beq
\bar\sigma \leftarrow \min\!\left(\max(\bar\sigma,\, \sigma_0),\, \sigma_{N_\nu-1}\right)
\eeq
The non-linear source term $S_{nl}$ is the discrete interaction approximation (DIA)
of Hasselmann et al.\ (1985), taken from WAM (cycle 7).
A fortran version directly from the original WAM legacy code, 
and two python versions (one with jax/jit acceleration, one pure numpy) are available.
The attribute ''snonlin\_backend'' determines which one to use ('fortran'/'jax'/'numpy'), default is 'fortran';
all versions give identical results.


\subsection*{Stokes velocity}

Stokes velocity is  
\beq
\v u_S(z) =  \int \int 2 (k/\sigma) e^{2kz} \v k   E(\v k)  d \v k  
= (2/g) \int_0^\infty  \sigma^3  e^{2k z } \int_0^{2 \pi}   \v s E  d\theta d \nu 
\eeq
with direction vector $\v s=(\sin \theta,\cos \theta)$. The frequency integral is calculated over the
prognostic range of $\nu$, and the contribution by the high frequency tail is added.
The latter part is given by
  \beq \label{eq:ushi}
  \v u_S^{hi} = 
  (2/g) \int_{\nu_c}^\infty \int_0^{2 \pi}    \sigma^3  e^{2z \sigma^2/g}    \v s  E(\nu_c,\theta)  \nu_c^5  \nu^{-5} d\theta d \nu 
  \approx  (2/g)  \nu_c  \int_0^{2 \pi}   \sigma_c^3 \v s  E(\nu_c,\theta)  d\theta
%  \\  \approx 
 %   (2/g) \int_{\nu_c}^\infty \sigma^3 \nu_c^5  \nu^{-5} d \nu  \int_0^{2 \pi}   \v s  E(\nu_c,\theta) d\theta 
 %   =   (2/g)  (2 \pi)^3\nu_c^5 /\nu_c  \int_0^{2 \pi}   \v s  E(\nu_c,\theta)  d\theta
 \eeq 
 assuming that the uppermost model level thickness consumes the factor $ e^{2z \sigma^2/g} $
 and with $\sigma_c = 2 \pi \nu_c$. Below the uppermost vertical level, $  \v u_S^{hi} =0$ in this
 approximation.
 For small effective vertical decay scale of the tail contribution or large
 vertical levl thickness, this approximation
 gets worse.
The frequency integral in \eq{eq:ushi} can in fact be carried out
exactly, retaining the full depth dependence of the exponential. With
$\sigma_c = 2\pi \nu_c$ the result is
%
\begin{equation}
  u_S^{hi}(z)
  = \frac{2}{g}\, \nu_c \int_0^{2\pi}\! \mathbf{s}\,E(\nu_c,\theta)\,
    \sigma_c^{3}\left[\,
      e^{2z\sigma_c^{2}/g}
      - \sigma_c\sqrt{-\frac{2\pi z}{g}}\;
        \erfc\!\left(\sigma_c\sqrt{-\frac{2z}{g}}\right)
    \right] d\theta ,
  \label{eq:ushi-exact}
\end{equation}
%
valid for $z\le 0$. The approximation \eq{eq:ushi} is recovered in the
near-surface limit $\sigma_c^{2}|z|/g \ll 1$, where the $\erfc$ term
vanishes and $e^{2z\sigma_c^{2}/g}\to 1$, leaving the bracket equal to
$\sigma_c^{3}$ and hence the prefactor $\frac{2}{g}\nu_c\,\sigma_c^{3}$ of
\eq{eq:ushi}. Note that WAM (cycle 7) uses the approximated version;
the difference grows with the ratio of the uppermost level thickness to the
vertical decay scale of the tail.
 
\ifdetail
 
%
Starting from the high-frequency tail of \eq{eq:ushi} with the Phillips
form $E_{hi}(f)=E(\nu_c,\theta)\,\nu_c^{5}\nu^{-5}$, the per-direction frequency
integral to evaluate is
%
\beq
  I(z) \;=\; \int_{\nu_c}^{\infty}\sigma^{3}\,
              e^{2z\sigma^{2}/g}\,\nu_c^{5}\,\nu^{-5}\,d\nu ,
  \qquad \sigma = 2\pi \nu .
\eeq
%
Substituting $\sigma^{3}=(2\pi)^{3}\nu^{3}$ and writing
$b \equiv 8\pi^{2}z/g \le 0$ (since $z\le 0$, so $e^{b\nu^{2}}$ decays),
%
\beq
  I(z) \;=\; (2\pi)^{3} \nu_c^{5}\int_{\nu_c}^{\infty} \nu^{-2}\,e^{b \nu^{2}}\,d\nu .
\eeq
%
The remaining integral is evaluated by parts, differentiating $e^{b\nu^2}$
and integrating $\nu^{-2}$:
%
\beq
  \int_{\nu_c}^{\infty} \nu^{-2} e^{b \nu^{2}}\,d\nu
  = \Bigl[-\nu^{-1} e^{b \nu^{2}}\Bigr]_{\nu_c}^{\infty}
    + 2b\int_{\nu_c}^{\infty} e^{b \nu^{2}}\,d\nu
  = \frac{e^{b \nu_c^{2}}}{\nu_c}
    - \sqrt{-\pi b}\;\erfc\!\left(\nu_c\sqrt{-b}\right),
\eeq
%
where the boundary term at infinity vanishes ($b<0$) and the Gaussian
integral gives
$\int_{\nu_c}^{\infty} e^{b \nu^{2}}\,d\nu
   = \frac{1}{2}\sqrt{-\pi/b}\,\erfc(\nu_c\sqrt{-b})$.
Hence
%
\beq
  I(z) = (2\pi)^{3} \nu_c^{5}
         \left[\frac{e^{b \nu_c^{2}}}{\nu_c}
               - \sqrt{-\pi b}\;\erfc\!\left(\nu_c\sqrt{-b}\right)\right].
\eeq
%
Re-expressing in terms of $\sigma_c = 2\pi \nu_c$ using
$b \nu_c^{2} = 2z\sigma_c^{2}/g$,
$\;\nu_c\sqrt{-b} = \sigma_c\sqrt{-2z/g}\,$, and
$\;(2\pi)^{3}\nu_c^{5}\sqrt{-\pi b}
   = \sigma_c^{3}\nu_c\cdot\sigma_c\sqrt{-2\pi z/g}\,$ gives
%
\beq
  I(z) = \sigma_c^{3} \nu_c
         \left[\,e^{2z\sigma_c^{2}/g}
               - \sigma_c\sqrt{-\frac{2\pi z}{g}}\;
                 \erfc\!\left(\sigma_c\sqrt{-\frac{2z}{g}}\right)\right].
\eeq
%
Inserting $I(z)$ back, together with the prefactor $2/g$ and the angular
integral $\int_0^{2\pi}\mathbf{s}\,E(\nu_c,\theta)\,d\theta$, yields
\eq{eq:ushi-exact}. As $z\to 0^{-}$ the exponential tends to unity and
the $\erfc$ argument to zero with its $\sqrt{-z}$ prefactor, so
$I\to\sigma_c^{3}\nu_c$, recovering \eq{eq:ushi}.
 \fi
 
\subsection*{Wave stress}

The stress force in the ocean is given, with the dissipation energy source
$\dot E_{out} = \gamma_{out} E$ in $(\theta,\nu)$ coordinates, by
\beq \label{eq:dztau}
 \p_z \v \tau %= - \int Q \v k S_{out} d \v k 
 =  - (2 /g) \int_0^\infty  e^{2kz}  \sigma^3 \int_0^{2 \pi}  \v s  \dot E_{out}  d \theta d\nu 
  \eeq
The stress itself is 
\beq \label{eq:tauo}
 \v \tau_o = %- \int  \v k S_{out} d \v k = 
  -   \int_0^\infty\sigma   \int_0^{2 \pi}     \v s \dot E_{out}  d \theta d\nu
   \eeq 
The contribution  by the high-frequency tail assuming the WAMDI dissipation scale is
\beq \label{eq:dztau-hi}
 \p_z \v \tau^{hi}  
 %=  - (2 /g) \int_{\nu_c}^\infty  e^{2kz}  \sigma^3  \gamma_{out} \int_0^{2 \pi}  \v s  E  d \theta d\nu \\ 
 %=  2 \nu_c^5  \beta_{o} g^{-5}  \bar E^2 \bar \sigma^7 (2 \pi)^5 \int_{\nu_c}^\infty  e^{2kz}   d\nu
% \int_0^{2 \pi}  \v s   E(\nu_c,\theta)   d \theta \\
 =  2 \nu_c^5  \beta_{o} g^{-5}  \bar E^2 \bar \sigma^7 (2 \pi)^5 \int_{\nu_c}^\infty  e^{ (8 \pi^2/g) z  \,  \nu^2 }   d\nu
 \int_0^{2 \pi}  \v s   E(\nu_c,\theta)   d \theta 
  \eeq
   with $E_{hi}(\nu)  =   E(\nu_c,\theta) \nu_c^5 \nu^{-5}$ and 
   $\gamma_{out} = -  \beta_{o} g^{-4}  \bar E^2 \bar \sigma^7 \sigma^2 $ and the finite integral
   \beq \label{eq:erfint}
    \int_a^{\infty \:}e^{-bx^2}dx = \frac{ \sqrt{\pi }- \sqrt{\pi }\mbox{erf}\left(a\sqrt{b}\right)} { 2\sqrt{b}}
   \eeq
  For the stress we find for the high-frequency tail 
   \beq \label{eq:tauo-hi}
 \v \tau_o^{hi}  =    
 %  -   \nu_c^5 \int_{\nu_c}^\infty\sigma  \gamma_{out}  \nu^{-5} d\nu \int_0^{2 \pi}     \v s  E(\nu_c,\theta)d \theta 
 %\\ = 
  (2 \pi)^3 \nu_c^4 \beta_{o} g^{-4}  \bar E^2 \bar \sigma^7 
 \int_0^{2 \pi}     \v s   E(\nu_c,\theta)  d \theta 
   \eeq 

An alternative formulation of the tail stress follows from the equilibrium
assumption that underlies the imposed $\nu^{-5}$ shape itself: for $\nu > \nu_c$
the sources balance, $\dot E_{in} + \dot E_{nl} + \dot E_{out} = 0$, so the
flux to the ocean, $-\dot E_{out}$, can be evaluated as
$\dot E_{in} + \dot E_{nl}$ instead of extending the WAMDI dissipation scale
(which is calibrated on the resolved spectrum and underestimates the
dissipation that maintains the equilibrium range). The input part is
\beq \label{eq:tauo-hi-in}
 \v \tau_o^{hi,in} = \int_{\nu_c}^\infty \sigma\, \gamma_{in}\, \nu_c^5 \nu^{-5}
   \, d\nu \int_0^{2\pi} \v s\, E(\nu_c,\theta)\, d\theta
\eeq
with $\gamma_{in}$ from \eq{eq:gammas}. Writing $28 (u^\star/c)
\cos(\theta - \theta^*) = a \nu$ with $a(\theta) = 28 u^\star 2\pi
\cos(\theta-\theta^*)/g$, the input is active only above
$\nu_1 = \max(\nu_c, 1/a)$, and the frequency integral is elementary,
$\int_{\nu_1}^\infty (a \nu^{-2} - \nu^{-3})\, d\nu = a/\nu_1 - 1/(2 \nu_1^2)$, giving
per direction
\beq \label{eq:tauo-hi-in2}
 \v \tau_o^{hi,in} = \frac{\rho_a}{\rho_o} \frac{(2\pi)^2 \nu_c^3}{4}
 \int_0^{2\pi} \v s\, E(\nu_c,\theta)
 \left\{ \begin{array}{ll} a \nu_c - \frac{1}{2}, & a \nu_c \ge 1\\[2pt]
 \frac{1}{2}\, (a \nu_c)^2, & 0 < a \nu_c < 1 \end{array} \right\}  d\theta
\eeq
(continuous at $a \nu_c = 1$; zero for $\cos(\theta-\theta^*) \le 0$). The
integrand decays like $\nu^{-2}$, i.e.\ the integral converges exactly as
benignly as the WAMDI form in \eq{eq:tauo-hi}; note that the corresponding
integral with the dissipation of Janssen (1994; see also Komen et al.\ 1994,
and the ECMWF IFS documentation, Part VII) diverges strongly because of its
$(k/\bar k)^4$ term -- its $(k/\bar k)^2$ term behaves like the WAMDI form --
which is the reason the WAMDI scale is used in \eq{eq:tauo-hi}. The nonlinear part uses the conservation of the exact
transfer over the full spectrum, $\int_0^\infty \sigma \int \v s\,
\dot E_{nl}\, d\theta\, d\nu = 0$, so that
\beq \label{eq:tauo-hi-nl}
 \v \tau_o^{hi,nl} = \int_{\nu_c}^\infty \sigma \int_0^{2\pi} \v s\,
 \dot E_{nl}\, d\theta\, d\nu
 = - \int_0^{\nu_c} \sigma \int_0^{2\pi} \v s\, \dot E_{nl}\, d\theta\, d\nu
\eeq
i.e.\ the cascade flux into the tail equals minus the resolved-range
integral of the computed nonlinear transfer (the ``cascade credit''). For
completeness, the resulting total stress on the ocean in this
formulation is
\beq \label{eq:tauo-in-tot}
 \v \tau_o = - \int_0^{\nu_c} \sigma \int_0^{2\pi} \v s
 \left( \dot E_{out} + \dot E_{nl} \right) d\theta\, d\nu
 \; + \; \v \tau_o^{hi,in}
\eeq
where the first term is the resolved dissipation of \eq{eq:tauo} restricted
to the prognostic range with the cascade credit \eq{eq:tauo-hi-nl} folded
in, and the second term is the analytic tail input \eq{eq:tauo-hi-in2}.

The depth-distributed version (the ``stress work'' of \eq{eq:dztau}) follows
with the same replacement of the tail rate. With
$-\dot E_{out} = \gamma_{in} E(\nu_c,\theta) \nu_c^5 \nu^{-5}$ above $\nu_c$ and
$\sigma^3 \gamma_{in} = \frac{\rho_a}{\rho_o} \frac{(2\pi)^4}{4}
\nu^4 (a \nu - 1)$,
\beq \label{eq:dztau-hi-in}
 \p_z \v \tau^{hi,in} = \frac{2}{g} \frac{\rho_a}{\rho_o}
 \frac{(2\pi)^4}{4} \nu_c^5 \int_0^{2\pi} \v s\, E(\nu_c,\theta)
 \int_{\nu_1}^\infty e^{ (8\pi^2/g) z\, \nu^2 }
 \left( a - \nu^{-1} \right) d\nu\, d\theta
\eeq
with $\nu_1 = \max(\nu_c, 1/a)$ as before. The $a$--term is the finite integral
\eq{eq:erfint}; the second term requires the companion
\beq \label{eq:e1int}
 \int_{a_0}^{\infty} x^{-1} e^{-b x^2} dx = \frac{1}{2} E_1( a_0^2\, b )
\eeq
with the exponential integral $E_1$. The vertical integral of
\eq{eq:dztau-hi-in} recovers \eq{eq:tauo-hi-in2} exactly. The cascade
contribution \eq{eq:tauo-hi-nl} is deposited at the surface: it is
dissipated by tail waves with $\nu \ge \nu_c$, whose deposition scales
$g/2\sigma^2$ are centimetres, far below any ocean grid spacing. Both
formulations are implemented, selected by the switch \verb|tail_flux|
$\in$ \{\verb|'wamdi'|, \verb|'input'|\}; the discrete treatment, including
the diagnostic bins inside the grid, is described below.

\subsubsection*{Discretisation of the wave--ocean coupling and the energy exchange}

All three coupling quantities of \eq{eq_Q} -- the Stokes velocity $\v u_S$,
the stress force $\p_z \v \tau$ and the Doppler velocity $\v U$ -- involve the
same vertical kernel $Q(k,z) = 2 k e^{2kz}$ with $\int_{-\infty}^0 Q \, dz =
1$: $\v u_S$ and $\p_z \v \tau$ assign spectral quantities to a depth profile,
while $\v U$ projects the Lagrangian velocity back onto the spectral grid. The
exchange term $C$ of \eq{eq_C} closes because $\v U$ is the adjoint of the
profile assignment, i.e.\ because the same $Q$ appears on both sides of the
depth integral. The discretisation preserves this structure: one common
discrete kernel is used in all three methods of
\verb|./pyOM3/wim/interface.py| (\verb|calc_stokes|, \verb|calc_stress_work|,
\verb|calc_doppler_vel|), given by the cell integral of $Q$ over each ocean
cell of thickness $\Delta z$ with centres $z_k$ and faces $z_{k \pm 1/2} = z_k
\pm \Delta z/2$ (uppermost face at $z=0$),
\beq \label{eq:Qdz}
 \bar Q_k \Delta z = \int_{z_{k-1/2}}^{z_{k+1/2}} Q \, dz 
 = e^{2 k z_{k+1/2}} - e^{2 k z_{k-1/2}} 
\eeq
in place of the point value $Q(z_k) \Delta z$. The cell integral is essential:
the decay scale $g/(2\sigma^2)$ of the kernel is of the order of centimetres
for the high frequencies of the prognostic range and for the tail, far below
any ocean level thickness, so that point sampling loses a large part of the
coupling amplitude (about half of the stress force for $\Delta z = 1 \, m$ and
a developing wind sea). The weights \eq{eq:Qdz} telescope over the column,
$\sum_k \bar Q_k \Delta z = 1 - e^{-2 k D}$ for column depth $D$, with three
consequences that hold exactly for any $\Delta z$: the column-integrated
stress force equals the surface stress $\v \tau_o$ of \eq{eq:tauo}, the
discrete Stokes transport $\sum_k \v u_S(z_k) \Delta z$ equals its analytic
value, and a depth-independent (barotropic) current is projected onto the
Doppler velocity without attenuation, $\v U = \v u_L (1 - e^{-2kD})$. The
frequency derivative $\p_\nu \v U$ entering the radial refraction term
\eq{eq:xdot_disc} is taken as the analytic derivative of the discrete kernel,
$\p_\nu (\bar Q_k \Delta z) = (2k/\nu) \, 2 ( z_{k+1/2} e^{2 k z_{k+1/2}} -
z_{k-1/2} e^{2 k z_{k-1/2}} )$, so that refraction sees the same Doppler
projection as the advection.

For the high-frequency tails the same principle is applied to the exact tail
profiles. The tail of the stress force \eq{eq:dztau-hi} contains the frequency
integral $W(t) = \int_{\nu_c}^\infty e^{-b t \nu^2} d\nu$ with $t = -z \ge 0$ and $b
= 8 \pi^2/g$, evaluated with \eq{eq:erfint}; its antiderivative is
\beq \label{eq:tailH}
 H(t) = \int_0^t W(t') dt' = \frac{\sqrt{\pi}}{b} \left[ \sqrt{b t} \, 
 \mbox{erfc} \left( \nu_c \sqrt{b t} \right) 
 + \frac{1 - e^{-\nu_c^2 b t}}{\nu_c \sqrt{\pi}} \right]
\eeq
and each cell receives $\left( H(-z_{k-1/2}) - H(-z_{k+1/2}) \right)/\Delta z$
in place of the point value. Since $H(\infty) = 1/(b \nu_c)$, the column sum of
the tail force equals its exact depth integral $\v \tau_o^{hi}$ of
\eq{eq:tauo-hi}, using $2 (2 \pi)^5 \nu_c^5 g^{-5} H(\infty) = (2 \pi)^3 \nu_c^4
g^{-4}$; the prefactor $(2 \pi)^5$ in \eq{eq:dztau-hi} is essential for this
consistency. The Stokes tail \eq{eq:ushi-exact} is treated the same way: with
$a = 2 \sigma_c^2/g$ its depth profile $B(t) = \nu_c [ e^{-a t} - \sqrt{\pi a t}
\, \mbox{erfc}(\sqrt{a t}) ]$ has the antiderivative ($S = \sqrt{t}$)
\beq \label{eq:Istokes}
 I(t) = \nu_c \left[ \frac{1 - e^{-a t}}{a} - 2 \sqrt{\pi a} \left( 
 \frac{S^3}{3} \mbox{erfc}(\sqrt{a} S) 
 + \frac{1 - (a S^2 + 1) e^{-a S^2}}{3 \sqrt{\pi} \, a^{3/2}} \right) \right]
\eeq
with $I(\infty) = \nu_c/(3 a)$, so that the discrete tail transport equals the
exact value $(2\pi/3) \nu_c^2 \int \v s E(\nu_c,\theta) d\theta$. In tests the
column sums reproduce $\v \tau_o$ and the
Stokes transport to machine precision; for a shallow column below a
developed wind sea, however, a substantial part of the retained fraction $1 - e^{-2kD}$
of the wave momentum flux can be deposited below the column bottom. For a
closed momentum budget the attribute \verb|closed_bottom| extends the
integration of the bottom cell to $z \to -\infty$, in the shared kernel
($e^{2 k z_{bot}} \to 0$ in \eq{eq:Qdz}) as well as in both tail
antiderivatives ($H, I \to H(\infty), I(\infty)$), so that $\sum_k \bar Q_k
\Delta z = 1$ exactly: the bottom cell absorbs the deep part of all coupling
profiles, no wave momentum is lost, and the barotropic Doppler projection
becomes exact, $\v U = \v u_L$, at all frequencies (with the open bottom the
low-frequency bins of a shallow column see an almost vanishing Doppler
current). Since the kernel remains shared, the adjoint identity
\eq{eq:Cdisc} below is unaffected. The closure amounts to assuming $\v u_L$
constant below the column floor and is the recommended setting for shallow
test columns; for deep columns ($2 k D \gg 1$) it makes no difference.

The energy exchanged between ocean and waves is controlled by the same
kernel. On the ocean side, with the adjoint pair \eq{eq:Qdz} the discrete
identity
\beq \label{eq:Cdisc}
 \sum_k \v u^L_h(z_k) \cdot \p_t \v u_S(z_k) \, \Delta z 
 = \int \v U \cdot \v k \; \p_t A \; d \v k
\eeq
holds exactly (to machine precision) for any spectral tendency $\p_t A$ and
any $\Delta z$, over the prognostic frequency range (the tails are diagnostic
and remain outside this closure). The ocean therefore receives, per coupling
step, exactly the energy $\int \v U \cdot \v k \, \p_t A \, d \v k$
corresponding to the actual discrete evolution of the wave field, and the
handoff itself introduces no discretisation error.

On the wave side, the counterpart of this energy must be extracted from (or
added to) the intrinsic wave energy $\int \sigma A \, d \v k$ by the
propagation operator. For steady $\v U$ the rays conserve $\Omega = \sigma +
\v U \cdot \v k$, so the continuous propagation operator satisfies $\int
\Omega \, \p_t A |_{prop} \, d \v k = 0$ and the intrinsic energy change $\int
\sigma \, \p_t A$ balances the ocean gain $-\int \v U \cdot \v k \, \p_t A$
identically. Discretely, the flux-form advection conserves the action $\int A
\, d\v k$ to machine precision, but not $\int \Omega A \, d \v k$: the
nonlinear Superbee limiter leaves a residual (the frequency advection velocity
$\dot \nu$ is already located on the $u$-centred frequency faces, so face
staggering is not the cause). In tests with a wave packet refracted across a jet
(moderate Doppler shift, $\max |\v U \cdot \v k|/\sigma \approx 0.3$) this
$\Omega$-conservation defect amounts to about $1 \, \%$ of the gross local
transfer rate $\int |\dot \sigma A| \, d\v k$ for the Superbee scheme,
consistent with the $\mathcal{O}(0.5 \, \%)$ energy overshoot observed for such a
packet. First-order upwind advection is far worse (about
$12 \, \%$, since its frequency diffusion moves action across the $\sigma$
gradient), an unlimited second-order flux reduces the defect to about $0.2 \,
\%$. The $x$-sweep is exactly energy-neutral for a zonal jet; the defect
resides in the $y$, $\nu$ and $\theta$ sweeps.

The residual is removed exactly by a conservative energy fixer
(attribute \verb|energy_consistent| in \verb|propagation.py|). Writing the
discrete advection in flux form per phase-space grid cell $i$ as
$\pt{(A_i \Delta_i)} = -\sum_{faces} \delta F_i$, with the phase-space
volume $\Delta_i$ of the cell and $F$ the total (limited) action transport
through a cell face in $\v x$, $\nu$ or $\theta$, the transports cancel
pairwise in the sum over all cells and action is conserved to machine
precision for any advection scheme. Multiplying by $\Omega_i$, summing over
all grid cells and shifting the sums equivalent to a discrete partial
integration yields
\beq \label{eq:edefect}
 \pt{} \sum_i \Omega_i A_i \Delta_i = \sum_i \sum_{faces} F_i \, \delta \Omega_i
 \equiv D
\eeq
with the forward difference $\delta \Omega$ of $\Omega$ across each cell
face. Boundary terms vanish for closed and periodic boundaries (not for
reflective ones, for which no correction is made). The possible leakage $D$
is split into contributions $\sum_{j,k} D_{j,k} = D$ per water column at
discrete position $(j,k)$, with the shared horizontal face terms split
half-half between neighbouring columns, and each $D_{j,k}$ is cancelled by a
small corrective action transport $\tilde F_m$ across the frequency faces of
that column,
\beq \label{eq:efix}
 \tilde F_m = \lambda \, \bar A_m \, \delta_\nu \Omega_m ~~,~~
 \lambda(\v x) = - \frac{ D_{j,k} }
                       { \sum_m \bar A_m \, (\delta_\nu \Omega_m)^2 }
\eeq
with $m$ running over all $\nu$ faces and directions $\theta$ at position
$(j,k)$, and the weight $\bar A_m$ given by the face average of the action
content of the two adjacent cells. Being a transport, $\tilde F$ leaves the
action conservation untouched, but by the same partial integration it changes
the energy by $\sum_m \tilde F_m \, \delta_\nu \Omega_m = -D_{j,k}$. The
frequency direction is chosen for the corrective flux since it is the only
one in which $\delta \Omega$ does not vanish without currents
($\delta_\nu \Omega \approx \delta_\nu \sigma > 0$ for small Doppler
shift), and it is where the leakage is generated by action crossing the
frequency gradient of $\Omega$. The factor $\bar A_m$ ties the correction to
the level of action, so that it acts as a small correction of the frequency
propagation velocity rather than as a source term and positivity is not
threatened, and the factor $\delta_\nu \Omega_m$ makes the denominator
sign-definite, so that $\lambda$ exists whenever action is present at that
position. Minimising the action-weighted amplitude
$\sum_m \tilde{F}_m^2 / \bar A_m$ subject to this constraint yields the same
form by the standard Lagrange-multiplier condition: among all conservative
frequency-flux corrections that restore the energy, \eq{eq:efix} is the
smallest one in the action-weighted least-squares sense, i.e.\ the correction
perturbs the solution as little as possible.

The correction is local (one scalar per position $\v x$, no global
communication) and dormant without currents: for $\v U = 0$ we have
$\Omega = \sigma$, no frequency transports exist, and $D_{j,k} = 0$ so that
$\lambda = 0$. In tests the correction amounts to about $1 \, \%$ of the
frequency transports, the tendency-level defect drops from $0.9 \, \%$ of the
gross transfer to machine precision, and in long integrations of a wave
packet through a jet the energy $\int \Omega A \, d \v k$ is conserved
exactly (Adams--Bashforth time stepping preserves this, since each stored
tendency is corrected individually and $\Omega$ is constant while $\v U$ is
frozen between coupling updates), whereas without the fixer the spurious
drift amounts to roughly half of the genuine exchange. The transfer from
ocean to waves is thus exact on both sides; a structurally energy-conserving
advection would achieve the same without an a posteriori correction and
remains an option for future work.
   
 \subsection*{The coupling to the ocean}
 
 
 Stokes velocity, wave stress, and wave stress force are calculated in \verb|./pyOM3/wim/interface.py|.
 The methods there may be called and results passed to the ocean model.
 The velocity variable of the coupled ocean model is interpreted as the
 three-dimensional Lagrangian mean velocity $\v u_L = \v u + \v u_S$, with the
 Eulerian mean velocity $\v u$ and the Stokes velocity $\v u_S$, which is the most
 natural choice for energetic consistency (it also follows the example of interpreting
 the model's velocity as the isopycnal thickness weighted velocity to account for the
 effect of meso-scale eddies). Starting from the Craik--Leibovich equations of motion,
 the momentum equation for $\v u_L$ becomes
\beq \label{eqM}
\p_t \v u_L  -  \v u_L \times \vn \times (\v u_L-\v u_S)
 + f \v z \times \v u_L  = - \vn \pi + \p_t \v u_S + (\p_z \v  \tau,0)
\eeq
with the three-dimensional operator $\vn = (\vn_{\v x}, \p_z)$ and a gradient force
$\vn \pi$ collecting the pressure gradient and wave-driven gradient components.
The ocean model thus receives the wave-driven force $\p_t \v u_S$, the vortex force
\beq \label{eq:vf_cont0}
\v F = - \v u_L \times \vn \times \v u_S
\eeq
with $\v u_L \cdot \v F = 0$, and the wave momentum flux (stress), either as the
surface stress $\v \tau_o$ into the uppermost grid cell or as the interior stress
force $\p_z \v \tau$; the tracers are advected by $\v u_L$.
Note that the Coriolis term acts on the full $\v u_L$: with $w = 0$, identical
results are obtained using $\v u$ as model velocity and adding the Stokes--Coriolis
force $-f \v z \times \v u_S$ instead of $\p_t \v u_S$, but with $w \neq 0$ the
Lagrangian formulation is required to stay energetically consistent. The scalar
product of \eq{eqM} with $\v u_L$ contains no transfer terms related to the vortex
force and the Coriolis--Stokes force; the direct transfer to or from the waves is by
$\v u_L \cdot \p_t \v u_S$ and by the stress work $\v u_h^L \cdot \p_z \v \tau$,
which together form the exchange $C$ of \eq{eq_C}.

\subsubsection*{Discrete vortex force for horizontally uniform $\v u_S$ (nh\_simple)}

For a horizontally uniform Stokes velocity $\v u_S = (u_S(z),v_S(z),0)$, as used in
the non-hydrostatic model nh\_simple, the wave-driven force of \eq{eq:vf_cont0} becomes
\beq
\v F =  - \v u_L \times \vn \times \v u_S =
 \Vec{ \; w \, \pz{u_S} \\  \;w \, \pz{v_S} \\ -  u \,\pz{u_S} - v\, \pz{v_S} }
 ~~,~~\v u_L \cdot\v F = 0
 \label{eq:vf_cont}
\eeq
writing $\v u_L = (u,v,w)$ in this section for brevity.
On the C-grid, $u$, $v$ and $w$ are staggered at the cell faces
$(i+\frac{1}{2},j,k)$, $(i,j+\frac{1}{2},k)$ and $(i,j,k+\frac{1}{2})$,
and the vertical Stokes shear is placed on the $w$-levels,
\beq
S^u_{k+\frac{1}{2}} = \frac{ u_{S,k+1}-u_{S,k} }{\Delta z} ~~,~~
S^v_{k+\frac{1}{2}} = \frac{ v_{S,k+1}-v_{S,k} }{\Delta z}
\eeq
with $S^u = S^v = 0$ at the rigid lid and the flat bottom.
Introducing the two-point averages
$\Bar{ \phi }^{x}_{i+\frac{1}{2}} = ( \phi_{i} + \phi_{i+1} )/2$
and analogous in $y$ and $z$, the discrete force $\v F=(F^u,F^v,F^w)$ is
\beq
F^u_{i+\frac{1}{2},j,k} = \Bar{ \; \Bar{w}^{x} \, S^u \; }^{z} ~~,~~
F^v_{i,j+\frac{1}{2},k} = \Bar{ \; \Bar{w}^{y} \, S^v \; }^{z}
~~,~~
F^w_{i,j,k+\frac{1}{2}} =
 - \Bar{ \; \Bar{u}^{x} \; }^{z} S^u  - \Bar{ \; \Bar{v}^{y} \; }^{z} S^v
\label{eq:vf_disc}
\eeq
i.e.\ the products $\Bar{w}^{x} S^u$ are formed on the cell edges
$(i+\frac{1}{2},j,k+\frac{1}{2})$ and averaged back to the $u$-points,
while for $F^w$ the horizontal velocities are averaged first to the cell
centers and then to the $w$-levels where the shear is located.
The averaging operators in \eq{eq:vf_disc} are pairwise transposed:
each average that moves $w$ towards $u$ in $F^u$ is mirrored by an average
moving $u$ towards $w$ in $F^w$. With the summation-by-parts identity for
the two-point average
\beq
\sum_i \psi_{i+\frac{1}{2}} \, \Bar{\phi}^{x}_{i+\frac{1}{2}}
= \sum_i \phi_{i} \, \Bar{\psi}^{x}_{i}
\label{eq:sbp}
\eeq
which holds without boundary terms in the periodic horizontal directions,
and in the vertical because the boundary contributions carry the factor
$S^u = S^v = 0$ at the lid and bottom faces, the total work of the discrete
force vanishes exactly,
\beq
\sum_{i,j,k} \lk u F^u + v F^v + w F^w \rk = 0
\eeq
for arbitrary $\v u_L$ and arbitrary profiles $u_S(z)$, $v_S(z)$; this has
been verified numerically to machine precision for random fields. Note that
the work vanishes only in the sum.

The vanishing work is a property of the spatial operator and does not by
itself guarantee numerical stability. Since $\v F$ is perpendicular to
$\v u_L$, it turns the velocity vector without changing its magnitude,
i.e.\ it acts locally as a rotation with angular frequency $|\pz{\v u_S}|$,
analogous to inertial oscillations with frequency $f$. An explicit
third-order Adams--Bashforth time stepping scheme is stable for such neutral
oscillations only if the angle covered per time step remains below
$|\pz{\v u_S}| \Delta t \lesssim 0.72$. Since the Stokes shear is largest in
the uppermost grid cell and grows with the developing sea, this number
should be monitored during integrations together with the advective Courant
numbers.

\subsubsection*{Discrete vortex force in the hydrostatic case (pe\_simple)}

In the hydrostatic model pe\_simple the Stokes velocity $\v u_S = (u_S,v_S,0)$
varies in all three dimensions, and the force \eq{eq:vf_cont0} is truncated to
\beq
\v F = \Vec{ - v \, \zeta_S \\ \; u \, \zeta_S }
~~,~~ \zeta_S = \px{v_S} - \py{u_S}
~~,~~ \v u_h^L \cdot \v F = 0
\label{eq:vf_pe}
\eeq
with $\v u_h^L = (u,v)$ in this section for brevity.
The Stokes-shear terms $w \pz{u_S}$ and $w \pz{v_S}$, part of the first and
second component of the total force $-\v u_L \times \vn \times \v u_S$,
respectively, are deliberately omitted. This is because the scalar product
of these terms with $\v u_h^L$ is only cancelled by $w$ times the vertical
force component $-u \pz{u_S} - v \pz{v_S}$, which the hydrostatic model does
not carry. Retained alone, the Stokes-shear terms would thus perform the
uncancelled work $w \lk u \pz{u_S} + v \pz{v_S} \rk$ on $\v u_L$.
The omitted terms are the so-called Stokes shear force of Suzuki and
Fox-Kemper (2016). Their neglect is also in agreement with quasi-geostrophic
scaling where the Rossby number $Ro$ becomes small: within the Stokes layer
of depth $h_S$ the vertical velocity of balanced flow is limited by the
surface kinematic condition and its $O(Ro)$ divergence to
$|w| \lesssim Ro \, U h_S/L$, such that
$|w \pz{u_S}| \sim Ro \, U U_S / L$ independent of $h_S$ and smaller by
$O(Ro)$ than the retained vortex force.

On the C-grid, $u_S$ and $v_S$ are given at the cell centers and averaged
to the $u$- and $v$-points, and the Stokes vorticity is formed at the cell
corners $(i+\frac{1}{2},j+\frac{1}{2})$ by centered differences,
\beq
\zeta_S = \frac{ \Bar{v_S}^{y}_{i+1} - \Bar{v_S}^{y}_{i} }{\Delta x}
        - \frac{ \Bar{u_S}^{x}_{j+1} - \Bar{u_S}^{x}_{j} }{\Delta y}
\eeq
The products are formed at the corners first,
and then averaged back to the velocity points,
\beq
F^u_{i+\frac{1}{2},j} = - \Bar{\zeta_S \, \Bar{v}^{x} }^{y} ~~,~~
F^v_{i,j+\frac{1}{2}} = \; \Bar{ \zeta_S \, \Bar{u}^{y}}^{x}
\label{eq:vf_pe_disc}
\eeq
with $\v F = (F^u,F^v)$.
Applying the summation-by-parts identity \eq{eq:sbp} once in $y$ to the
first and once in $x$ to the second term of the work moves both averages
onto the velocities, and the two contributions meet at the corners,
\beq
\sum u F^u + \sum v F^v
= \sum_{\mathrm{corners}} \zeta_S
 \lk - \Bar{v}^{x} \, \Bar{u}^{y} + \Bar{u}^{y} \, \Bar{v}^{x} \rk = 0
\eeq
i.e.\ the cancellation is not only global as in nh\_simple but holds
corner by corner. The boundary terms of \eq{eq:sbp} vanish because the
velocities vanish on the channel walls and on land, so the identity holds
in the presence of topography; this has been verified to machine precision
for random fields on a masked channel.

\subsubsection*{Discrete aspects of the momentum coupling}

The discretisation of the remaining coupling terms in the ocean model
(C-grid) is chosen such that the energetic properties of the continuous
system are preserved where possible, and monitored where not.

\emph{Stokes--Coriolis force.} Continuously, the combined rotational forces are
$(f+\zeta)\hat z\times\v u_L$, which does no work against $\v u_L$. For the
Coriolis part this property carries over exactly to the discrete level if the
Stokes velocity is first interpolated from T to the U/V points and then passed
through the \emph{same} discrete Coriolis operator $C$ as the Eulerian
velocity (the energy-conserving form of the main code, with $f$ evaluated at
T points inside the adjoint-paired averaging). By linearity the total discrete
Coriolis force is then $C(\v u_L)$, and the antisymmetry of the operator gives
$\sum \v u_L\cdot C(\v u_L)=0$ to machine precision for arbitrary $f(y)$. Note
that adding instead a separate Stokes--Coriolis term with its own (single)
T-to-U/V averaging leaves an $O(\Delta x^2)$ residual in the $\v u_L$ energy
budget, since $\v u_S$ then enters the budget once as a T-point field and once
through a double interpolation; the difference is a commutator between the
$1$-$2$-$1$ smoothing and $f$.

\emph{Vortex force.} With the discretisations of the previous section the
$\v u_L$-work of the applied vortex force vanishes exactly (globally in
nh\_simple, corner by corner in pe\_simple). The coupled driver nevertheless
computes $\sum \v u_L\cdot \v F$ (volume integrated, with the vortex force
exactly as applied) every coupling step and stores it in the coupling
output, so that it can be verified to remain at round-off level.

\emph{Stokes advection of tracers.} The advection of buoyancy by $\v u_S$ is
added in flux form: $\v u_S$ is averaged from T to the U/V faces, the vertical
Stokes velocity $w_S$ is diagnosed from the discrete Stokes continuity
equation integrated upward from the bottom ($w_S(-D)=0$), and the fluxes are
differenced like the model's own tracer fluxes, with zero buoyancy flux
through the surface face. This form is exactly conservative (the fluxes
telescope, so $\int b\,dV$ is preserved to round-off) and
constancy-preserving in the interior (a constant $b$ yields zero tendency,
since the diagnosed $w_S$ closes the three-dimensional Stokes divergence
discretely); both properties were verified to machine precision. In the top
cell the surface no-flux condition places the convergence of the Stokes
transport (the Stokes pumping $w_S$ at the surface face) into the cell budget,
which is the physically correct location: it is compensated by the Eulerian
anti-Stokes flow driven by the Stokes--Coriolis force. Centered fluxes are
used; the Stokes Courant number is $O(10^{-3})$, far below the regime where
flux limiting matters.

\emph{Stress and vertical placement.} The wave stress is kinematic
(the spectrum is the variance density, so no density factor appears; the
air--water density ratio enters through the growth rate of the input source
term) and is applied either to the surface face of the top wet cell
($\tau/\Delta z$) or depth-distributed with the same cell-integrated vertical
kernel that underlies the Stokes velocity and the Doppler projection, which
makes the discrete energy-exchange identity (the $C$-closure of the previous
section) hold to machine precision. Note that with a top-cell thickness large
compared to the Stokes decay scale $1/(2k)$, the stored top-cell Stokes
velocity is a cell average and correspondingly smaller than the surface value;
resolving the near-surface Stokes shear requires $\Delta z$ of a few metres.

\emph{High-frequency tail fluxes to the ocean (tail\_flux switch).} Above
the prognostic cutoff $\nu_c$ the spectrum is diagnostic ($E \propto \nu^{-5}$,
imposed by the tail parameterisation), and its momentum and energy fluxes to
the ocean must be added analytically to the resolved-range dissipation
integrals. Two formulations are implemented, selected by the attribute
\verb|tail_flux|:

(i) \verb|'wamdi'| (default): the WAMDI dissipation rate
continued over the tail,
$\int_{\nu_c}^\infty \gamma_{out}\,E_{hi}\, \sigma\, d\nu$ with $\gamma_{out} \propto
\bar\sigma (\sigma/\bar\sigma)^2 (\bar\alpha/\alpha_{PM})^2$. The integrand
decays like $\nu^{-2}$ and converges; note that the corresponding integral with
the dissipation of Janssen (1994) diverges strongly because of its
$(k/\bar k)^4$ term (its $(k/\bar k)^2$ term behaves like the WAMDI form),
which is why the WAMDI form is used for this term. However,
$\gamma_{out}$ is the whitecapping rate calibrated on the resolved spectrum and
underestimates the dissipation that actually maintains the $\nu^{-5}$
equilibrium range --- the very reason the tail must be imposed diagnostically
--- so this form systematically underestimates the tail fluxes (in tests by
a substantial factor in the momentum flux).

(ii) \verb|'input'|: in the equilibrium range $S_{in} + S_{nl} + S_{out} = 0$
holds, so the flux to the ocean, $-S_{out}$, equals $S_{in} + S_{nl}$. The
input part is the analytic integral of the WAMDI growth rate over the tail:
with $a(\theta) = 28 u_\star 2\pi \cos(\theta-\theta_w)/g$ (the input
parameter evaluated at $\nu_c$), the input being active only above
$\nu_1 = 1/a$,
\beq
I_\tau(\theta) = \frac{\epsilon}{4} (2\pi)^2 \nu_c^3\, E_c(\theta)
\left\{ \begin{array}{ll} a \nu_c - \frac{1}{2}, & a \nu_c \ge 1\\[2pt]
\frac{1}{2} (a \nu_c)^2, & 0 < a \nu_c < 1 \end{array} \right.
\qquad
I_\Phi(\theta) = \frac{\epsilon}{4} (2\pi) \nu_c^2 g\, E_c(\theta)
\left\{ \begin{array}{ll} \frac{a \nu_c}{2} - \frac{1}{3}, & a \nu_c \ge 1\\[2pt]
\frac{1}{6} (a \nu_c)^3, & 0 < a \nu_c < 1 \end{array} \right.
\eeq
($\epsilon = \rho_a/\rho_w$, $E_c(\theta) = E(\nu_c,\theta)$; continuous at
$a \nu_c = 1$; zero upwind; the integrands again decay like $\nu^{-2}$). The same booking applies to the
\emph{in-grid} diagnostic bins (between the per-column prognostic cut-off
$m_0$ and the last grid frequency, where the $\nu^{-5}$ shape is imposed):
there the flux is evaluated discretely as $\gamma_{in} E$ with the actual
imposed spectrum, while the prognostic bins keep the true $\gamma_{out} E$
dissipation --- booking the in-grid diagnostic range at $\gamma_{out}$ would be
inconsistent with the equilibrium assertion that the imposed tail itself
encodes (and underestimates the flux substantially: the in-grid diagnostic
range can carry several times the beyond-grid tail). The nonlinear part uses the
conservation of the exact transfer over the full spectrum,
$\int_{tail} S_{nl}\, d\Omega = -\int_{resolved} S_{nl}\, d\Omega$: the
\emph{prognostic}-range integrals of the stored $S_{nl}$ (momentum weight
$\sigma$, energy weight $g$) are credited to the ocean fluxes (the ``cascade
credit''; the in-grid diagnostic $S_{nl}$ is part of the equilibrium booking
there, and the algebra of the two ranges collapses exactly to the
prognostic-masked integral).
On the truncated grid the discrete DIA is not exactly conservative at the
cutoff; the residual is deliberately included in the credit --- whatever the
discrete operator removes from the resolved range is deposited in the tail.
For the depth-distributed stress (method \verb|calc_stress_work|, as opposed
to the surface stress of \verb|calc_stress|) the tail is
deposited with its exact profile: with $b = 8\pi^2/g$, $t=-z$,
$\beta_L = b\,t \nu_L^2$, the cumulative-from-surface tail stress is
$V(t) = C_0 [\, a/\nu_L - 1/(2\nu_L^2) - (a J_2 - J_3)\,]$ with
$J_2 = e^{-\beta_L}/\nu_L - \sqrt{\pi b t}\,\mathrm{erfc}\sqrt{\beta_L}$ and
$J_3 = \frac{1}{2}( e^{-\beta_L}/\nu_L^2 - b t\, E_1(\beta_L))$, so the cell
weights $(V(t_{bot})-V(t_{top}))/\Delta z$ sum to the surface tail exactly;
the cascade credit is deposited in the uppermost cell (its carrier waves
have deposition scales $g/2\sigma^2$ of centimetres).

\emph{Dissipation energy flux to the ocean.} The term
$\int \sigma S_{out}\, d\v k$ of the wave energy budget \eq{eq_W} is
the energy the breaking waves hand to near-surface turbulence. It is diagnosed
locally as
\beq
\Phi_{oc}(\v x) = -\,g \int \dot E_{out}\, d\nu\, d\theta \;\ge\, 0
\eeq
(kinematic; $\rho_w \Phi_{oc}$ in W\,m$^{-2}$; WAM's PHIOC), computed from the
same stored dissipation rate $\gamma_{out}$ and with the same spectral
measures as the wave stress, and written to the coupling output both as the
local field and as its domain integral in the energy-diagnostic record. The
dissipation of the diagnostic $\nu^{-5}$ tail is included according to the
\verb|tail_flux| switch described above (under \verb|'input'| together with
the nonlinear cascade credit). The
field is the surface boundary condition for the flux $F$ of the turbulent
kinetic energy (TKE) closure following Gaspar et al.\ (1990) in the ocean
component ($\rho_w \Phi_{oc}$ as the surface TKE flux, in place of the
$u_\star^3$-based parameterisation), which routes this term into the ocean's
turbulent kinetic energy budget and thereby closes the coupled energy cycle.
The TKE closure receives, besides the usual shear production by the Eulerian
mean flow, a shear production term related to the mean Stokes velocity,
$c_u K (\p_z \bar{\v u}) \cdot \p_z \bar{\v u}_S$, resulting from the
down-gradient closure of the eddy stresses; the corresponding energy is
exchanged with the mean Lagrangian kinetic energy.

\emph{Vertical friction and mixing.} For consistency with the TKE closure,
the vertical friction operator acts on the Eulerian mean velocity,
$\p_z A_v \p_z (\v u_L - \v u_S) = \p_z A_v \p_z \v u$: the work of this
viscous term against $\v u_L$ then exactly matches the shear production
terms of the TKE closure, including the Stokes-shear production. Mixing the
Lagrangian instead of the Eulerian momentum would break this energy
consistency.

\emph{Coupling frequency and the exchange term.} The two models exchange
fields once per coupling interval $\Delta t_c$, chosen as a common multiple of
the two time steps: the component with the smaller time step sub-cycles
$\Delta t_c/\Delta t$ steps between exchanges. When the wave step is the
smaller one, $\Delta t_c = \Delta t_{\rm ocean}$ and the wave model
sub-cycles; when $\Delta t_{\rm ocean} < \Delta t_{\rm wave}$ (e.g.\ a
non-hydrostatic or very high resolution ocean) the roles reverse and the ocean
sub-cycles, with $\Delta t_c = \Delta t_{\rm wave}$. In either case all
coupling fields are piecewise constant over the coupling interval, so the
exchange work $C=\int \v u_L\cdot\p_t\v u_S\,dV$ is evaluated with
$\p_t\v u_S$ as the difference quotient across the coupling update, i.e.\ over
$\Delta t_c$ rather than over the time step of either component. The
domain-integrated energy diagnostics --- wind input to the waves
$g\int\gamma_{\rm in}E\,d\Omega\,dA$, the stress work on the ocean (computed
with the identical stencils as the applied forcing), the exchange term $C$,
and the vortex-force work --- are written to the coupling output at every
coupling step, so the discrete energy pathway of the coupled system is
monitored throughout the integration rather than assumed. The choice of
$\Delta t_c$ is constrained by the physics rather than by stability: the
coupling fields should be well sampled, i.e.\ $\Delta t_c$ small compared to
the evolution time scales of the exchanged quantities (translation of the wind
systems across a grid cell for the stress and Stokes fields, and the evolution
of the resolved ocean currents for the refraction terms).

  
\subsection*{Propagation, boundary conditions, land mask}

Action is propagated since it avoids to calculate the energy exchange with the mean flow
explicitely. 
Advection uses the second-order scheme with Superbee limiter (\verb|tracer_flux_superbee|)
taken from MITgcm. 
A first-order upwind option is also available (\verb|tracer_flux_upwind|,
selected by the attribute ''f\_limiter''$=$''upwind''); it shares the identical
finite-volume (face-flux) structure --- the same land masking, halo exchange,
closed frequency faces and reflection --- so conservation and the energy fixer
of \eq{eq:efix} still apply, but the flux limiter is set to zero. Upwind is more
diffusive (it smears sharp swell fronts) but robust to grid-scale noise.
The refraction drift velocities $\dot \nu$ and $\dot \theta$ can optionally be
clipped to a CFL bound (attribute \verb|limit_fdot|, default off, with the
bound set by \verb|refraction_cfl|, default $0.5$): strong current shear can
otherwise drive the spectral advection beyond its stability limit locally.
The clipping breaks the exact ray equations where it engages and is meant as
a safeguard, not a substitute for a proper time step.



Lateral boundary conditions in $x$ and $y$ are periodic (attribute ''periodic\_in\_x''/''periodic\_in\_y''),
closed, or reflective. Boundary conditions in frequency are always closed, in direction
always periodic.
Lateral 
reflection boundary conditions at coast lines are implemented  the same way as in the 
IDEMIX2 in pyOM2.
These are activated by setting the attribute ''enable\_reflection'' to true value'. 
If undefined or false, the default of closed boundary are used where energy may pile up near the 
boundary.
Land mask is set in the template method   ''set\_topography'' in \verb|./pyOM3/wim/main.py|.


\subsection*{Time stepping scheme}

Each time step is split into an explicit propagation substep for the wave action
$A$ followed by an implicit source substep for the wave energy $E = \sigma A$.
The propagation tendency $\dot A|_{\rm prop}$ (the sum of the flux divergences in
$x$, $y$, $\nu$ and $\theta$) is advanced with a multi-step Adams--Bashforth
scheme,
\beq \label{eq:ab3}
A^{n+1} = A^{n} + \Delta t\left( c_A\, \dot A^{\,n}|_{\rm prop}
                              + c_B\, \dot A^{\,n-1}|_{\rm prop}
                              + c_C\, \dot A^{\,n-2}|_{\rm prop}\right).
\eeq
The scheme used is the stabilised second- or third-order Adams--Bashforth  pair
\beq \label{eq:ab2eps}
(c_A,c_B,c_C) &=& \left(\frac{3}{2}+\epsilon,\; -\frac{1}{2}-\epsilon,\; 0\right),
\qquad \epsilon \approx 0.05, ~~\mbox{for second order}
\\
(c_A,c_B,c_C) &=& \left(\frac{23}{12},\; -\frac{16}{12},\; \frac{5}{12}\right)
~~\mbox{for third order}
\eeq
both started with a forward Euler step $(1,0,0)$. The time loop uses the
AB2$+\epsilon$ pair; the third-order coefficients are given for reference but
are deliberately not used (see below). The recent tendencies are held in a
rotating three-level buffer (the pointers ''tau'', ''taum1'', ''taum2''). Only
the propagation is integrated this way; the source terms are treated separately
below. After the propagation update the action is subject to the lateral
boundary conditions and, on land, is masked to zero.

The choice of the time scheme for the propagation is more delicate than usual
because the advection uses a flux \emph{limiter}. TVD limiters (Superbee,
upwind) guarantee a monotone, non-oscillatory update --- in particular
$A \ge 0$ --- only when the limited flux is applied in a \emph{single forward
Euler step} (or a strong-stability-preserving Runge--Kutta combination of such
steps). Multi-step Adams--Bashforth schemes combine tendencies from several time
levels with \emph{negative} weights, which destroys this guarantee at any CFL
number: each stored tendency is limited with respect to its own state, but the
negatively-weighted combination is free to over- and undershoot. For the
third-order scheme $(23/12,\,-16/12,\,5/12)$ the negative weight is large and
the violation is severe: a smooth, sign-definite action field develops
undershoots ($A < 0$) in an alternating grid-scale ($2\Delta x$) checkerboard
pattern at sharp fronts (storm edges, swell fronts). Unlike in the ocean
model --- where the same combination is harmless because tracers are not
sign-definite, explicit (biharmonic and vertical) diffusion damps grid-scale
noise every step, and no term amplifies small wiggles --- the wave model has no
explicit diffusion, its action must remain non-negative, and the source terms
amplify any $\pm$ pattern strongly non-linearly (through
$\alpha = \bar E\,\bar\sigma^4$ and the cubic DIA transfer), so the checkerboard
grows over $\mathcal{O}(1)$ day until floating-point overflow. The AB2$+\epsilon$
coefficients \eq{eq:ab2eps} reduce the negative weight to $-(1/2+\epsilon)$,
which in practice (CFL${}\lesssim 0.25$) is small enough for the limiter's
built-in slack to absorb the violation: the action remains non-negative to
machine precision in long integrations where AB3 fails within a day. This is an
empirical rather than a rigorous property; forward Euler $(1,0,0)$ or an SSP
Runge--Kutta scheme restores the formal TVD guarantee at the cost of temporal
accuracy or an additional flux evaluation per step.

Two safeguards back this up. First, a positivity floor (attribute
\verb|floor_action|, default on) clips the action at zero after each
propagation step, as is standard in WAM. The clipped mass --- the action the
floor removes, which equals the spurious action it injects and hence the
energy-consistency cost of the clip --- is measured \emph{before} clipping and
a warning is issued whenever it exceeds round-off scale relative to the total
action, so a silently active floor cannot mask a genuine positivity violation;
switching the floor off restores the raw scheme. Second, $\min A$ is monitored
automatically every $O(100)$ steps against a relative round-off threshold and a
warning names the remedy (larger $\epsilon$ or forward Euler): any
limiter/time-scheme violation shows up there long before a blow-up.


Writing the source terms as growth rates acting on the energy, the source substep
integrates $\partial_t E = S$ with $S = \dot E_{nl} + (\gamma_{in} +
\gamma_{out})E$, where $\dot E_{nl}$ is the non-linear (DIA) transfer. A fully
explicit update would require a very small time step because the growth rate
increases with frequency (this is also the reason for the high-frequency
cut-off). Instead a linearised-implicit (semi-implicit) step is used: linearising
$S$ about the current state, $S(E^{n+1}) \approx S(E^n) + \partial_E S\,
(E^{n+1}-E^n)$, and solving for $E^{n+1}$ gives
\beq \label{eq:implicit_source}
E^{n+1} = E^{n} + \frac{S\, \Delta t}{1 - \Delta t\, \partial_E S},
\qquad
\partial_E S = \gamma_{in} + \gamma_{out} + \partial_E \dot E_{nl},
\eeq
where the diagonal of the non-linear functional derivative
$\partial_E \dot E_{nl}$ is returned by the DIA backend alongside the transfer
itself. This is unconditionally stable for a purely dissipative feedback
($\partial_E S < 0$), for which the denominator exceeds unity and the update is
damped. The growth can additionally be bounded as in WAM by a growth limiter
(attribute \verb|limit_growth|, default off) that caps $S$ (not $\partial_E S$)
so that the energy cannot grow by more than a prescribed fraction per step,
\beq \label{eq:growthlimit}
S \leftarrow \min\!\left(S,\; 10^{-7}\, g\,
\max(u_\star,\,u_\star^{\min})\, \bar\sigma\,/\,(2\pi \nu^4)\right),
\qquad u_\star^{\min} = 0.01\ {\rm m\,s^{-1}}.
\eeq
The friction-velocity floor $u_\star^{\min}$ in the cap is essential and easy
to overlook: without it, the cap vanishes wherever the wind vanishes and the
limiter degenerates to $S \leftarrow \min(S,0)$ in calm cells. This clips only
the \emph{positive} lobes of the nonlinear transfer, destroying the
conservative $\pm$ structure of $S_{nl}$ exactly there, and acts as a small
but systematic negative-action source. In a calm region that accumulates swell
without a sink (e.g.\ a walled basin with reflective coasts outside the wind
band), the resulting violation grows linearly in time --- independent of the
time step and of the time-stepping scheme --- until, after $O(100)$ days, the
nonlinear source terms of the accumulated spectrum activate and the
integration fails abruptly. The floor keeps the cap strictly positive
everywhere while leaving the physics unchanged: $u_\star^{\min}$ is far below
the input threshold of the slowest resolved wave ($c \approx 28\,u_\star$
requires $u_\star \gtrsim 0.06$ m\,s$^{-1}$ for the highest resolved
frequency), so the wind input remains identically zero in floored cells; only
the growth cap becomes a finite positive bound instead of zero.
Because the limiter bounds only the numerator, the denominator
$1 - \Delta t\, \partial_E S$ can still vanish or change sign where
$\partial_E S$ is large and positive --- most notably at high frequency during
spin-up from a small seed, where the DIA derivative is steep. When this happens
the update overshoots and the highest bins become unstable. Two safeguards are
used: the high-frequency cut-off removes the steepest bins from the prognostic
integration (replacing them by the diagnostic $\nu^{-5}$ tail), and, optionally,
only the dissipative part of the feedback is kept in the implicit denominator,
\beq \label{eq:denomclamp}
\partial_E S \;\rightarrow\; \min\!\left(\partial_E S,\, 0\right),
\eeq
guaranteeing $1 - \Delta t\, \partial_E S \ge 1$. This is activated 
by setting the attribute  \verb|implicit_dissipation_only|  to
\verb|True| and equals  the standard WAM/WW3
construction --- the stabilising dissipative feedback is treated implicitly while
the growth is treated explicitly and bounded by the limiter. It leaves the
equilibrium ($S = 0$) unchanged and therefore does not affect the energy balance;
it only alters the transient integration of the growth phase.

After the source substep the energy is masked to zero on land, the
high-frequency tail is re-imposed through \verb|imphftail|, and the action is
recovered as $A = E/\sigma$ for the next propagation step.


The time step of the propagation is limited by the advective CFL number based
on the fastest group velocity, i.e.\ the lowest resolved frequency (Superbee is
stable to CFL${}\lesssim 0.5$, upwind to CFL${}<1$; in practice CFL${}\approx
0.2$--$0.25$ is used, which is also the regime in which the AB2$+\epsilon$
scheme preserves the non-negativity of the action, see the time-stepping
section). Note that the grid-scale checkerboard described there is \emph{not} a
CFL violation --- it appears at any Courant number when a multi-step scheme
with a large negative weight (AB3) is combined with the flux limiter --- so it
cannot be cured by reducing the time step; the cure is the time scheme.
Substantially larger time steps than the advective CFL allows would require an
implicit or semi-Lagrangian treatment of propagation, at the cost --- for plain
semi-Lagrangian --- of the exact (face-flux) energy conservation.


\subsection*{Model parameters and options}

All parameters and options are attributes of the model instance. They are set
either through the parameter dictionary passed at construction or by assigning
the attribute on the instance before the model is stepped; options marked with
a default may simply be left unset. The tables below list every attribute the
wave model reads, with its default where one exists.

\subsubsection*{Required parameters (no default)}

\begin{longtable}{@{}llp{8.6cm}@{}}
\hline
Attribute & Unit & Meaning \\
\hline
\endhead
\verb|nx|, \verb|ny| & --- & global horizontal grid size (cells, without halo
  points) \\
\verb|dx|, \verb|dy| & m & horizontal grid spacing \\
\verb|nf| & --- & number of frequency bins (local on every PE) \\
\verb|nd| & --- & number of direction bins (local on every PE) \\
\verb|f0| & 1/s & lowest frequency of the geometric frequency grid (WAM value
  $0.035$) \\
\verb|df0| & --- & relative frequency increment $\Delta \nu/\nu$ of the
  geometric grid (WAM value $0.1$) \\
\verb|dt| & s & time step \\
\verb|periodic_in_x| & --- & lateral boundary in $x$: periodic if true, closed
  if false (the outermost ring of cells is then set to land) \\
\verb|periodic_in_y| & --- & as above for $y$ \\
\hline
\end{longtable}

\noindent
The constructor additionally takes the keyword arguments \verb|n_pes_x|,
\verb|n_pes_y| (processor grid, default $1\times 1$; their product must equal
the communicator size), \verb|halo_points| (default 1) and \verb|comm_in|
(default: the MPI world communicator).

\subsubsection*{Forcing and coupling fields}

\begin{longtable}{@{}llp{8.6cm}@{}}
\hline
Attribute & Default & Meaning \\
\hline
\endhead
\verb|ustar| & required & friction velocity magnitude $u_\star$; scalar or
  horizontal field (scalars are broadcast) \\
\verb|udir| & required & wind direction $\theta^*$ (radians); scalar or field \\
\verb|ocean_U|, \verb|ocean_V| & unset & Doppler velocity $\v U(\nu)$ provided
  by the coupler. If unset, all current effects (advection by currents,
  refraction, the energy fixer) are dormant \\
\verb|ocean_Udf|, \verb|ocean_Vdf| & $0$ & frequency derivative
  $\p_\nu \v U$ of the Doppler velocity (radial group-velocity correction
  $R$ and frequency refraction) \\
\hline
\end{longtable}

\subsubsection*{Options with defaults}

\begin{longtable}{@{}llp{8.2cm}@{}}
\hline
Attribute & Default & Meaning \\
\hline
\endhead
\verb|snonlin_backend| & \verb|'fortran'| & DIA implementation:
  \verb|'fortran'| (original WAM code; requires double precision),
  \verb|'jax'| or \verb|'numpy'|; all give identical results \\
\verb|grav| & $9.81$ & gravitational acceleration \\
\verb|rho_o| & $1024$ & reference water density. Currently unused inside the
  wave model: the air--water density ratio of the source terms is fixed at
  $1.225/1000$ \\
\verb|z0| & $0.0002$ & roughness length; read only by the optional WAM input
  parameterisation, inactive with the default WAMDI input \\
\verb|mask_energy| & off & prognostic cut-off $\nu_c$ per column with the
  imposed $\nu^{-5}$ diagnostic tail (\verb|imphftail|); prerequisite for the
  equilibrium booking of \verb|tail_flux| $=$ \verb|'input'| \\
\verb|limit_growth| & off & WAM growth limiter, \eq{eq:growthlimit} \\
\verb|implicit_dissipation_only| & off & keep only the dissipative part of the
  source feedback in the implicit denominator, \eq{eq:denomclamp} \\
\verb|floor_action| & on & positivity floor $A \leftarrow \max(A,0)$ after the
  propagation step, with the clipped mass monitored and warned about \\
\verb|f_limiter| & \verb|'superbee'| & advection flux limiter:
  \verb|'superbee'| (second order) or \verb|'upwind'| (first order,
  diffusive, robust) \\
\verb|energy_consistent| & off & conservative energy fixer of the propagation
  operator, \eq{eq:efix} \\
\verb|enable_reflection| & off & specular reflection of the advective fluxes
  at coastlines \\
\verb|limit_fdot| & off & clip the refraction velocities $\dot\nu, \dot\theta$
  to a CFL bound \\
\verb|refraction_cfl| & $0.5$ & the CFL bound used by \verb|limit_fdot| \\
\verb|closed_bottom| & off & extend the bottom cell of the coupling kernel to
  $z \to -\infty$ (exact barotropic projection, no momentum lost below the
  column floor) \\
\verb|tail_flux| & \verb|'wamdi'| & booking of the diagnostic-tail fluxes to
  the ocean: \verb|'wamdi'| (dissipation-rate continuation) or
  \verb|'input'| (equilibrium-range $\gamma_{in}$ integral plus nonlinear
  cascade credit) \\
\hline
\end{longtable}

\subsubsection*{Snapshot diagnostics}

The netCDF snapshot class reads the attribute \verb|snapint| (snapshot
interval in seconds, required) and takes the setup arguments
\verb|snap_file_name| (default \verb|'snap.cdf'|) and
\verb|show_each_time_step| (default on). The land mask is defined by
overloading the template method \verb|set_topography|; the default is an
all-ocean domain, with closed outer boundaries generated automatically from
\verb|periodic_in_x|/\verb|periodic_in_y|.


\end{document}

\clearpage\phantomsection\label{doc:wim_realistic}\addcontentsline{toc}{section}{wim\_realistic}
\documentclass[12pt]{article}
\usepackage{lineno}
%\usepackage{ametsoc}
\usepackage{natbib}
\usepackage{a4wide}
\usepackage{longtable}
\usepackage{amsfonts}
\usepackage{amssymb}

\usepackage{color}
\usepackage[pdftex]{graphicx}
\usepackage[pdftex,colorlinks, citecolor=black,linkcolor=black,urlcolor=black]{hyperref}
\definecolor{darkgreen}{rgb}{0.1,0.4,0.1}
\definecolor{darkblue}{rgb}{0.1,0.1,0.3}



\newif\ifdetail
\global\detailfalse
%\global\detailtrue

\def\bem#1{\colorbox{green}{\em #1}}

\newcommand{\beq}  { \begin{eqnarray} }
\newcommand{\eeq}  { \end{eqnarray}}
\newcommand{\beeq}  { \begin{eqnarray*} }
\newcommand{\eeeq}  { \end{eqnarray*}}
\newcommand{\eq }  [1]{Eq.~(\ref{#1})}  % Eq. (#1)
\newcommand{\eqs }  [1]{Eqs.~(\ref{#1})}
\newcommand{\sect}  [1]{Sec.~(\ref{#1})}  % Eq. (#1)
\newcommand{\fig}  [1]{Fig.~\ref{#1}}   % Fig. #1
\renewcommand{\v}[1]{{\mbox{\boldmath$ #1 $}}}  % Vektor durch Fettdruck
\newcommand{\m}    [1]{{\bf  #1 }}              % Matrix in sans serif
\newcommand{\gz}[1] {{{d #1} \over {dz}}}
\newcommand{\pt}[1]{{\partial_t #1}}
\newcommand{\ptt}[1]{{\partial_{tt} #1}}
\newcommand{\pttt}[1]{{\partial_{ttt} #1}}
\newcommand{\pttx}[1]{{{\partial^3 #1} \over {\partial t^2 \partial x}}}
\newcommand{\ptty}[1]{{{\partial^3 #1} \over {\partial t^2 \partial y}}}
\newcommand{\ptx}[1]{{{\partial^2 #1} \over {\partial t \partial x}}}
\newcommand{\pty}[1]{{{\partial^2 #1} \over {\partial t \partial y}}}
\newcommand{\pz}[1]{{\partial_z #1}}
\newcommand{\pzz}[1]{{\partial_{zz} #1}}
\newcommand{\px}[1]{{\partial_x #1}}
\newcommand{\py}[1]{{\partial_y #1}}
\newcommand{\pxi}[1]{{\partial_i #1}}
\newcommand{\pxj}[1]{{\partial_j #1}}
\newcommand{\pxk}[1]{{\partial_k #1}}
\newcommand{\pb}[1]{{{\partial #1} \over {\partial b}}}
%\newcommand{\pxy}[1]{{{\partial^2 #1} \over {\partial x \partial y}}}
\newcommand{\pxy}[1]{\partial_{xy} #1}
\newcommand{\pyy}[1]{\partial_{yy} #1}
\newcommand{\pxx}[1]{\partial_{xx} #1}
\newcommand{\kx}{{\bf k} \times}
\newcommand{\dt}[1]{{{d #1} \over {d t}}}
\newcommand{\Dt}[1]{{{D #1} \over {D t}}}
\newcommand{\Div}[1]{ \nabla \cdot #1 }
\renewcommand{\div}[1]{ \nabla_{yz} \cdot #1 }
\newcommand{\Bar}[1]{ \overline{ #1} }
\newcommand{\rvec} [1]{
    \raisebox{-1.5ex}{$\stackrel{\textstyle #1}{\neg}$} }% rotierter Vektor
\newcommand{\nablq}   {\rvec\nabla}                    % rotated nabla operator
\newcommand{\Vec}[1]{\left(\begin{array}{c} #1 \end{array}\right) }
%\renewcommand{\mho}{\mathcal{W}}
\newcommand{\pn}[1]{{{\partial #1} \over {\partial n}}}
\newcommand{\pnn}[1]{{{\partial #1} \over {\partial m}}}
\newcommand{\ps}[1]{{{\partial #1} \over {\partial s}}}
\newcommand{\p}{{\partial}}
\newcommand{\vn}{{\v \nabla}}
\newcommand{\lk}{\left(}
\newcommand{\rk}{\right)}

\newcommand{\tl}{\tilde}

\newcommand{\I}{\mathcal{I}}
\newcommand{\E}{\mathcal{E}}
\newcommand{\A}{\mathcal{A}}

\newcommand{\ep}{\epsilon}
\newcommand{\erfc}{\mathrm{erfc}}
 
 
\begin{document}

 \section*{The spectral wave model wim\_realistic}

wim\_realistic is wim\_simple on the grid of pe\_realistic.
This text describes only the differences to wim\_simple.
Everything that is not mentioned here is as described in the documentation of wim\_simple.

The differences are the following.
The horizontal grid spacings may vary.
The horizontal coordinates may be longitude and latitude on the sphere.
The propagation then contains the metric factors and the turning of the rays on the sphere.
The coupling to the ocean allows a vertical grid with varying cell thickness.
The wind input of Janssen and the dissipation of WAM4 are available as an option.

Most methods of wim\_realistic are those of wim\_simple.
They are inherited or assigned from wim\_simple.
wim\_simple itself is not changed by wim\_realistic.
The source terms of wim\_simple, the nonlinear transfer, the time stepping, the reflection, the open coasts and the JAX backend are therefore the same in both models.

\subsection*{Grid}

The model domain consists of $n_x \times n_y$ cells.
The horizontal coordinates $x$ and $y$ are either Cartesian coordinates in metres or longitude $\lambda$ and latitude $\phi$ in degrees.
The attribute \verb|coord_degree| selects longitude and latitude.
The spectral grid in frequency $\nu$ and direction $\theta$ is that of wim\_simple.

The widths of the cells are given by the attributes \verb|dxt| and \verb|dyt|.
Each of them is either a number or an array of length $n_x$ or $n_y$ without halo points.
The numbers \verb|dx| and \verb|dy| of wim\_simple are accepted instead.
With \verb|coord_degree| the widths are given in degrees.

The horizontal grid is built as in pe\_realistic.
A cell has its centre at a $T$ point and its faces at $U$ points.
In $x$ the cell $i$ has the centre $x^t_i$, the width $\Delta x^t_i$ and the eastern face $x^u_i$.
The distance between the centres $x^t_i$ and $x^t_{i+1}$ is $\Delta x^u_i$.
The same holds in $y$.
The faces and centres follow from the widths by the recursion of pyOM2,
\beq \label{eq:ucentred_w}
 x^u_i = x^u_{i-1} + \Delta x^t_i ~,~~
 x^t_{i+1} = 2 x^u_i - x^t_i ~,~~
 \Delta x^u_i = x^t_{i+1} - x^t_i
\eeq
The first interior $U$ point is placed at \verb|x_origin| and \verb|y_origin| (default 0).

The recursion \eq{eq:ucentred_w} can give $\Delta x^u \le 0$ when the widths jump.
This holds also across the seam of a periodic direction.
The widths must therefore vary smoothly.

With \verb|coord_degree| the widths are converted to metres after the coordinates are built.
The factor is $a \pi /180$ with the radius of the earth $a = 6370\,$km (attribute \verb|radius|).
The metric factors are $\cos \phi$ at the $T$ rows, $\cos \phi$ at the $U$ rows and $\tan \phi /a$ at the $T$ rows.
In Cartesian coordinates they are 1, 1 and 0.
The area of a cell is $A^t = \cos \phi^t \Delta x^t \Delta y^t$ with $\Delta x$ and $\Delta y$ in metres.

The attributes \verb|dx| and \verb|dy| are set to None after the grid is built.
Code from wim\_simple that still uses them therefore stops with an error.

The grid has no treatment of the poles and no reduced grid.
The domain must therefore end before the poles.
The zonal width $\cos \phi\, \Delta x^t$ of a cell becomes small at high latitude.
The Courant number in $x$ then grows with $1/\cos \phi$, and the time step has to be chosen accordingly.

\subsection*{Propagation}

On the sphere the action equation of wim\_simple becomes
\beq
 \pt A + \frac{1}{a \cos \phi} \lk \p_\lambda (\dot x A) + \p_\phi (\cos \phi\, \dot y A) \rk
 + \p_\theta (\dot \theta A) + \p_\nu (\dot \nu A) = S
\eeq
with $\lambda$ and $\phi$ in radians.
The velocities $\dot x$ and $\dot y$ are the eastward and northward velocities in m/s.
They are those of wim\_simple, $\dot{\v x} = (c/2) \v s + \v U + R \v s$.
The direction $\theta$ is counted from north, with $\v s = (\sin \theta, \cos \theta)$ as in wim\_simple.

The discrete form is the flux form of wim\_simple.
The flux in $y$ is multiplied by $\cos \phi^u$.
The flux divergence is divided by $\cos \phi^t \Delta x^t$ and $\cos \phi^t \Delta y^t$, with the widths in metres.
The fluxes in frequency and direction are those of wim\_simple.
The Courant numbers of the Superbee scheme are $|\dot x| \Delta t/(\cos \phi^t \Delta x^t)$ and $|\cos \phi^u \dot y| \Delta t /(\cos \phi^t \Delta y^t)$.

The refraction by the Doppler velocity $\v U$ needs the horizontal gradients of $\v U$.
On the sphere these are $(a \cos \phi)^{-1} \p_\lambda$ and $a^{-1} \p_\phi$.
The discrete differences of wim\_simple are therefore divided by $\cos \phi^t \Delta x^u$ and $\Delta y^u$ instead of $\Delta x$ and $\Delta y$.

On the sphere a ray that keeps its direction relative to north does not follow a great circle.
The direction therefore changes along the ray even without currents.
With currents the frequency changes as well.
Both effects add terms to $\dot \theta$ and $\dot \sigma$ of wim\_simple,
\beq
 \dot \theta &=& \dot \theta_{simple} + \lk c/2 + \v s \cdot \v U \rk \sin \theta \frac{\tan \phi}{a} \\
 \dot \sigma &=& \dot \sigma_{simple} - \frac{\sigma}{2} \lk \v U \cdot \v s_\perp \rk \sin \theta \frac{\tan \phi}{a}
\eeq
with $\v s_\perp = (\cos \theta, -\sin \theta)$ and $\dot \nu = \dot \sigma /(2 \pi)$.
Here $\dot \theta_{simple}$ and $\dot \sigma_{simple}$ are the expressions of wim\_simple with the gradients on the sphere.
The added term of $\dot \theta$ is evaluated at the faces in direction, as $\dot \theta$ of wim\_simple.
The added term of $\dot \sigma$ is evaluated at the centres in direction.

Without currents only the term $(c/2) \sin \theta \tan \phi /a$ remains.
It turns the ray onto the great circle.
The frequency is then constant along the ray.

The two terms follow exactly from the dispersion relation in spherical coordinates.
The derivation is checked symbolically in \verb|development/derive_spherical_rays.py|.
The documentation of ECMWF (IFS Part VII, Eq.\ 2.25c) keeps only the term $(c/2) \sin \theta \tan \phi/a$.
The terms with $\v U$ are of the order $|\v U| \tan \phi /a$.
This is about $2 \cdot 10^{-7}\,$s$^{-1}$ at $60^\circ$ latitude for $|\v U| = 1\,$m/s.

The option \verb|limit_fdot| of wim\_simple also bounds the added terms.

The energy fixer is that of wim\_simple.
The defect $D$ of wim\_simple is here the integral over the area of the cells.
The terms of the faces in $x$ are multiplied by $\Delta y^t$.
The terms of the faces in $y$ are multiplied by $\Delta x^t$, since the flux already carries $\cos \phi^u$.
The terms of the faces in frequency and direction are multiplied by $A^t$.
On a uniform Cartesian grid this reduces to the fixer of wim\_simple.

The open coasts are those of wim\_simple.
The outflow through a face in $x$ is weighted with $\Delta y^t$ and through a face in $y$ with $\Delta x^t$.

The Fortran kernels of the propagation (\verb|propagation_backend| $=$ 'fortran') are those of wim\_simple.
The metric factors and the added terms above are handed to them.
They give the same result as the array form.

\subsection*{Coupling to an ocean with varying cell thickness}

The coupling of wim\_simple assumes ocean cells of constant thickness $\Delta z$.
pe\_realistic has cells of varying thickness $\Delta z^t$.
wim\_realistic therefore accepts $\Delta z^t$ as an array in the coupling.

The Stokes velocity, the stress and the Doppler velocity of wim\_simple use the kernel $Q = 2 k e^{2 k z}$ integrated over each ocean cell,
\beeq
 \int_{z_b}^{z_t} Q\, dz = e^{2 k z_t} - e^{2 k z_b}
\eeeq
with the top face $z_t$ and the bottom face $z_b$ of the cell.
For an array $\Delta z^t$ the faces are the partial sums of $\Delta z^t$ from the surface downward.
The array is given from the bottom to the surface, as in pe\_realistic.
The top face of a cell is then exactly the bottom face of the cell above.
The sum of the kernel over the column is therefore exactly $1 - e^{-2 k H}$ for a column of depth $H$.
With \verb|closed_bottom| the lowest cell also takes the rest $e^{-2 k H}$, so that the sum is exactly 1, as in wim\_simple.

\verb|calc_stokes| and \verb|calc_stress_work| divide the integrated kernel by the thickness $\Delta z^t$ of each cell.
The Doppler velocities need no division and are those of wim\_simple.
With a number $\Delta z$ the arithmetic is that of wim\_simple, bit for bit.

\verb|calc_phioc_work| of wim\_simple is not generalised.
It works only with a number $\Delta z$.

\subsection*{Wind input of Janssen and dissipation of WAM4}

The attribute \verb|input_scheme| selects the source terms.
The default 'wamdi' uses the input and dissipation of wim\_simple.
The value 'janssen' uses the wind input of Janssen (1991) and the dissipation of WAM4 (Bidlot et al.\ 2007).
These are the source terms of ecWAM for \verb|IPHYS| $=0$.
They are ported from the routines \verb|sinput_jan|, \verb|sdissip_jan|, \verb|stresso|, \verb|tau_phi_hf|, \verb|z0wave|, \verb|fkmean|, \verb|femeanws| and \verb|frcutindex| of ecWAM.
The nonlinear transfer is that of wim\_simple in both cases.

The friction velocity $u_*$ and its direction $\theta_u$ are prescribed, as in wim\_simple.
ecWAM computes $u_*$ from the wind speed and the roughness.
This part of ecWAM is not ported.

The input is $S_{in} = \gamma_{in} E$ with the growth rate
\beq
 \gamma_{in} = \frac{\beta_{max}}{\kappa^2} \frac{\rho_a}{\rho_w}\, \sigma\, x^2 \mu \lk \ln \mu \rk^4
\eeq
Here $x = (u_*/c + z_\alpha) \cos(\theta - \theta_u)$ and $\ln \mu = \ln(k z_0) + \kappa / x$.
$\kappa = 0.4$ is the constant of von K\'arm\'an.
The input is zero where $\cos(\theta - \theta_u) \le 0.01$ or $\ln \mu \ge 0$.
The cells with input define the wind sea used below.

The roughness length depends on the momentum that the waves take from the wind,
\beq
 z_0 = \frac{\alpha}{g} \frac{u_*^3}{\sqrt{u_*^2 - \tau_w}}
\eeq
with the constant $\alpha$ of Charnock.
The momentum flux $\tau_w$ into the waves is divided by the density of air $\rho_a$.
It is the input of momentum over all frequency bins plus the analytic tail described below,
\beeq
 \tau_w = \frac{\rho_w}{\rho_a} \left| \int \gamma_{in} E\, \sigma\, \v s\, d\theta\, d\nu + \v \tau_{tail} \right|
\eeeq
It is limited to $u_*^2/(1+\epsilon_1)$.
Without waves $\tau_w = 0$ and $z_0 = \alpha u_*^2/g$ is the roughness of Charnock.

$z_0$ and $\tau_w$ depend on each other.
In each source step $z_0$, $\gamma_{in}$ and $\tau_w$ are therefore computed \verb|janssen_ncall| times in turn, as ecWAM does.
The first pass starts from $\tau_w$ of the previous step.

The dissipation is $S_{out} = \gamma_d E$ with
\beq
 \gamma_d = c_{dis}\, 2 \pi \nu_1\, E_{tot}^2\, \bar k^4\, X \lk (1-\delta) + \delta X \rk ~,~~ X = k/\bar k
\eeq
It does not depend on direction.
$E_{tot}$ is the total energy, $\nu_1$ the mean frequency weighted with $\nu$ and $\bar k$ the mean wavenumber from the moment with $\sqrt k$.

The means are sums over the frequency bins plus an analytic tail $E_N (\nu_N/\nu)^5$ beyond the last bin $\nu_N$,
\beeq
 E_{tot} &=& \sum E_d \Delta \nu + \nu_N E_N /4 \\
 \nu_1 &=& \lk \sum \nu E_d \Delta \nu + \nu_N^2 E_N/3 \rk / E_{tot} \\
 \bar k &=& \lk \lk \sum \sqrt{k}\, E_d \Delta \nu + \frac{2 \pi}{\sqrt g} \nu_N^2 E_N/3 \rk / E_{tot} \rk^2
\eeeq
with $E_d = \int E\, d\theta$ and $E_N$ of the last bin.

The prognostic range ends at $2.5 \max(\bar \nu, \bar \nu_{ws})$, at the nearest bin in $\ln \nu$.
Here $\bar \nu$ is the mean frequency weighted with $1/\nu$ and $\bar \nu_{ws}$ the same mean over the wind sea.
The lowest frequency bin is always prognostic.
The term of Pierson and Moskowitz in the cut-off of wim\_simple is not used, as in ecWAM.

Beyond the prognostic range the spectrum is $E_c (\nu_c/\nu)^5$, as in wim\_simple.
The stress and the energy flux of this tail are computed as in \verb|tau_phi_hf| of ecWAM.
The integral over the tail uses the rule of Simpson with \verb|jtot_tauhf| $=19$ points in $\ln(\omega \sqrt{z_0/g})$.
It runs from $\max(\omega_c, x_0 g/u_*)$ to $\omega = \sqrt{g/z_0}$.
$x_0$ is the value of $u_*/c$ below which the growth rate vanishes for the roughness of Charnock.
The growth in the tail is the one along the wind times $\cos^2(\theta - \theta_u)$ for $\cos(\theta - \theta_u) > 0$.

The tail of $\tau_w$ starts at the last frequency bin of the model.
The frequency ghost bin is still zero when the input is computed.
The tails of the stress and of $\Phi_{oc}$ with \verb|tail_flux| $=$ 'input' start at the ghost bin, as in wim\_simple.
\verb|calc_stress_work| puts the tail stress into the ocean cells point by point of the rule of Simpson, with the kernel $2 k_j e^{2 k_j z}$ of each point $j$.

The port differs from ecWAM in three points.
The widths of the first and last frequency bins are those of wim\_simple, while ecWAM uses half bins there.
This changes the means by $0.05$ to $0.5\,\%$.
The tail of $\tau_w$ starts at the last model bin, as stated above.
The tail stress is booked along each direction $\v s(\theta)$.
Its component along the wind is the stress of ecWAM, and ecWAM drops the component across the wind.

The following parts of ecWAM are not ported.
These are the gustiness, the negative input, the normalisation of the growth rate, the roughness of gravity-capillary waves, the sheltering and the viscous dissipation.
The viscous term is in the code with $c_{disvis} = 0$, as in ecWAM for \verb|IPHYS| $=0$.

$\tau_w$ and $z_0$ are carried from step to step.
They are not part of the restart of the model.
A restart starts with $\tau_w = 0$, i.e.\ with the roughness of Charnock.
The iteration of $z_0$ and $\tau_w$ corrects this within the first step.

The port is compared with the routines of ecWAM compiled with f2py.
The test is \verb|janssen_ref_test.py| in \verb|development/|.
All quantities agree to round-off.

\section*{Model parameters and options}

The parameters and options of wim\_simple are also those of wim\_realistic.
The tables list only the differences.

\begin{longtable}{@{}llp{7.6cm}@{}}
\hline
Attribute & Default & Meaning \\
\hline
\endhead
\verb|dxt|, \verb|dyt| & required & widths of the cells, a number or an array without halo.
  Degrees with \verb|coord_degree|, metres otherwise.
  Replace \verb|dx| and \verb|dy| of wim\_simple, which are still accepted \\
\verb|coord_degree| & off & longitude and latitude in degrees instead of Cartesian coordinates \\
\verb|x_origin|, \verb|y_origin| & $0$ & position of the first interior $U$ point \\
\verb|radius| & $6370\,$km & radius of the earth (m) \\
\verb|input_scheme| & 'wamdi' & source terms. 'wamdi' as in wim\_simple, 'janssen' the input of Janssen and the dissipation of WAM4 \\
\hline
\end{longtable}

\noindent
With \verb|input_scheme| $=$ 'janssen' the following constants are read.
Their defaults are those of ecWAM for \verb|IPHYS| $=0$.

\begin{longtable}{@{}llp{8.6cm}@{}}
\hline
Attribute & Default & Meaning \\
\hline
\endhead
\verb|betamax| & $1.20$ & $\beta_{max}$ of the input \\
\verb|zalp| & $0.008$ & $z_\alpha$ of the input \\
\verb|alpha_charnock| & $0.0065$ & constant $\alpha$ of Charnock \\
\verb|kappa| & $0.40$ & constant of von K\'arm\'an \\
\verb|cdis| & $-1.33$ & $c_{dis}$ of the dissipation \\
\verb|delta_sdis| & $0.5$ & $\delta$ of the dissipation \\
\verb|cdisvis| & $0$ & viscous dissipation \\
\verb|rnu| & $1.5 \cdot 10^{-6}$ & kinematic viscosity (m$^2$/s) of the viscous dissipation \\
\verb|eps1| & $10^{-5}$ & $\epsilon_1$ in the limit of $\tau_w$ \\
\verb|tailfactor| & $2.5$ & factor of the prognostic cut-off \\
\verb|janssen_ncall| & $2$ & number of passes of $z_0$, $\gamma_{in}$ and $\tau_w$ per step \\
\verb|jtot_tauhf| & $19$ & number of points of the tail integral \\
\verb|rho_a| & $1.225$ & density of air (kg/m$^3$) \\
\verb|rho_w| & $1000$ & density of water (kg/m$^3$) \\
\hline
\end{longtable}

\noindent
The densities \verb|rho_a| and \verb|rho_w| act only in the source terms of Janssen.
The input of wim\_simple and the tails of the coupling use $1.225/1000$.

\subsection*{Snapshot diagnostics}

The snapshot classes are those of wim\_simple.
They write the coordinates $x^t$ and $y^t$ of the grid, in degrees with \verb|coord_degree|.
The fields are those of wim\_simple.

\subsection*{Tests}

The tests are in \verb|pyOM3/wim_realistic/development/|.
\verb|reduction_test.py| compares wim\_realistic on a uniform Cartesian grid with wim\_simple.
The action agrees bit for bit, and to $10^{-15}$ with the energy fixer.
\verb|sphere_test.py| tests the propagation on the sphere.
The total action is conserved to $10^{-16}$ over three days.
A wave packet follows the great circle to $0.15^\circ$ in position and $0.6^\circ$ in direction after $36^\circ$ of arc.
The error of a uniform state decreases with the second power of the grid spacing.
The absolute energy closes to $2 \cdot 10^{-16}$ with a sheared current.
\verb|janssen_ref_test.py| compares the source terms of Janssen with ecWAM, see above.

\subsection*{Limits}

\begin{itemize}
\item The domain must end before the poles.
\item The widths of the cells must vary smoothly, also across a periodic seam.
\item The terms with $\v U$ in the turning on the sphere are not tested by a dynamic test.
\item \verb|calc_phioc_work| works only with a number $\Delta z$.
\item With \verb|input_scheme| $=$ 'janssen' and \verb|tail_flux| $=$ 'wamdi' the tails of the stress use the dissipation of wim\_simple.
\item $\tau_w$ and $z_0$ of the scheme of Janssen are not part of the restart.
\end{itemize}

\end{document}

\clearpage\phantomsection\label{doc:wim_snl_exact}\addcontentsline{toc}{section}{The exact $S_{nl}$ of wim\_simple}
\documentclass[12pt]{article}
\usepackage{lineno}
%\usepackage{ametsoc}
\usepackage{natbib}
\usepackage{a4wide}
\usepackage{longtable}
\usepackage{amsfonts}
\usepackage{amssymb}

\usepackage{color}
\usepackage[pdftex]{graphicx}
\usepackage[pdftex,colorlinks, citecolor=black,linkcolor=black,urlcolor=black]{hyperref}
\definecolor{darkgreen}{rgb}{0.1,0.4,0.1}
\definecolor{darkblue}{rgb}{0.1,0.1,0.3}



\newif\ifdetail
\global\detailfalse
%\global\detailtrue

\def\bem#1{\colorbox{green}{\em #1}}

\newcommand{\beq}  { \begin{eqnarray} }
\newcommand{\eeq}  { \end{eqnarray}}
\newcommand{\beeq}  { \begin{eqnarray*} }
\newcommand{\eeeq}  { \end{eqnarray*}}
\newcommand{\eq }  [1]{Eq.~(\ref{#1})}  % Eq. (#1)
\newcommand{\eqs }  [1]{Eqs.~(\ref{#1})}
\newcommand{\sect}  [1]{Sec.~(\ref{#1})}  % Eq. (#1)
\newcommand{\fig}  [1]{Fig.~\ref{#1}}   % Fig. #1
\renewcommand{\v}[1]{{\mbox{\boldmath$ #1 $}}}  % Vektor durch Fettdruck
\newcommand{\m}    [1]{{\bf  #1 }}              % Matrix in sans serif
\newcommand{\gz}[1] {{{d #1} \over {dz}}}
\newcommand{\pt}[1]{{\partial_t #1}}
\newcommand{\ptt}[1]{{\partial_{tt} #1}}
\newcommand{\pttt}[1]{{\partial_{ttt} #1}}
\newcommand{\pttx}[1]{{{\partial^3 #1} \over {\partial t^2 \partial x}}}
\newcommand{\ptty}[1]{{{\partial^3 #1} \over {\partial t^2 \partial y}}}
\newcommand{\ptx}[1]{{{\partial^2 #1} \over {\partial t \partial x}}}
\newcommand{\pty}[1]{{{\partial^2 #1} \over {\partial t \partial y}}}
\newcommand{\pz}[1]{{\partial_z #1}}
\newcommand{\pzz}[1]{{\partial_{zz} #1}}
\newcommand{\px}[1]{{\partial_x #1}}
\newcommand{\py}[1]{{\partial_y #1}}
\newcommand{\pxi}[1]{{\partial_i #1}}
\newcommand{\pxj}[1]{{\partial_j #1}}
\newcommand{\pxk}[1]{{\partial_k #1}}
\newcommand{\pb}[1]{{{\partial #1} \over {\partial b}}}
%\newcommand{\pxy}[1]{{{\partial^2 #1} \over {\partial x \partial y}}}
\newcommand{\pxy}[1]{\partial_{xy} #1}
\newcommand{\pyy}[1]{\partial_{yy} #1}
\newcommand{\pxx}[1]{\partial_{xx} #1}
\newcommand{\kx}{{\bf k} \times}
\newcommand{\dt}[1]{{{d #1} \over {d t}}}
\newcommand{\Dt}[1]{{{D #1} \over {D t}}}
\newcommand{\Div}[1]{ \nabla \cdot #1 }
\renewcommand{\div}[1]{ \nabla_{yz} \cdot #1 }
\newcommand{\Bar}[1]{ \overline{ #1} }
\newcommand{\rvec} [1]{
    \raisebox{-1.5ex}{$\stackrel{\textstyle #1}{\neg}$} }% rotierter Vektor
\newcommand{\nablq}   {\rvec\nabla}                    % rotated nabla operator
\newcommand{\Vec}[1]{\left(\begin{array}{c} #1 \end{array}\right) }
%\renewcommand{\mho}{\mathcal{W}}
\newcommand{\pn}[1]{{{\partial #1} \over {\partial n}}}
\newcommand{\pnn}[1]{{{\partial #1} \over {\partial m}}}
\newcommand{\ps}[1]{{{\partial #1} \over {\partial s}}}
\newcommand{\p}{{\partial}}
\newcommand{\vn}{{\v \nabla}}
\newcommand{\lk}{\left(}
\newcommand{\rk}{\right)}

\newcommand{\tl}{\tilde}

\newcommand{\I}{\mathcal{I}}
\newcommand{\E}{\mathcal{E}}
\newcommand{\A}{\mathcal{A}}

\newcommand{\ep}{\epsilon}
\newcommand{\erfc}{\mathrm{erfc}}
 

\begin{document}

\section*{The exact scattering integral $S_{nl}$ in wim\_simple: \\ formulation, discretisation and validation}

This note documents the exact evaluation of the non-linear transfer $S_{nl}$
implemented in the sub-package snl\_exact of wim\_simple
(''snonlin\_backend''$\,=\,$'exact').
It complements the model documentation, which contains only a summary,
with the underlying conventions, the derivation chain of the coupling
coefficient, the geometry of the resonance loci, the discretisation on the
spectral grid, the discrete conservation properties, and the validation
against independent implementations.
Only deep water is considered.

\subsection*{Conventions and the kinetic equation}

The kinetic equation below is written for the wave action spectral density
$A(\v k)$ per $d^2 k$. It is related to the surface elevation spectrum
$F(\v k)$, with $\langle \zeta^2 \rangle = \int F \, d\v k$, by
\beq \label{eq:convert}
 A \;=\; \frac{g}{\sigma}\, F   ~~,~~ \sigma(\v k) = \sqrt{g |\v k|}
\eeq
The model carries its spectra as functions of cyclic frequency $\nu$ and
direction $\theta$. Its spectral variables are the variance density
$E(\nu,\theta)$ per Hz and radian and the action $E/\sigma$, which is also a
density per Hz and radian. This model action is not the density $A$ per
$d^2 k$ of \eq{eq:convert}. The model hands $E$ to the exact backend.
With $k = (2\pi\nu)^2/g$ and $dk/d\nu = 8\pi^2 \nu / g$ the backend converts
$E$ to the action density per $d^2 k$ by
\beq
 A(k,\theta) \;=\; \frac{g \, E(\nu,\theta)}{\sigma \, k \, dk/d\nu}
\eeq
and converts the computed transfer back to $E$ with the same factor.

The evolution of $A$ by resonant four-wave interactions is the kinetic
equation of Hasselmann (1962), written here in the Zakharov form
\beq \label{eq:kinetic}
 S_{nl}(\v k_1) = 4\pi \int T^2_{1234}\,
 \delta(\v k_1 + \v k_2 - \v k_3 - \v k_4)\,
 \delta(\sigma_1 + \sigma_2 - \sigma_3 - \sigma_4)\; \mathcal{B}\;
 d\v k_2 \, d\v k_3 \, d\v k_4
\eeq
with the density bracket
\beq \label{eq:bracket}
 \mathcal{B} \;=\; A_3 A_4 (A_1 + A_2) - A_1 A_2 (A_3 + A_4)
\eeq
Two properties of $\mathcal{B}$ are used repeatedly below: it vanishes
identically for constant $A$, and it vanishes on the resonance manifold for
the Rayleigh--Jeans equilibrium $A \propto 1/\sigma$, since
\beq \label{eq:rj}
 \mathcal{B}\big|_{A = c/\sigma} \;=\;
 c^3\, \frac{(\sigma_1+\sigma_2) - (\sigma_3+\sigma_4)}
           {\sigma_1 \sigma_2 \sigma_3 \sigma_4} \;=\; 0
\eeq
Both properties are preserved by the discretisation (see below) and provide
sharp correctness tests.

The kernel $T_{1234}$ in \eq{eq:kinetic} is the canonical four-wave kernel of
the Zakharov equation
\beq \label{eq:zakharov}
 i\, \p_t b_1 \;=\; \int T_{1234}\, b_2^* b_3 b_4\,
 \delta(\v k_1 + \v k_2 - \v k_3 - \v k_4)\, d\v k_2\, d\v k_3\, d\v k_4
\eeq
in the Fourier convention without $2\pi$ prefactors, for which
$\langle b(\v k) b^*(\v k') \rangle = A(\v k)\,\delta(\v k - \v k')$ holds
with the physical action density \eq{eq:convert} directly, and
$T(\v k,\v k,\v k,\v k) = k^3$.
In the convention of Krasitskii (1994), with a factor $1/2\pi$ per transform,
the same kernel reads $T^{K} = T/(2\pi)^2$ and
$T^{K}(\v k,\v k,\v k,\v k) = k^3/4\pi^2$; the collision prefactor $4\pi$
in \eq{eq:kinetic} is form-invariant under this change of convention.
The combination
\beq \label{eq:webb}
 G \;=\; 4\pi\, T^2 \;=\; 4\pi (2\pi)^4 \lk T^{K} \rk^2
\eeq
is the deep-water coupling coefficient of Webb (1978).

\subsection*{Derivation and verification of the kernel}

The kernel is not transcribed from the literature but derived from the
water-wave Hamiltonian, truncated at fourth order in wave steepness,
\beq
 H_3 &=& \frac{1}{2} \int \zeta \lk (\vn \psi)^2 - (G_0 \psi)^2 \rk d\v x \\
 H_4 &=& \frac{1}{2} \int \zeta \, (G_0 \psi)
         \lk G_0 (\zeta\, G_0 \psi) + \zeta \, \nabla^2 \psi \rk d\v x
\eeq
with the surface potential $\psi$, the flat-surface Dirichlet--Neumann
operator $G_0$ (symbol $|\v k|$ in deep water), and normal variables defined by
\beq
 \hat\zeta(\v k) = \sqrt{\frac{k}{2\sigma}} \lk a_{\v k} + a^*_{-\v k} \rk
 ~~,~~
 \hat\psi(\v k) = -i \sqrt{\frac{\sigma}{2 k}} \lk a_{\v k} - a^*_{-\v k} \rk
\eeq
(here written in the $2\pi$-free convention).
Both truncated Hamiltonians were verified exactly against the analytic
kinetic energy of a prescribed harmonic flow, with residuals scaling as the
neglected orders $\varepsilon^3$, $\varepsilon^4$, $\varepsilon^5$.
This check was made during the development. It is not part of the tests
shipped with the sub-package (test\_kernel.py, test\_locus.py,
test\_assembly.py and test\_physics.py).
The canonical (normal-form) transformation then removes the non-resonant
cubic terms; on the four-wave resonance manifold the two classes of induced
quartic contributions coincide pairwise, which is exploited by keeping both
with weight $1/4$ so that the $(12)\!\leftrightarrow\!(34)$ symmetry of $T$
is manifest also off the manifold.

The resulting kernel satisfies, at machine precision, the on-shell
Hermiticity and index symmetries, rotation invariance, cubic homogeneity
$T(\lambda \v k_i) = \lambda^3 T(\v k_i)$, and independence of $g$.
It reproduces the two-wave degenerate-quartet interaction coefficients of
Longuet-Higgins/Leblanc (as compiled by Stuhlmeier, 2024) to about
$10^{-6}$.
These coefficients are values of $T$ at degenerate quadruplets such as
$(\v k, \v k, \v k, \v k)$. There the expression for $T$ contains vanishing
difference wavenumbers and cannot be evaluated directly.
The tests therefore evaluate $T$ at nearby quadruplets
$(\v k_i, \v k_j, \v k_i + \v d, \v k_j - \v d)$ with a small offset
$\v d$ of length $10^{-5}$ or $10^{-6}$ times a fixed vector of order one.
The agreement of about $10^{-6}$ is set by this offset, not by the kernel.
The deviation decreases linearly with the offset and is
$1.2 \cdot 10^{-9}$ for $(\v k,\v k,\v k,\v k)$ at an offset factor of
$10^{-8}$.
During the development the kernel also reproduced the Stokes
second-harmonic and mean-flow response and the bichromatic super- and
sub-harmonic amplitudes. These checks are not part of the shipped tests.
Finally, the identity
\beq \label{eq:xcwebb}
  \mathrm{xc\_webb} \;=\; 64 \pi^5 g^2 \lk T^{K}\rk^2 \;=\; g^2 \, G
\eeq
connecting \eq{eq:webb} to the coupling routine of the exact-NL code of
van Vledder (mod\_xnl4v5, as distributed with WAVEWATCH III and SWAN) was
verified to a relative accuracy of $3 \cdot 10^{-13}$ on resonant quadruplets,
so that the absolute normalisation of \eq{eq:kinetic} is anchored to an
independent production code at the level of the coupling coefficient.

\subsection*{WRT reduction and the resonance loci}

Following Webb (1978) and Tracy and Resio (1982), the wavenumber delta
function in \eq{eq:kinetic} eliminates $\v k_2$, and for fixed
$(\v k_1, \v k_3)$ the frequency delta function reduces the remaining
$\v k_2$-plane integral to a line integral along the resonance locus
\beq \label{eq:locus}
 z(\v k_2) \;=\; q + \sigma(\v k_2) - \sigma(\v k_2 + \v P) \;=\; 0
 ~~,~~ \v P = \v k_1 - \v k_3 ~,~ q = \sigma_1 - \sigma_3
\eeq
with $\v k_4 = \v k_2 + \v P$, so that
\beq \label{eq:wrt}
 S_{nl}(\v k_1) \;=\; \int d\v k_3 \oint_{z=0}
 \frac{ds}{|\vn_{\v k_2} z|}\; G \, \mathcal{B}
\eeq
Note that $z(\v k_3) = 0$ identically: the trivial quadruplet
$\v k_2 = \v k_3$, $\v k_4 = \v k_1$ lies on every locus, with bounded $G$;
the vanishing of the integrand there resides entirely in $\mathcal{B}$.

In deep water the loci form a one-parameter family: scaling lengths by
$|\v P|$ and frequencies by $\sqrt{g |\v P|}$, the locus shape depends only on
\beq
 \hat q = \frac{q}{\sqrt{g |\v P|}} \in (0,1)
\eeq
(the labels $1,3$ are chosen such that $q>0$; $\hat q < 1$ holds for every
pair since $(\sqrt{k_1}-\sqrt{k_3})^2 < k_1 - k_3 \le |\v k_1 - \v k_3|$).
In the normalised frame with $\hat{\v P} = (1,0)$ the two intersections of
the locus with the symmetry axis are analytic,
\beq \label{eq:crossings}
 x_+ = \lk \frac{1-\hat q^2}{2 \hat q} \rk^2
 ~~,~~
 x_- = - \lk \frac{\sqrt{2-\hat q^2} - \hat q}{2} \rk^2
\eeq
which provide exact anchors for the tracing.
The code contains two tracers. Both describe the closed curve by its radius
$r(\theta)$ about the interior midpoint of \eq{eq:crossings}. Here $\theta$
is the polar angle about this midpoint, not the wave direction.
Both compute the derivative of the radius along the curve analytically from
the gradient,
\beq
 r'(\theta) = - \frac{\p z / \p \theta}{\p z / \p r}
 ~~,~~ ds = \sqrt{r^2 + r'^2}\, d\theta
\eeq
The operator uses the routine trace\_batch of assembly.py.
It places the $n_{\rm pts}$ points of a locus uniformly in $\theta$, with
spacing $2\pi/n_{\rm pts}$, and finds the radius of each point by bisection
along its ray.
The periodic trapezoid rule in $\theta$ then yields spectrally accurate
quadrature weights $w_s = (ds/d\theta)\,(2\pi/n_{\rm pts}) / |\vn z|$.
The routine trace\_locus of locus.py is used in the locus tests.
It traces the curve densely by radial continuation from ray to ray and
resamples it to equal arc length $\Delta s$, with the weights
$w_s = \Delta s / |\vn z|$.
For $q = 0$ (equal frequencies, $\v k_1 \ne \v k_3$) the locus degenerates to
the perpendicular bisector $\{ \v k_2 : \v k_2 \cdot \v P = -|\v P|^2/2 \}$;
these same-frequency interactions are unbounded in extent but their line
integral converges through the decay of the spectrum, and they are included
with a separate straight-line quadrature.
The dimensional line measure follows from the scale invariance,
\beq \label{eq:measure}
 \frac{ds}{|\vn z|} \;=\; \frac{|\v P|^{3/2}}{\sqrt{g}}\;
 \frac{d\hat s}{|\hat\vn \hat z|}
\eeq
and the coupling scales as $G \propto |\v k|^6$ by the homogeneity of $T$.
Combining the scalings of \eq{eq:wrt}, the transfer obeys the deep-water
self-similarity
\beq \label{eq:selfsim}
 S_{nl}\lk \v k; \lambda A(\cdot/\mu) \rk \;=\;
 \lambda^3 \mu^{19/2}\, S_{nl}\lk \v k/\mu;  A \rk
\eeq
which the discretisation reproduces on the geometric grid for the closed
loci in the interior of the grid (see below).

\subsection*{Discretisation}

The spectral grid of wim\_simple is geometric in frequency and uniform in
direction, hence geometric in $k$ with uniform spacing $h$ in $\ln k$.
Two consequences are exploited.
First, a rotation of the pair $(\v k_1, \v k_3)$ by a multiple of the
direction spacing $\Delta\theta$ rotates its locus by the same angle.
The locus geometry, the coupling values, the quadrature weights and the
interpolation stencils are therefore tabulated for $\theta_1 = 0$ only.
The table covers every pair $(i_1, i_3)$ of frequency indices and every
direction-index difference $j_3 - j_1$.
Each of these configurations is traced once, exactly.
In the stencils the $\ln k$ indices are absolute grid indices, and the
direction indices are offsets relative to $j_1$.
The table is applied to every $\theta_1$ by shifting all direction indices
by $j_1$.
This renders the discrete operator exactly equivariant under rotations of
the spectrum by grid directions.
The size of the table scales as $n_f^2\, n_d\, n_{\rm pts}$, with $n_f$
frequencies, $n_d$ directions and $n_{\rm pts}$ points per locus.
There is no interpolation in $\hat q$; the discrete grid realises finitely
many configurations, and all of them are solved.
Second, $\hat q$ depends only on the ratio $k_3/k_1$ and the direction
difference.
Let $\rho$ be the ratio of neighbouring frequencies, so that $h = 2\ln\rho$.
A shift of $i_1$ and $i_3$ by one ring multiplies all wavenumbers of the
quadruplets by $\rho^2$, and the stencils move by one index in $\ln k$.
In the interior of the grid this renders the closed-locus part of the
operator equivariant under shifts along the geometric $k$-axis
--- the discrete form of \eq{eq:selfsim}.
Near the grid boundaries this equivariance is broken.
There the strict truncation described below discards different quadruplets
for a spectrum and for its shifted copy.
The same-frequency line loci break it as well.
Each line is cut at $|\v k_2| = 1.02\, k_{max}$, with $k_{max}$ the largest
wavenumber of the grid, and not at a fixed multiple of $|\v P|$.
The shipped test of \eq{eq:selfsim} therefore leaves out the line loci and
evaluates only interior rings.

Action densities at the off-grid locus points $\v k_2, \v k_4$ are obtained
by bilinear interpolation in $(\ln k, \theta)$.
Locus points within a small fraction of $|\v P|$ of the trivial point
$\v k_2 = \v k_3$ are removed: there the exact bracket \eq{eq:bracket}
vanishes but its interpolant does not; the computed transfer is insensitive
to the filter width at the percent level, which is verified as part of the
test battery.
Above the discrete frequency range the spectrum is extended by a parametric
$E \propto \nu^{-5}$ tail over a number of additional geometric rings before
evaluation (following WRT practice in WAVEWATCH III), and the transfer is
restricted to the original band afterwards; quadruplets with any member
beyond the extended range are discarded entirely (strict truncation).
Strict truncation is what preserves the equilibria: every retained quadruplet
carries all four action values, so $\mathcal{B} \equiv 0$ termwise for
constant $A$, and \eq{eq:rj} holds up to interpolation error for
$A \propto 1/\sigma$.
This holds for the operator itself, evaluated without the tail.
The backend in wim\_simple appends the tail by default, with eight rings
(pad\_hi$\,=\,$8, see below).
The tail continues an equilibrium spectrum of the band by
$E \propto \nu^{-5}$, so the extended spectrum is no equilibrium.
The backend therefore returns a nonzero transfer for an equilibrium
spectrum of the band.
For constant $A$ on a band of 20 rings this transfer reaches about $0.2$ of
the one-sided (gross) rate, also in the interior of the band.

With cell areas $\Delta_i = k_i\, \Delta k_i \, \Delta\theta$, where
$\Delta k_i = k_i \lk \rho - 1/\rho \rk = 2 k_i \sinh(h/2)$ is the width of
ring $i$, the discrete form of
\eq{eq:wrt} for one quadruplet instance $(i_1 j_1, i_3 j_3, s)$ carries the
measure
\beq
 \mu_q \;=\; \Delta_{3}\; w_s \; G_s
\eeq
and two assemblies of the same integral are provided.
In the \emph{direct} mode the rate $\mu_q \mathcal{B}$ is accumulated at
$\v k_1$ only --- the discrete analogue of \eq{eq:wrt} and of other WRT
implementations.
In the \emph{conservative} mode each instance deposits the total-action
increments
\beq \label{eq:scatter}
 \lk \delta \mathcal{A}_1, \delta \mathcal{A}_2,
     \delta \mathcal{A}_3, \delta \mathcal{A}_4 \rk
 \;=\; \lk +1,+1,-1,-1 \rk \, \frac{\Delta_1 \mu_q \mathcal{B}}{4}
\eeq
distributed over the grid through the same bilinear stencils used for the
gather.
Dividing the deposit at the interpolated legs by the effective cell area
$\bar\Delta$, interpolated in $\ln k$ with the same weights, conserves the
deposited action exactly, since the stencil weights $w_t$ and the
$\theta$-independent areas satisfy $\sum_t w_t \Delta_{c(t)} = \bar\Delta$
identically.
The factor $1/4$ in \eq{eq:scatter} averages the four role-estimators of the
enumeration over ordered pairs; the total action budget of the discrete
transfer, including the action exchanged with the parametric tail, then
closes at machine precision by construction.
The two modes are independent discretisations of \eq{eq:wrt}; their
difference is $O(h)$, decreases under grid refinement, and serves as a
built-in convergence diagnostic.
Energy and momentum are conserved by the continuous quadruplets and hence to
interpolation accuracy by the discrete transfer.
The code does not compute their residuals.
The backend stores action diagnostics of the last evaluated cell in the
dictionary Snonlin.diag.
In the conservative mode these are action\_rate, the sum of $\Delta_i$
times the transfer over the extended grid, and leak, the action of deposits
that fall outside the grid.
In both modes they are band\_action\_rate and tail\_action\_rate, the same
sum over the band and over the tail rings.

For the semi-implicit time stepping a diagonal derivative is accumulated
alongside,
\beq
 \frac{\p S_{nl}(\v k_1)}{\p A_1} \;=\;
 \sum_q \mu_q \lk A_3 A_4 - A_2 (A_3 + A_4) \rk
\eeq
It is the derivative of the direct rate at $\v k_1$ with respect to $A_1$,
with $A_2$, $A_3$ and $A_4$ held fixed.
The code returns it in both modes.
In the conservative mode, the default, it is not the derivative of the
returned transfer.
There $A_1$ enters also through the deposits of the instances in which
$\v k_1$ plays the role 2, 3 or 4, and the role-1 deposit carries the
factor $1/4$.
At the peak cell of a JONSWAP spectrum and at a cell above the peak the
derivative of the conservative transfer, estimated by finite differences,
differs from the sum above by about a factor 2.
In the direct mode the two differ by about 15\% at these cells, since
$A_1$ also enters through the interpolated $A_2$ and $A_4$.
The sum serves as the diagonal factor of the semi-implicit step.
It is invariant under the conversion between $A$ and $E$ on the diagonal
and is used directly in the update
$E \leftarrow E + S\, \Delta t / (1 - \Delta t\, \p S/\p E)$.

\subsection*{The backend in wim\_simple}

The model selects the exact transfer by
''snonlin\_backend''$\,=\,$'exact'.
In each source step the model hands the variance density $E$ of a cell to
the backend.
The backend returns the transfer $S_{nl}$ of $E$ and the diagonal factor
$\p S/\p E$.
The settings of the backend are attributes of the object Snonlin in the
module pyOM3.wim\_simple.snl\_exact.snonlin\_exact.
They must be set before the setup of the model, since that setup builds the
interaction table.
The attributes and their defaults are
\begin{itemize}
\item g$\,=\,$9.81, the gravitational acceleration of the backend.
It is not taken from the attribute grav of the model.
A model with another value of grav needs the same value here.
\item npts$\,=\,$24, the number $n_{\rm pts}$ of points per closed locus.
The 128 points per locus quoted for the locus tests below refer to the
routine trace\_locus, not to the operator.
\item npts\_line$\,=\,$None, which gives $4\, n_{\rm pts}$ points per
same-frequency line locus.
\item pad\_hi$\,=\,$8, the number of tail rings.
\item dfilt$\,=\,$0.025, the radius of the trivial-point filter in units
of $|\v P|$.
\item cache$\,=\,$None.
A file name here stores the table in that file.
A later setup with the same grid and the same settings reloads it.
\item mode$\,=\,$'conservative', the assembly mode.
The other choice is 'direct'.
\item points$\,=\,$None, the list of cells $(i_x, i_y)$ to evaluate.
None means all cells.
\item max\_cells$\,=\,$64, the largest number of cells per call.
With more cells the backend stops with an error.
\item verbose$\,=\,$True, progress output during the table build.
\end{itemize}
Cells not in points get $S_{nl} = 0$ and $\p S/\p E = 0$.
With the numpy backend of pyOM3 the indices $(i_x, i_y)$ count the interior
cells of the local MPI rank, without halo.
The frequency axis of the model must be geometric.
The setup of the backend checks this and stops otherwise.

\subsection*{Validation}

The implementation is validated layer by layer; the quoted accuracies refer
to the test battery shipped with the sub-package (test grids of 28--30
frequencies and 24 directions).

\begin{itemize}
\item \emph{Kernel}: degenerate two-wave interaction coefficient to $4\cdot
10^{-6}$ at the test offsets; on-shell symmetries, rotation, homogeneity and $g$-independence to
$10^{-13}$; identity \eq{eq:xcwebb} with the WAVEWATCH III coupling routine
to $3\cdot 10^{-13}$.
\item \emph{Loci}: residuals $|z|$ at returned points $4 \cdot 10^{-14}$;
crossings \eq{eq:crossings} exact; resonance closure of dimensional
quadruplets $10^{-14}$ for both signs of $q$; line-delta quadrature against a
brute-force two-dimensional integral to $5\cdot 10^{-5}$ (grid limited);
measure scaling \eq{eq:measure} to $4 \cdot 10^{-15}$; quadrature of
trace\_locus converged to $10^{-9}$ with 128 points per locus.
\item \emph{Assembly}: action budget (transfer $+$ leak) closed to
$10^{-15}$; rotation equivariance, isotropy and cubic homogeneity at
machine precision; discrete self-similarity \eq{eq:selfsim} to
$1.2 \cdot 10^{-7}$ for the closed loci on interior rings;
trivial-filter insensitivity $0.6\%$ between filter widths $0.025$ and
$0.05$.
\item \emph{Equilibria}: constant $A$ null at $2\cdot 10^{-19}$ of the
one-sided (gross) rate; Rayleigh--Jeans null at $7\cdot 10^{-5}$ with $O(h^2)$
convergence; near-vanishing transfer for the Zakharov--Filonenko spectrum
$E \propto \nu^{-4}$.
\item \emph{Magnitude and shape}: mean JONSWAP transfer with the classic
$(+,-,+)$ lobe structure.
The maximum of the positive low-frequency lobe is
$0.8$ times that of the WAM DIA.
The DIA was itself tuned against exact computations
(Hasselmann et al., 1985).
This agreement rules out any residual power of $2\pi$ or $g$ in the
normalisation chain.
\item \emph{Independent exact code}: on identical input, and with the
default interaction filtering of that code disabled, the transfer lobes agree
with the exact-NL code of van Vledder (WAVEWATCH III / SWAN) at the level of
a few percent (positive lobe ratio $0.96$, negative lobe ratio $1.02$ in the
reference case), with the independent result bracketed between the direct and
conservative modes, as expected for consistent discretisations of the same
integral.
\end{itemize}

\subsection*{References}

\noindent Hasselmann, K., 1962: On the non-linear energy transfer in a
gravity-wave spectrum. Part 1. \emph{J. Fluid Mech.}, 12, 481--500. \\
Hasselmann, S., K. Hasselmann, J.H. Allender and T.P. Barnett, 1985:
Computations and parameterizations of the nonlinear energy transfer in a
gravity-wave spectrum. Part II. \emph{J. Phys. Oceanogr.}, 15, 1378--1391. \\
Krasitskii, V.P., 1994: On reduced equations in the Hamiltonian theory of
weakly nonlinear surface waves. \emph{J. Fluid Mech.}, 272, 1--20. \\
Stuhlmeier, R., 2024: Deterministic wave forecasting with the Zakharov
equation. arXiv:2401.03539. \\
Tracy, B.A. and D.T. Resio, 1982: Theory and calculation of the nonlinear
energy transfer between sea waves in deep water. WES Report 11, US Army
Corps of Engineers. \\
van Vledder, G.Ph., 2006: The WRT method for the computation of non-linear
four-wave interactions in discrete spectral wave models.
\emph{Coastal Engineering}, 53, 223--242. \\
Webb, D.J., 1978: Non-linear transfers between sea waves.
\emph{Deep-Sea Res.}, 25, 279--298. \\
Zakharov, V.E., 1968: Stability of periodic waves of finite amplitude on the
surface of a deep fluid. \emph{J. Appl. Mech. Tech. Phys.}, 9, 190--194.

\end{document}


\end{document}
